% ============================================================
%  ROZDZIAŁ 3 — PROBLEM I SPOSÓB JEGO ROZWIĄZANIA
%  Detekcja zainfekowanych tabel (node detection).
%  Włączany do Thesis.tex przez % ============================================================
%  ROZDZIAŁ 3 — PROBLEM I SPOSÓB JEGO ROZWIĄZANIA
%  Detekcja zainfekowanych tabel (node detection).
%  Włączany do Thesis.tex przez % ============================================================
%  ROZDZIAŁ 3 — PROBLEM I SPOSÓB JEGO ROZWIĄZANIA
%  Detekcja zainfekowanych tabel (node detection).
%  Włączany do Thesis.tex przez % ============================================================
%  ROZDZIAŁ 3 — PROBLEM I SPOSÓB JEGO ROZWIĄZANIA
%  Detekcja zainfekowanych tabel (node detection).
%  Włączany do Thesis.tex przez \input{03_problem_i_architektura}
% ============================================================

\chapter{Problem i~sposób jego rozwiązania}

Rozdział precyzuje formalną definicję problemu detekcji węzłów
o~zerwanym lineage (sekcja~\ref{sec:sformulowanie}). Sekcja~\ref{sec:dane}
charakteryzuje wykorzystane dane, czyli zbiór DLG-DG-23 oraz~zewnętrzny
graf DUT wraz z~adapterem, sekcja~\ref{sec:symulacja} opisuje
metodykę symulacji broken lineage przez~usuwanie zadań,
a~sekcja~\ref{sec:architektura} architekturę potoku przetwarzania.
Zaimplementowane metody detekcji omówiono w~sekcji~\ref{sec:algorytmy},
realizację techniczną w~sekcji~\ref{sec:implementacja}, testy jednostkowe
i~odtwarzalność wyników w~sekcji~\ref{sec:testy}, a~dostępność implementacji
i~pochodzenie algorytmów w~sekcji~\ref{sec:repozytorium}.

% ============================================================
\section{Sformułowanie problemu}
\label{sec:sformulowanie}
% ============================================================

Graf lineage modeluje się jako skierowany graf DAG przepływu danych
$G = (V, E)$, gdzie $V = V_{J} \cup V_{T} \cup V_{F}$ obejmuje węzły
typu zadanie (\emph{Data Job}), tabela (\emph{Data Table}) i~pole,
a~$E$ zawiera krawędzie przepływu danych (\emph{DATA\_FLOW}) między zadaniami
a~tabelami. Zjawisko broken lineage odpowiada usunięciu z~grafu pewnego
zadania $J$. Tabele z~nim incydentne, czyli jego poprzedniki (tabele
karmiące $J$) oraz~następniki (tabele karmione przez~$J$), tracą część
informacji o~swoim pochodzeniu i~stają się~\emph{zainfekowane}
(ang.~\emph{infected}).

Obserwowany graf ma~postać $G_{\mathrm{obs}} = G \setminus \{J_1,\dots,J_m\}$,
powstałą przez~usunięcie podzbioru zadań. Zadaniem detekcji jest wyznaczenie
funkcji, która każdej tabeli $t \in V_{T}$ przypisuje wynik
$s(t) \in [0, 1]$ interpretowany jako~oszacowane prawdopodobieństwo
$p(\text{tabela } t \text{ jest zainfekowana})$. Wynikiem nie~jest~decyzja
binarna przy~ustalonym progu, lecz~ranking tabel. Tabele o~najwyższym
wyniku są~pierwszymi kandydatami do~ręcznej inspekcji.

Etykiety referencyjne (\emph{ground truth}) generowane są~automatycznie.
Tabela jest zainfekowana wtedy i~tylko wtedy, gdy była incydentna
do~co~najmniej jednego usuniętego zadania. Zbiór kandydatów ograniczono
do~węzłów typu \emph{Data Table}. Węzły typu pole (\emph{field}) łączy
z~tabelami wyłącznie relacja \emph{PARENT\_CHILD}, dlatego nie~podlegają one detekcji.

Zasadnicza trudność zadania polega na~tym, że~osierocona tabela-następnik
po~usunięciu zadania wygląda jak~prawdziwa tabela źródłowa
($\mathrm{in\_df} = 0$). Rozróżnienie tych~dwóch przypadków wymaga cech
opisujących pozycję węzła w~grafie DAG oraz~uczenia indukcyjnego
(cross-graph), a~nie~samego stopnia węzła.

% ============================================================
\section{Zbiory danych}
\label{sec:dane}
% ============================================================

W~pracy wykorzystano dwa źródła danych o~odmiennym pochodzeniu i~przeznaczeniu.
Wszystkie eksperymenty główne, obejmujące trening i~ewaluację metod
w~protokole indukcyjnym, przeprowadzono na~rzeczywistych grafach zbioru
DLG-DG-23. Uzupełnieniem jest~zewnętrzny graf DUT, zbudowany na~podstawie
benchmarku Dutkiewicza et~al.~\cite{Dutkiewicz2026}, który służy wyłącznie
do~sprawdzenia przenośności modeli wytrenowanych na~DLG-DG-23. Oba źródła
opisano w~kolejnych podsekcjach.

\subsection{Zbiór DLG-DG-23}
\label{sec:dane-dlg}

Zbiór DLG-DG-23 pochodzi z~pracy Chen et~al.~\cite{Chen2024}
i~zawiera 18~grafów lineage pobranych ze~środowiska produkcyjnego
Huawei Cloud. Dane zostały zanonimizowane, a~nazwy tabel i~atrybutów
zastąpiono identyfikatorami w~formacie UUID. Schematów bazy danych
nie~udostępniono. Zbiór opublikowano w~2024~roku jako pierwszy
otwarty zestaw rzeczywistych grafów lineage.

Każdy graf zawiera węzły trzech typów: \emph{Data Job} (zapytanie SQL
lub~zadanie ETL), \emph{Data Table}
(tabela lub~widok przechowujący dane) oraz~\emph{field} (atrybut lub~kolumna
należąca do~tabeli). Łączą je~krawędzie dwóch typów: \emph{DATA\_FLOW}
(skierowana krawędź przepływu danych między tabelą a~zadaniem)
oraz~\emph{PARENT\_CHILD} (relacja przynależności pola do~tabeli). Pełny
graf jest~zatem \emph{heterogeniczny}. Przedmiotem detekcji są~wyłącznie
węzły typu \emph{Data Table}, a~sygnał czerpany jest z~podgrafu przepływu
danych (krawędzie \emph{DATA\_FLOW} między węzłami \emph{Data Table}
i~\emph{Data Job}). Podgraf ten~ma~strukturę dwudzielną i~skierowaną,
interpretowaną jako graf DAG. Grafy różnią się skalą, a~liczba węzłów waha się
od~278 do~17~085.

Zbiór DLG-DG-23 wybrano, ponieważ jest pierwszym publicznie udostępnionym
zestawem rzeczywistych grafów lineage~\cite{Chen2024}, a~w~trakcie prowadzenia
badań nie~zidentyfikowano innego zbioru tego rodzaju. Anonimizacja nazw
węzłów ma~bezpośredni wpływ na~dobór metod, ponieważ metryki podobieństwa
oparte na~nazwach obiektów (np.~miara Jaro–Winklera stosowana przez Boińskiego
et~al.~\cite{Boinski2025}) są~nieaplikowalne. W~niniejszej pracy sygnał
ograniczono zatem do~atrybutów topologicznych grafu. Ograniczenie
to~jest zarazem częścią pytania badawczego o~to, czy~sama struktura grafu
lineage wystarcza do~wskazania tabel o~zerwanym pochodzeniu danych.
Metoda skuteczna przy~tak ograniczonym wglądzie w~graf ma~przy~tym większą
szansę na~zastosowanie praktyczne niż~podejścia wymagające pełnych
metadanych.

Tabela~\ref{tab:dlg-stats} przedstawia statystyki poszczególnych
grafów zbioru DLG-DG-23.

\begin{table}[ht]
\centering
\caption{Charakterystyka grafów zbioru DLG-DG-23.}
\label{tab:dlg-stats}
\begin{threeparttable}
\begin{tabular}{llrl}
\toprule
\textbf{ID} & \textbf{Węzły} & \textbf{Krawędzie DATA\_FLOW} & \textbf{Skala} \\
\midrule
DLG1  &    298 &    24 & mały   \\
DLG2  &    464 &    37 & mały   \\
DLG3  &    603 &    57 & mały   \\
DLG4  &    415 &    32 & mały   \\
DLG5  & 1\,299 &   175 & średni \\
DLG6  & 13\,840 & 1\,443 & duży  \\
DLG7  &    574 &    43 & mały   \\
DLG8  &    546 &    34 & mały   \\
DLG9  &    756 &    62 & mały   \\
DLG10 & 5\,378 &   497 & duży   \\
DLG11 & 2\,024 &   121 & średni \\
DLG12 & 5\,645 &   469 & duży   \\
DLG13 & 2\,157 &   229 & średni \\
DLG14 &    453 &    24 & mały   \\
DLG15 &    558 &    41 & mały   \\
DLG16 & 5\,574 &   745 & duży   \\
DLG17 &    651 &    38 & mały   \\
DLG18 & 17\,085 & 1\,694 & duży \\
\bottomrule
\end{tabular}
\begin{tablenotes}
\small
\item Skala: mały $<$ 1\,000 węzłów, średni 1\,000–5\,000, duży $>$ 5\,000.
\end{tablenotes}
\end{threeparttable}
\end{table}

\subsection{Zewnętrzny graf DUT}
\label{sec:dane-dut}

Drugim źródłem danych jest~benchmark opracowany przez Dutkiewicza
et~al.~\cite{Dutkiewicz2026} w~zespole badawczym promotora
(sekcja~\ref{sec:literatura}). Obejmuje on~syntetyczne scenariusze broken
lineage, wygenerowane dla~przykładowej bazy \emph{Northwind}~\cite{Northwind}.
Lineage zapisany jest~w~nim na~poziomie pojedynczych wierszy, w~postaci grafu
wiedzy RDF oraz~plików CSV, w~których każdy rekord wiąże wartość źródłową
z~wartością docelową, która z~niej powstała. Jest to~zatem reprezentacja
drobnoziarnista (sekcja~\ref{sec:lineage_def}), niezgodna z~gruboziarnistym
modelem tabela–zadanie, na~którym operują metody opisane w~niniejszej pracy.

Aby~ocenić na~tych danych modele wytrenowane na~DLG-DG-23 bez~zmian w~ich
implementacji, opracowano adapter (moduł \texttt{dutkiewicz\_adapter.py}),
który agreguje zależności wierszowe do~poziomu tabel. Dla~każdej tabeli
docelowej tworzony jest~węzeł typu \emph{Data Job}, do~którego prowadzą
krawędzie \emph{DATA\_FLOW} od~wszystkich tabel źródłowych, z~których
pochodzą jej~wiersze. Z~węzła tego prowadzi następnie krawędź \emph{DATA\_FLOW}
do~samej tabeli docelowej. Kolumny kluczy występujące w~zapisie lineage
dodawane są~jako~węzły typu \emph{field}, połączone z~tabelą relacją
\emph{PARENT\_CHILD}. Wynikowy graf ma~te~same typy węzłów i~krawędzi
co~grafy DLG-DG-23, dzięki czemu przechodzi przez potok przetwarzania
identycznie jak~one.

Graf DUT liczy 21~tabel, 14~zadań i~33~krawędzie przepływu danych. Jest więc
o~rzędy wielkości mniejszy od~grafów DLG-DG-23 i~ma~odmienną, łańcuchową
topologię, w~której tabele rzadko współdzielą zadania. Adapter nie~odtwarza
przy~tym pełnych schematów tabel, ponieważ benchmark ich~nie~zawiera,
dlatego liczba pól jest~mniejsza niż~w~DLG-DG-23. Poprawność odwzorowania
weryfikuje sześć testów jednostkowych (tabela~\ref{tab:testy}). Graf DUT
nie~jest~wykorzystywany w~treningu żadnej z~metod, a~wyniki uzyskane na~nim
przedstawiono w~sekcji~\ref{sec:cross_dataset}.

% ============================================================
\section{Metodyka symulacji broken lineage przez usuwanie zadań}
\label{sec:symulacja}
% ============================================================

Ponieważ zbiór DLG-DG-23 zawiera w~pełni zaobserwowane grafy lineage,
zjawisko broken lineage symulowane jest przez~kontrolowane usuwanie
zadań (nie~pojedynczych krawędzi). Podejście to~odwzorowuje
przyczynę broken lineage, ponieważ zniknięcie tabeli tymczasowej lub~ukrycie
logiki przez UDF odpowiada utracie całego węzła-zadania wraz z~jego krawędziami.

\subsection{Procedura symulacji i~konstrukcja etykiet}
\label{sec:procedura-symulacji}

Kandydatami do~usunięcia są~\emph{zadania łączące}, czyli zadania mające
co~najmniej jedną wejściową i~jedną wyjściową krawędź \emph{DATA\_FLOW}
($\mathrm{in\_df} > 0 \land \mathrm{out\_df} > 0$). Usunięcie takiego zadania
zrywa lineage po~obu stronach, ponieważ tabele karmiące je~tracą konsumenta,
a~tabele przez~nie~karmione tracą producenta.

Instancję klasyfikacji węzłów buduje się~z~grafu $G$ w~kilku krokach.
Najpierw losowany jest~podzbiór zadań łączących przeznaczonych do~usunięcia,
którego udział reguluje parametr \texttt{removal\_ratio}. W~eksperymentach
przyjęto wartość 10\%, a~wpływ tego parametru zbadano w~analizie wrażliwości
(sekcja~\ref{sec:przeciek}). Zbiór tabel zainfekowanych wyznacza się~jako~sumę
tabel incydentnych do~usuniętych zadań. Graf obserwowany $G_{\mathrm{obs}}$
powstaje następnie przez~usunięcie wybranych zadań wraz z~ich krawędziami.
Kandydatami do~detekcji są~wszystkie tabele pozostałe w~$G_{\mathrm{obs}}$,
przy~czym etykieta~1 oznacza tabelę zainfekowaną, a~etykieta~0 tabelę czystą.

Dodatkowo zaimplementowano wariant kontrolny, w~którym każda instancja
usuwa dokładnie jedno zadanie, co~daje czystą interpretację
przyczynową (jedno zdarzenie broken lineage $=$ jedno brakujące zadanie).

\subsection{Kontrola przecieku – wykluczenie węzłów izolowanych}
\label{sec:exclude-isolated}

Usunięcie zadania może~pozbawić tabelę wszystkich krawędzi
przepływu danych ($\mathrm{in\_df} = 0 \land \mathrm{out\_df} = 0$).
Ponieważ w~oryginalnych grafach tabele izolowane nie~występują, trywialna
reguła „tabela izolowana $\Rightarrow$ zainfekowana'' pozwala rozpoznać
znaczną część pozytywów bez~żadnego uczenia i~sztucznie zawyża wynik.
Aby~zapewnić uczciwą ewaluację, wprowadzono flagę \texttt{exclude\_isolated},
która usuwa z~puli kandydatów tabele w~pełni izolowane. Wariant ten~stawia
trudniejsze i~bardziej realistyczne zadanie, polegające na~wykryciu tabel
o~częściowo zerwanym lineage, które nadal mają część połączeń. W~pracy raportowane
są~oba warianty, a~wariant uczciwy (z~wykluczeniem izolowanych) traktowany
jest jako~główny.

Tabela~\ref{tab:symulacja} przedstawia skalę symulacji dla~czterech wartości
parametru \texttt{removal\_ratio}. Łącznie w~18~grafach zbioru DLG-DG-23
znajduje się 1453~zadania łączące (od~3 do~404 na~graf), spośród których usuwana
jest~losowa część. Liczby podano dla~wariantu przed wykluczeniem tabel
izolowanych, ponieważ tylko wtedy można zmierzyć ich~udział wśród pozytywów.

\begin{table}[ht]
\centering
\caption[Charakterystyka symulacji broken lineage]{Charakterystyka symulacji broken lineage w~zależności
od~parametru \texttt{removal\_ratio}. Wartości uśrednione po~18~grafach
i~trzech ziarnach losowości (42--44). W~ostatniej kolumnie podano odsetek
tabel zainfekowanych, które utraciły wszystkie krawędzie przepływu
danych i~są~rozpoznawalne trywialną regułą izolacji.}
\label{tab:symulacja}
\begin{tabular}{ccccc}
\toprule
\texttt{removal\_ratio} & \textbf{Usuniętych} & \textbf{Tabel} & \textbf{Udział} & \textbf{Pozytywów} \\
 & \textbf{zadań} & \textbf{zainfekowanych} & \textbf{pozytywów} & \textbf{izolowanych} \\
\midrule
0{,}05 & 4{,}2 & 15{,}1 & 11{,}4\% & 19{,}7\% \\
0{,}10 & 8{,}0 & 25{,}1 & 19{,}0\% & 22{,}9\% \\
0{,}20 & 16{,}2 & 45{,}2 & 34{,}1\% & 29{,}8\% \\
0{,}30 & 24{,}3 & 62{,}6 & 47{,}3\% & 34{,}5\% \\
\bottomrule
\end{tabular}
\end{table}

Dwie obserwacje z~tabeli mają~znaczenie dla~konstrukcji eksperymentu.
Przyjęta w~biegu kanonicznym wartość $0{,}10$ daje średnio 25~tabel
zainfekowanych na~graf przy~19\% pozytywów. Zadanie pozostaje więc~wyraźnie
niezrównoważone, ale~klasa pozytywna jest~dostatecznie liczna, by~metryki
rankingowe były wiarygodne. Udział pozytywów trywialnych (tabel w~pełni
izolowanych) rośnie z~kolei z~19{,}7\% do~34{,}5\% wraz ze~wzrostem
odsetka usuwanych zadań. Im~silniejsze zaburzenie, tym~większa część zadania
daje się~rozwiązać bez~żadnego uczenia. Liczby te~określają skalę przecieku,
któremu zapobiega opisana wyżej flaga \texttt{exclude\_isolated}. Przy~biegu
kanonicznym usuwa ona~z~puli kandydatów blisko jedną czwartą pozytywów,
i~to~tych najłatwiejszych.

% ============================================================
\section{Architektura rozwiązania}
\label{sec:architektura}
% ============================================================

Potok przetwarzania składa się z~czterech komponentów: modułu wczytywania
grafów (wraz z~adapterem grafu DUT), modułu symulacji broken lineage (budowa
instancji klasyfikacji węzłów), metod detekcji oraz~modułu ewaluacji indukcyjnej.

Rysunek~\ref{rys:potok} przedstawia schemat potoku przetwarzania, od~wczytania
grafów, przez symulację broken lineage i~ekstrakcję cech, po~scoring detektora
i~wyznaczenie metryk rankingowych. Oba źródła danych trafiają do~wspólnego
potoku. Grafy DLG-DG-23 wczytywane są~bezpośrednio z~plików JSON, natomiast
graf DUT przechodzi najpierw przez adapter, który sprowadza go~do~tego samego
modelu danych. Od~etapu symulacji oba rodzaje grafów przetwarzane są~identycznie,
różni je~jedynie rola w~ewaluacji.

\begin{figure}[ht]
  \centering
  \includegraphics[width=\linewidth]{rys_potok.pdf}
  \caption[Schemat potoku przetwarzania]{Schemat potoku przetwarzania. Grafy
           DLG-DG-23 oceniane są~w~protokole indukcyjnym leave-one-graph-out
           (trening na~17~grafach, test na~grafie niewidzianym). Graf DUT,
           sprowadzony adapterem do~modelu tabela–zadanie, służy wyłącznie
           do~testu modeli wytrenowanych na~DLG-DG-23.}
  \label{rys:potok}
\end{figure}

\subsection{Moduł wczytywania danych}

Dane wejściowe każdego grafu przechowywane są~w~dwóch plikach JSON:
\texttt{DLGi-node.json} (opisy węzłów z~atrybutem \texttt{asset\_type})
oraz~\texttt{DLGi-edge.json} (krawędzie z~atrybutami \texttt{relation\_type}
i~\texttt{relation\_id}). Moduł \texttt{loader.py} buduje z~nich obiekt
\texttt{nx.DiGraph} (NetworkX) i~zwraca słownik wszystkich 18~grafów.
Graf DUT wczytywany jest~przez moduł \texttt{dutkiewicz\_adapter.py}, który
zwraca obiekt tego samego typu, z~atrybutami węzłów i~krawędzi zgodnymi
z~konwencją DLG-DG-23 (sekcja~\ref{sec:dane-dut}). Dalsze moduły potoku
nie~rozróżniają więc~źródła grafu.

\subsection{Moduł symulacji usuwania zadań}

Moduł \texttt{job\_removal.py} realizuje metodykę z~sekcji~\ref{sec:symulacja}:
identyfikuje zadania łączące, usuwa ich podzbiór, wyznacza zbiór zainfekowanych
tabel i~zwraca instancję (graf $G_{\mathrm{obs}}$, listę kandydatów,
etykiety, listę usuniętych zadań). Obsługuje warianty \texttt{side},
określający stronę incydencji, po~której tabele uznaje się~za~zainfekowane
(strona producenta, konsumenta lub~obie), oraz~\texttt{exclude\_isolated}. Funkcja \texttt{inductive\_split} dzieli
identyfikatory grafów na~rozłączne zbiory treningowy i~testowy (cross-graph).

\subsection{Moduł ewaluacji i~metryki rankingowe}

Moduł \texttt{node\_metrics.py} wyznacza metryki rankingowe: AUC-ROC,
AUC-PR, Precision@k, Recall@k, Hits@k oraz~Brier score. Wyniki agregowane
są~po~grafach testowych w~protokole indukcyjnym. Dla~stabilności oszacowań
stosuje się powtórzenia z~różnymi ziarnami i~trymowanie wartości skrajnych
(min/max).

Ewaluację prowadzi się~w~protokole indukcyjnym cross-graph.
Model trenuje się~na~podzbiorze grafów, a~testuje na~grafach niewidzianych
(leave-one-graph-out lub~podział zbioru 18~grafów). Zbiory treningowy
i~testowy są~rozłączne, co~sprawdza zdolność metody do~generalizacji
na~nowe grafy, a~nie~jedynie do~zapamiętania jednego. Konsekwencją tego
protokołu jest~wykluczenie metody Node2Vec~\cite{Grover2016}, która uczy zanurzeń
węzłów transduktywnie, dla~konkretnego grafu, przez~co~embeddingi
są~geometrycznie nieporównywalne między grafami~\cite{Hamilton2017}.
Obserwacja ta~motywuje zastosowanie indukcyjnych
grafowych sieci neuronowych (GraphSAGE~\cite{Hamilton2017},
R-GCN~\cite{Schlichtkrull2018}, HGT~\cite{Hu2020}), wskazanych w~rozdziale~5 jako~kierunek dalszych badań.

% ============================================================
\section{Zaimplementowane metody detekcji}
\label{sec:algorytmy}
% ============================================================

Zaimplementowano trzy grupy metod detekcji zainfekowanych tabel, różniące się
ilością wykorzystywanej informacji i~stopniem złożoności. Heurystyki strukturalne
(sekcja~\ref{sec:heurystyki}) przypisują tabeli wynik na~podstawie prostych
reguł i~nie~wymagają uczenia, dlatego wyznaczają dolny punkt odniesienia.
Klasyczne klasyfikatory uczenia maszynowego (sekcja~\ref{sec:ml}) uczą się
na~grafach treningowych zależności między wektorem cech topologicznych
a~etykietą, co~pozwala ocenić, ile~wnosi samo uczenie indukcyjne. Autorski
detektor LineageDetector (sekcja~\ref{sec:lineage-detector}) rozszerza ten~wektor
o~relacyjny score kompletności i~uzupełnia model o~kalibrację prawdopodobieństw.

Wszystkie metody mają wspólny interfejs. Dla~każdej tabeli-kandydata zwracają
wynik z~przedziału $[0, 1]$, na~podstawie którego budowany jest~ranking. Dzięki
temu każda z~nich jest~oceniana w~tym samym potoku, na~tych samych instancjach
i~tymi samymi metrykami, a~różnice wyników można przypisać samym metodom.

% ============================================================
\subsection{Heurystyki strukturalne}
\label{sec:heurystyki}
% ============================================================

Heurystyki strukturalne przypisują tabeli wynik podejrzliwości wyłącznie
na~podstawie jej pozycji w~grafie, bez~etapu uczenia. Są~naturalnym
punktem odniesienia (\emph{baseline}) dla~metod uczonych. Miary oparte
na~wspólnych sąsiadach (Common Neighbors, Jaccard,
Adamic-Adar)~\cite{LibenNowell2007} są~w~grafie
dwudzielnym nieskuteczne, ponieważ węzły tego~samego typu nigdy nie~mają
wspólnych sąsiadów bezpośrednich. Z~tego powodu wykluczono je~z~implementacji.

W~pracy zaimplementowano cztery heurystyki bazowe. Piątym punktem
odniesienia w~tej~grupie jest~sam \emph{score kompletności}
(sekcja~\ref{sec:lineage-detector}), oceniany dodatkowo w~oderwaniu
od~detektora, co~pozwala rozdzielić wkład sygnału relacyjnego od~wkładu
uczenia. W~zestawieniach rozdziału~4 występuje on~zatem jako~osobna metoda
obok czterech poniższych:

\begin{itemize}
  \item \textbf{Anomalia stopnia} (\texttt{degree\_anomaly}): wynik
        $1/(1+\deg(t))$, gdzie $\deg(t)$ jest stopniem całkowitym tabeli.
        Heurystyka opiera się~na~założeniu, że~tabela, która utraciła
        zadanie, ma~mniej połączeń niż~przed zerwaniem lineage.
  \item \textbf{Brzeg grafu} (\texttt{boundary}): suma wskaźników
        $\mathbf{1}[\mathrm{in\_df}=0] + \mathbf{1}[\mathrm{out\_df}=0]$.
        Tabele na~brzegu przepływu (korzenie i~liście) są~podejrzane,
        ponieważ zerwanie lineage często przesuwa tabelę na~brzeg grafu.
  \item \textbf{Niska łączność z~zadaniami}
        (\texttt{low\_job\_connectivity}): wynik $1/(1+n_{J}(t))$,
        gdzie $n_{J}(t)$ to~liczba unikatowych zadań sąsiadujących
        z~tabelą przez krawędzie \emph{DATA\_FLOW}.
  \item \textbf{Reguła korzenia} (\texttt{rule\_root}): binarny wskaźnik
        $\mathbf{1}[\mathrm{in\_df}=0]$, zgodnie z~którym każda tabela
        bez~producenta jest uznawana za~podejrzaną. To~celowo trywialny punkt odniesienia,
        obrazujący siłę najprostszego możliwego symptomu.
\end{itemize}

Heurystyki te~mierzą wyłącznie lokalny symptom (stopień, rola).
Ich~ograniczenie, przede wszystkim nieodróżnianie osieroconej tabeli
od~prawdziwego źródła, motywuje metody uczone oraz~autorski detektor
LineageDetector (sekcja~\ref{sec:lineage-detector}), który wzbogaca opis
węzła o~relacyjny score kompletności (peer-consistency).

% ============================================================
\subsection{Klasyczne metody uczenia maszynowego}
\label{sec:ml}
% ============================================================

Klasyczne metody ML traktują detekcję jako~binarną klasyfikację węzła, w~której
wejściem jest wektor cech opisujących pozycję tabeli w~grafie DAG,
a~wyjściem prawdopodobieństwo zainfekowania. Ze~względu na~anonimizację
(UUID) jako~cechy przyjęto wyłącznie atrybuty topologiczne, ponieważ metryki
oparte na~nazwach są~nieaplikowalne.

\subsubsection{Wektor cech węzła}

Tabela~\ref{tab:cechy} przedstawia pełny wektor 13~cech topologicznych
wyznaczanych dla~każdej tabeli-kandydata. Cechy opisują pozycję węzła
w~grafie przepływu danych, bo~to~ona, a~nie~sam stopień, pozwala odróżnić
prawdziwą tabelę źródłową od~tabeli osieroconej po~usunięciu zadania.

\begin{table}[ht]
\centering
\caption[Wektor cech topologicznych węzła]{Wektor cech topologicznych węzła (moduł \texttt{node\_features}).}
\label{tab:cechy}
\small
\begin{tabular}{cl>{\raggedright\arraybackslash}p{5.2cm}>{\raggedright\arraybackslash}p{4.6cm}}
\toprule
\textbf{Nr} & \textbf{Cecha} & \textbf{Opis} & \textbf{Rola w~detekcji} \\
\midrule
1 & \texttt{in\_df} & Liczba wejściowych krawędzi \emph{DATA\_FLOW}
    (od~zadań) & Brak producenta ($=0$) to~główny symptom zerwania \\
2 & \texttt{out\_df} & Liczba wyjściowych krawędzi \emph{DATA\_FLOW}
    (do~zadań) & Symetryczny symptom, czyli utrata konsumenta \\
3 & \texttt{degree} & Stopień całkowity węzła (wszystkie typy krawędzi)
    & Ogólna łączność, spadająca po~usunięciu zadania \\
4 & \texttt{role} & Rola w~przepływie (źródło, pośrednia, ujście,
    poza przepływem) & Pozycja funkcjonalna tabeli w~DAG \\
5 & \texttt{flow\_depth} & Głębokość w~DAG (odległość od~źródeł)
    & Osierocona tabela ma~sztucznie małą głębokość \\
6 & \texttt{flow\_reach} & Zasięg w~dół DAG (odległość do~ujść)
    & Symetrycznie sztucznie mały zasięg \\
7 & \texttt{depth\_plus\_reach} & Suma głębokości i~zasięgu
    & Długość ścieżki przez węzeł, wykrywa „skrócone'' łańcuchy \\
8 & \texttt{is\_root} & Wskaźnik $\mathbf{1}[\mathrm{in\_df}=0]$
    & Binarna wersja symptomu utraty producenta \\
9 & \texttt{is\_leaf} & Wskaźnik $\mathbf{1}[\mathrm{out\_df}=0]$
    & Binarna wersja symptomu utraty konsumenta \\
10 & \texttt{n\_field\_children} & Liczba pól tabeli (krawędzie
    \emph{PARENT\_CHILD}) & Rozmiar tabeli, sygnał niezależny
    od~przepływu \\
11 & \texttt{n\_job\_neighbors} & Liczba unikatowych zadań sąsiadujących
    & Łączność z~warstwą zadań \\
12 & \texttt{avg\_job\_in\_df} & Średni stopień wejściowy sąsiednich zadań
    & Charakterystyka otoczenia węzła \\
13 & \texttt{avg\_job\_out\_df} & Średni stopień wyjściowy sąsiednich
    zadań & Jak~wyżej, od~strony wyjść \\
\bottomrule
\end{tabular}
\end{table}

Detektor LineageDetector (sekcja~\ref{sec:lineage-detector})
rozszerza ten wektor o~trzy cechy kompletności
(\texttt{peer\_producer\_score}, \texttt{peer\_consumer\_score},
\texttt{completeness\_anomaly}), opisane w~sekcji~\ref{sec:lineage-detector}.

\subsubsection{Random Forest}

Random Forest (RF)~\cite{Breiman2001} to~metoda uczenia zespołowego oparta na~uśrednianiu
predykcji drzew decyzyjnych trenowanych na~losowych podpróbkach danych
i~cech (\emph{bagging}). Dla~obsługi niezrównoważonych klas zastosowano
\texttt{class\_weight=\textquotedbl{}balanced\textquotedbl{}}.

\subsubsection{RUSBoost}

RUSBoost (\emph{Random Under-Sampling Boosting})~\cite{Seiffert2010} łączy boosting AdaBoost
z~losowym podpróbkowaniem klasy większościowej przed~każdą iteracją.
Zaimplementowano go~dla~bezpośredniego nawiązania do~wyników Boińskiego
et~al.~\cite{Boinski2025}, przy~czym w~niniejszej pracy klasyfikuje on~węzły,
a~nie~pary obiektów.

\subsubsection{LightGBM}

LightGBM~\cite{Ke2017} to~implementacja gradientowego boostingu drzew decyzyjnych,
budująca kolejne drzewa sekwencyjnie, każde minimalizując błąd resztkowy
poprzednika. Niezbalansowanie klas kompensuje parametr
\texttt{class\_weight}~$=$~\texttt{balanced}, przypisujący obserwacjom wagi
odwrotnie proporcjonalne do~liczności klasy.

% ============================================================
\subsection{Autorski detektor LineageDetector}
\label{sec:lineage-detector}
% ============================================================

Autorski detektor \texttt{LineageDetector} (klasa
\texttt{LineageHealthDetector}) łączy w~jednym skalibrowanym modelu trzy
elementy. Podstawę stanowią cechy pozycji w~DAG, wyznaczane przez moduł
\texttt{node\_features} i~opisujące węzeł lokalnie. Uzupełnia je~score
kompletności \emph{peer-consistency}, opisany w~dalszej części sekcji,
który wnosi sygnał relacyjny i~stanowi główny element wkładu własnego pracy.
Całość zamyka kalibracja prawdopodobieństw metodą izotoniczną
(\texttt{CalibratedClassifierCV}), dzięki której wynik detektora można
interpretować jako~oszacowanie $p(\text{zainfekowana})$, a~nie~jedynie
jako~wartość porządkującą ranking.

\subsubsection{Score kompletności peer-consistency}

Zaproponowany w~pracy score kompletności opiera się na~założeniu,
że~tabela o~zerwanym lineage jest~\emph{strukturalnie niespójna}
ze~swoimi sąsiadami (peerami) w~grafie DAG. Dla~tabeli-korzenia
($\mathrm{in\_df} = 0$) rozważa się~zbiór peerów-wejść, czyli inne tabele
wpadające do~tych~samych zadań downstream. Jeśli~peerzy mają
producenta, a~rozpatrywana tabela nie, to~jej brak producenta jest~podejrzany,
ponieważ prawdziwe źródło danych sąsiaduje zwykle z~innymi źródłami. Symetrycznie
dla~tabeli-liścia ($\mathrm{out\_df} = 0$) i~jej peerów po~stronie konsumpcji.
Formalnie dla~tabeli $t$ definiuje się:
%
\[
  \mathrm{prod}(t) = \frac{1}{|P(t)|}\sum_{p \in P(t)} \mathbf{1}[\mathrm{in\_df}(p) > 0],
  \qquad
  \mathrm{cons}(t) = \frac{1}{|C(t)|}\sum_{p \in C(t)} \mathbf{1}[\mathrm{out\_df}(p) > 0],
\]
%
gdzie $P(t)$ i~$C(t)$ oznaczają zbiory peerów odpowiednio po~stronie
producenta i~konsumenta, a~$\mathbf{1}[\cdot]$ funkcję wskaźnikową.
Bramkowany sygnał anomalii aktywuje się~wyłącznie na~brzegu grafu:
%
\[
  \mathrm{anomaly}(t) =
  \max\bigl(\mathbf{1}[\mathrm{in\_df}(t)=0]\cdot\mathrm{prod}(t),\;
            \mathbf{1}[\mathrm{out\_df}(t)=0]\cdot\mathrm{cons}(t)\bigr).
\]
%
Sygnał ten~ma~charakter relacyjny. Różni się~jakościowo od~lokalnego
stopnia węzła i~dostarcza informacji niedostępnej klasyfikatorowi
korzystającemu wyłącznie z~cech lokalnych. Definicja jest~symetryczna
względem obu~stron przepływu, jednak na~zbiorze DLG-DG-23 aktywna
okazuje się~wyłącznie strona producenta, ponieważ każde zadanie w~tym~zbiorze
zapisuje dokładnie jedną tabelę, przez~co~zbiór $C(t)$ jest~zawsze pusty.
Konsekwencje tej~własności omówiono w~sekcji~\ref{sec:waznosc}. Trzy wynikające z~niego cechy
(\texttt{peer\_producer\_score}, \texttt{peer\_consumer\_score},
\texttt{completeness\_anomaly}) rozszerzają wektor cech z~tabeli~\ref{tab:cechy}
i~są~dostępne wyłącznie detektorowi LineageDetector.

Bazowym estymatorem jest~domyślnie Random Forest, a~wariantem
alternatywnym LightGBM. Interfejs \texttt{fit}/\texttt{score} jest~zgodny z~pozostałymi
klasyfikatorami, co~umożliwia bezpośrednie porównanie w~tym~samym potoku.
Konstrukcja detektora umożliwia badania ablacyjne. Wyłączenie kalibracji
(\texttt{calibrate=False}) izoluje jej~wpływ na~Brier score, a~wyłączenie
cech kompletności (\texttt{use\_completeness=False}) wpływ sygnału
relacyjnego na~AUC.

% ============================================================
\section{Implementacja}
\label{sec:implementacja}
% ============================================================

Sekcja opisuje realizację techniczną potoku z~sekcji~\ref{sec:architektura}:
zastosowane narzędzia, organizację kodu oraz~te~fragmenty implementacji,
które są~istotne dla~zrozumienia metody, a~nie~wynikają wprost z~jej~opisu
formalnego.

\subsection{Stos technologiczny}

Całość zaimplementowano w~języku Python. Tabela~\ref{tab:stos} zestawia
wykorzystane biblioteki wraz z~wersjami, w~których wykonano bieg kanoniczny.
Podanie wersji jest~istotne dla~odtwarzalności, ponieważ implementacje algorytmów
uczenia maszynowego bywają modyfikowane między wydaniami, co~może~wpływać
na~wyniki liczbowe.

\begin{table}[ht]
\centering
\caption{Stos technologiczny i~wersje bibliotek użyte w~biegu kanonicznym.}
\label{tab:stos}
\begin{tabular}{lll}
\toprule
\textbf{Składnik} & \textbf{Wersja} & \textbf{Rola w~pracy} \\
\midrule
Python~\cite{Python}                 & 3.10.2  & Środowisko uruchomieniowe \\
NetworkX~\cite{NetworkX}             & 3.4.2   & Reprezentacja i~przetwarzanie grafów \\
NumPy~\cite{NumPy}                   & 1.26.4  & Operacje macierzowe, wektory cech \\
pandas~\cite{pandas}                 & 1.5.1   & Agregacja wyników, eksport CSV \\
scikit-learn~\cite{scikit-learn}     & 1.6.1   & Random Forest, kalibracja izotoniczna, metryki \\
LightGBM~\cite{Ke2017}               & 4.6.0   & Gradientowy boosting drzew \\
imbalanced-learn~\cite{imblearn}     & 0.14.1  & RUSBoost \\
SciPy~\cite{SciPy}                   & 1.15.3  & Test kolejności par Wilcoxona \\
Matplotlib~\cite{Matplotlib}         & 3.8.0   & Rysunki \\
pytest~\cite{pytest}                 & 9.0.2   & Testy jednostkowe \\
\bottomrule
\end{tabular}
\end{table}

\subsection{Organizacja repozytorium}

Kod podzielono na~pakiety odpowiadające etapom potoku, a~skrypty
uruchomieniowe oddzielono od~modułów bibliotecznych, aby~każdy eksperyment
z~rozdziału~4 dawał się~odtworzyć jednym poleceniem. Listing~\ref{lst:repo}
przedstawia strukturę katalogów.

\begin{lstlisting}[caption={Struktura repozytorium (fragmenty istotne
dla~detekcji węzłów).},label={lst:repo},numbers=none,language={},float=htbp]
src/
  data/
    loader.py               # wczytanie 18 grafow DLG-DG-23 z JSON
    dutkiewicz_adapter.py   # adapter benchmarku zewnetrznego (graf DUT)
  detection/
    job_removal.py          # symulacja broken lineage, etykiety, podzial
    node_features.py        # 13 cech topologicznych wezla
    completeness.py         # score peer-consistency (wklad wlasny)
    lineage_detector.py     # LineageHealthDetector
    node_classifier.py      # heurystyki + klasyczne klasyfikatory
    node_metrics.py         # metryki rankingowe
    run_detection_experiments.py   # bieg kanoniczny (rozdz. 4.2)
    run_ablation.py                # badanie ablacyjne (rozdz. 4.5)
    run_cross_dataset.py           # test cross-dataset (rozdz. 4.9)
    run_removal_sensitivity.py     # analiza wrazliwosci (rozdz. 4.6)
scripts/
  test_istotnosci.py        # testy Wilcoxona (rozdz. 4.4)
  tabela_symulacji.py       # tabela charakterystyki symulacji
tests/                      # 83 testy jednostkowe (pytest)
results/                    # wyniki biegow w formacie CSV
\end{lstlisting}

\subsection{Format danych i~wczytywanie}

Każdy graf zbioru DLG-DG-23 opisany jest~parą plików JSON: lista węzłów
z~atrybutem \texttt{asset\_type} oraz~lista krawędzi z~atrybutami
\texttt{relation\_id} i~\texttt{relation\_type} (listing~\ref{lst:json}).
Moduł \texttt{loader.py} przekształca je~w~obiekt \texttt{nx.DiGraph}
i~przenosi typ węzła oraz~typ relacji do~atrybutów grafu. Dzięki temu
pozostałe moduły filtrują podgraf przepływu danych bez~ponownego sięgania
do~plików.

\begin{lstlisting}[caption={Fragment plików wejściowych
\texttt{DLGi-node.json} i~\texttt{DLGi-edge.json} (identyfikatory skrócone).},
label={lst:json},numbers=none,language={},float=htbp]
// DLG1-node.json
{"nodes": [
  {"asset_id": "3f2a...c81", "asset_type": "Data Table"},
  {"asset_id": "9b07...4de", "asset_type": "Data Job"},
  {"asset_id": "a154...77f", "asset_type": "Data Field"}
]}

// DLG1-edge.json
{"edges": [
  {"relation_id": "r-0012", "relation_type": "DATA_FLOW",
   "source": "3f2a...c81", "target": "9b07...4de"},
  {"relation_id": "r-0013", "relation_type": "PARENT_CHILD",
   "source": "3f2a...c81", "target": "a154...77f"}
]}
\end{lstlisting}

% Listingi 3.1 i 3.2 maja zostac zlozone przed ta podsekcja, aby nie
% rozcinaly jej pierwszego akapitu.
\FloatBarrier
\subsection{Wyznaczanie score'u peer-consistency}

Najistotniejszym fragmentem implementacji jest~wyznaczanie relacyjnego
score'u kompletności, zdefiniowanego wzorami w~sekcji~\ref{sec:lineage-detector}.
Listing~\ref{lst:peer} pokazuje jego~rdzeń. Sam wzór nie~precyzuje, że~zbiory
peerów budowane są~przez zadania. Peerami tabeli
po~stronie producenta są~inne tabele wpadające do~tych~samych zadań
downstream, w~grafie dwudzielnym nie~istnieje bowiem bezpośrednia krawędź
tabela–tabela. Widoczne w~listingu bramkowanie wskaźnikami odpowiada wprost
drugiemu ze~wzorów z~sekcji~\ref{sec:lineage-detector}. Bez~niego score
mierzyłby jedynie własności otoczenia, a~nie~niespójność tabeli
z~jej~otoczeniem.

\begin{lstlisting}[caption={Rdzeń wyznaczania score'u peer-consistency
(moduł \texttt{completeness.py}).},label={lst:peer},language=Python,float=htbp]
# Strona producenta: peerzy = inne wejscia do downstream-jobow tej tabeli.
downstream_jobs = [v for _, v, d in G.out_edges(t, data=True)
                   if d.get("relation_type") == "DATA_FLOW"
                   and G.nodes[v].get("asset_type") == "Data Job"]
prod_peers = set()
for j in downstream_jobs:
    prod_peers.update(job_inputs.get(j, ()))
prod_peers.discard(t)
prod_score = (np.mean([1.0 if in_df[p] > 0 else 0.0 for p in prod_peers])
              if prod_peers else 0.0)

# Bramkowana anomalia: liczy sie tylko na brzegu (korzen/lisc).
anomaly = 0.0
if in_df[t] == 0:
    anomaly = max(anomaly, prod_score)
if out_df[t] == 0:
    anomaly = max(anomaly, cons_score)
\end{lstlisting}

\FloatBarrier
\subsection{Protokół indukcyjny w~kodzie}

Pętla ewaluacji realizuje schemat leave-one-graph-out, w~którym dla~każdego grafu
testowego model uczy się~na~wszystkich pozostałych. Listing~\ref{lst:logo}
przedstawia jej~szkielet. Kluczowa dla~poprawności protokołu jest~kolejność
operacji. Instancje wszystkich grafów budowane są~przed pętlą,
przy~wspólnym ziarnie, a~pula treningowa powstaje przez~wykluczenie grafu
testowego z~gotowego słownika. Model nie~ma~więc~dostępu do~żadnej
informacji z~grafu testowego, także pośrednio przez~dobór usuwanych zadań.

\begin{lstlisting}[caption={Szkielet pętli leave-one-graph-out
(moduł \texttt{run\_detection\_experiments.py}).},label={lst:logo},
language=Python,float=htbp]
datasets = build_datasets(graphs, removal_ratio, seed, side,
                          exclude_isolated=exclude_isolated)
for test_g in datasets:
    ds_test = datasets[test_g]
    train_pool = [datasets[g] for g in datasets if g != test_g]

    for name in FIT_MODELS:
        clf = _make_model(name, seed).fit(train_pool)
        m = node_detection_metrics(ds_test["labels"], clf.score(ds_test))
        rows.append({"graph": test_g, "algorithm": name, **m})
\end{lstlisting}

Wyniki z~kolejnych ziaren agregowane są~średnią z~trymowaniem wartości
skrajnych (odrzucenie najmniejszej i~największej), co~ogranicza wpływ
pojedynczego nietypowego losowania na~oszacowanie.

\subsection{Złożoność obliczeniowa}

Koszt przetwarzania jednego grafu jest~liniowy względem liczby krawędzi.
Wyznaczenie stopni \emph{DATA\_FLOW}, ekstrakcja cech oraz~budowa zbiorów
peerów wymagają stałej liczby przejść po~krawędziach, a~metadane przepływu
(głębokość i~zasięg w~DAG) wyznaczane są~przeszukiwaniem grafu, również
w~czasie $O(|V| + |E|)$. Dominującym składnikiem kosztu nie~jest~zatem
przetwarzanie grafu, lecz~trenowanie klasyfikatorów. W~protokole
leave-one-graph-out model uczony jest~od~nowa dla~każdego grafu testowego,
co~przy~18~grafach, trzech ziarnach i~czterech wariantach detektora daje
216~pełnych treningów. Bieg ablacyjny z~sekcji~\ref{sec:ablacja} zajął
z~tego~powodu 412~sekund, podczas gdy sama ekstrakcja cech dla~wszystkich
18~grafów trwa poniżej sekundy. Skala zbioru DLG-DG-23 (do~17~085~węzłów
na~graf) nie~stanowi więc~ograniczenia wydajnościowego.

% ============================================================
\section{Testy jednostkowe i~odtwarzalność}
\label{sec:testy}
% ============================================================

\subsection{Zakres testów}

Poprawność implementacji zweryfikowano zestawem 83~testów jednostkowych
uruchamianych narzędziem \texttt{pytest}. W~wersji implementacji
opisanej w~pracy wszystkie testy kończą się~powodzeniem.
Tabela~\ref{tab:testy} przedstawia ich~rozkład na~moduły.

Poza poprawnością techniczną testy sprawdzają dwie własności istotne
metodycznie. Pierwszą jest~brak przecieku między zbiorem treningowym
a~testowym, weryfikowany przez sprawdzenie, czy~identyfikatory grafów w~obu
zbiorach są~rozłączne. Drugą jest~determinizm symulacji, który oznacza,
że~przy~ustalonym ziarnie losowości powtórne wykonanie daje identyczny zbiór
usuniętych zadań i~identyczne etykiety. Bez~tej własności liczby raportowane
w~rozdziale~4 nie~byłyby~odtwarzalne.

\begin{table}[ht]
\centering
\caption{Rozkład testów jednostkowych na~moduły.}
\label{tab:testy}
\small
\begin{tabular}{lr>{\raggedright\arraybackslash}p{7.2cm}}
\toprule
\textbf{Plik testowy} & \textbf{Liczba} & \textbf{Weryfikowane własności} \\
\midrule
\texttt{test\_loader.py} & 29 &
  Poprawność typów węzłów i~krawędzi, zgodność liczby obiektów z~plikami
  źródłowymi, wczytanie wszystkich 18~grafów \\
\texttt{test\_splitter.py} & 17 &
  Rozłączność podziałów, zgodność liczności podzbiorów \\
\texttt{test\_job\_removal.py} & 12 &
  Identyfikacja zadań łączących, poprawność etykiet, rozłączność zbiorów
  treningowego i~testowego, działanie flagi \texttt{exclude\_isolated} \\
\texttt{test\_scenario\_splitter.py} & 9 &
  Scenariusze usuwania krawędzi (badanie wstępne) \\
\texttt{test\_dutkiewicz\_adapter.py} & 6 &
  Mapowanie benchmarku Dutkiewicza et~al.\ na~model
  tabela–zadanie–\emph{DATA\_FLOW} (graf DUT) \\
\texttt{test\_node\_features.py} & 4 &
  Wymiar wektora cech, zgodność z~\texttt{feature\_names()},
  wartości brzegowe dla~węzłów izolowanych \\
\texttt{test\_completeness.py} & 3 &
  Poprawność score'u peer-consistency, bramkowanie sygnału anomalii \\
\texttt{test\_lineage\_detector.py} & 3 &
  Interfejs \texttt{fit}/\texttt{score}, zakres zwracanych
  prawdopodobieństw, warianty ablacyjne \\
\midrule
\textbf{Razem} & \textbf{83} & \\
\bottomrule
\end{tabular}
\end{table}

\subsection{Odtwarzalność wyników}

Wszystkie liczby przedstawione w~rozdziale~4 pochodzą z~jednego
\emph{biegu kanonicznego} o~ustalonej konfiguracji: 18~grafów,
\texttt{removal\_ratio}~$=0{,}10$, \texttt{side}~$=$~\texttt{both},
trzy powtórzenia z~ziarnami 42--44, protokół indukcyjny leave-one-graph-out,
próg wiarygodności $\mathrm{MIN\_POS} = 5$, wariant uczciwy
(\texttt{exclude\_isolated}). Wyniki pośrednie zapisywane są~w~postaci plików
CSV w~katalogu \texttt{results/}, osobno dla~wariantu uczciwego
i~zawyżonego, wraz z~opisem konfiguracji. Eksperymenty poboczne
(ablacja, analiza wrażliwości, test cross-dataset) mają~własne skrypty
i~własne pliki wynikowe, wymienione w~listingu~\ref{lst:repo}, co~pozwala
odtworzyć każdą tabelę rozdziału~4 niezależnie od~pozostałych.

% ============================================================
\section{Dostępność implementacji i~pochodzenie algorytmów}
\label{sec:repozytorium}
% ============================================================

Całość implementacji opisanej w~niniejszej pracy, obejmującą moduły
wczytywania danych, symulacji broken lineage, ekstrakcji cech, metod
detekcji oraz~ewaluacji, wraz z~testami jednostkowymi i~skryptami
odtwarzającymi eksperymenty z~rozdziału~4, udostępniono w~publicznym
repozytorium Git
pod~adresem \texttt{https://github.com/Catcuss/magisterka-data-lineage}
(dostęp zweryfikowano 19~sierpnia 2026~r.). Repozytorium zawiera instrukcję
uruchomienia (dodatek~\ref{dod:uruchomienie}) oraz~pliki konfiguracyjne biegu
kanonicznego, co~zapewnia odtwarzalność raportowanych wyników.

Metody z~grupy heurystyk oraz~klasyczne klasyfikatory nie~są~oryginalnym
wkładem pracy, lecz~stanowią punkty odniesienia (\emph{baseline}) zaczerpnięte
z~literatury i~z~publicznych bibliotek. Wkładem własnym jest~ich~adaptacja
do~zadania detekcji zainfekowanych węzłów w~grafie lineage, autorski detektor
LineageDetector oraz~adapter grafu DUT. Pochodzenie i~zakres adaptacji
poszczególnych składników zestawiono poniżej:

\begin{itemize}
  \item \textbf{Heurystyki strukturalne}: inspirowane rodziną miar
        predykcji krawędzi w~grafach (Common Neighbors, Jaccard, Adamic–Adar,
        Preferential Attachment)~\cite{LibenNowell2007,Martinez2016}. Ponieważ graf
        przepływu danych jest~dwudzielny, miary oparte na~wspólnych sąsiadach
        są~w~nim nieskuteczne. Zaadaptowano je~więc~do~postaci scoringu
        pojedynczego węzła opartego na~stopniu i~roli w~DAG.
  \item \textbf{Random Forest}: implementacja z~biblioteki
        \texttt{scikit-learn}~\cite{scikit-learn}. Adaptacja: wejściem jest~
        autorski wektor cech topologicznych węzła (tab.~\ref{tab:cechy})
        zamiast cech pary obiektów, a~niezbalansowanie klas obsłużono przez
        \texttt{class\_weight}.
  \item \textbf{RUSBoost}: implementacja z~biblioteki
        \texttt{imbalanced-learn}~\cite{imblearn}, wybrana dla~bezpośredniego
        nawiązania do~pracy Boińskiego et~al.~\cite{Boinski2025}. Adaptacja:
        klasyfikacja węzła zamiast pary, na~cechach topologicznych.
  \item \textbf{LightGBM}: biblioteka gradientowego boostingu
        firmy Microsoft~\cite{Ke2017}. Adaptacja: klasyfikacja
        węzła na~cechach topologicznych, z~kompensacją rozkładu klas przez
        \texttt{class\_weight}. Hiperparametrów nie~strojono, a~wszystkie
        klasyfikatory porównywane są~w~konfiguracjach domyślnych
        (por.~sekcja~\ref{sec:trafnosc}).
  \item \textbf{LineageDetector}, wkład własny: łączy powyższe cechy
        z~autorskim relacyjnym score'em kompletności peer-consistency
        (sekcja~\ref{sec:lineage-detector}) i~kalibracją izotoniczną.
        Nie~ma~bezpośredniego odpowiednika w~literaturze.
  \item \textbf{Adapter grafu DUT}, wkład własny: przekształca benchmark
        Dutkiewicza et~al.~\cite{Dutkiewicz2026} z~wierszowej reprezentacji
        lineage do~modelu tabela–zadanie–\emph{DATA\_FLOW}
        (sekcja~\ref{sec:dane-dut}). Sam benchmark wraz ze~scenariuszami
        pochodzi z~pracy źródłowej.
\end{itemize}

Rozgraniczenie komponentów zapożyczonych, z~jawnie wskazanym źródłem,
od~elementów autorskich pozwala jednoznacznie ocenić rzeczywisty wkład
pracy i~odróżnić go~od~narzędzi ogólnego przeznaczenia. Ma~ono również
znaczenie dla~interpretacji wyników rozdziału~4. Klasyfikatory bazowe
i~LineageDetector korzystają z~tego samego wektora cech topologicznych,
tych samych implementacji bibliotecznych i~tego samego potoku ewaluacji.
Detektor różni się~od~klasyfikatora Random Forest, na~którym jest~oparty,
przede wszystkim elementami autorskimi, czyli score'em kompletności
i~kalibracją. Wpływ każdego z~tych elementów na~wynik zmierzono bezpośrednio
w~badaniu ablacyjnym (sekcja~\ref{sec:ablacja}), co~pozwala oddzielić
wkład własny od~skuteczności samych metod bibliotecznych.

% ============================================================

\chapter{Problem i~sposób jego rozwiązania}

Rozdział precyzuje formalną definicję problemu detekcji węzłów
o~zerwanym lineage (sekcja~\ref{sec:sformulowanie}). Sekcja~\ref{sec:dane}
charakteryzuje wykorzystane dane, czyli zbiór DLG-DG-23 oraz~zewnętrzny
graf DUT wraz z~adapterem, sekcja~\ref{sec:symulacja} opisuje
metodykę symulacji broken lineage przez~usuwanie zadań,
a~sekcja~\ref{sec:architektura} architekturę potoku przetwarzania.
Zaimplementowane metody detekcji omówiono w~sekcji~\ref{sec:algorytmy},
realizację techniczną w~sekcji~\ref{sec:implementacja}, testy jednostkowe
i~odtwarzalność wyników w~sekcji~\ref{sec:testy}, a~dostępność implementacji
i~pochodzenie algorytmów w~sekcji~\ref{sec:repozytorium}.

% ============================================================
\section{Sformułowanie problemu}
\label{sec:sformulowanie}
% ============================================================

Graf lineage modeluje się jako skierowany graf DAG przepływu danych
$G = (V, E)$, gdzie $V = V_{J} \cup V_{T} \cup V_{F}$ obejmuje węzły
typu zadanie (\emph{Data Job}), tabela (\emph{Data Table}) i~pole,
a~$E$ zawiera krawędzie przepływu danych (\emph{DATA\_FLOW}) między zadaniami
a~tabelami. Zjawisko broken lineage odpowiada usunięciu z~grafu pewnego
zadania $J$. Tabele z~nim incydentne, czyli jego poprzedniki (tabele
karmiące $J$) oraz~następniki (tabele karmione przez~$J$), tracą część
informacji o~swoim pochodzeniu i~stają się~\emph{zainfekowane}
(ang.~\emph{infected}).

Obserwowany graf ma~postać $G_{\mathrm{obs}} = G \setminus \{J_1,\dots,J_m\}$,
powstałą przez~usunięcie podzbioru zadań. Zadaniem detekcji jest wyznaczenie
funkcji, która każdej tabeli $t \in V_{T}$ przypisuje wynik
$s(t) \in [0, 1]$ interpretowany jako~oszacowane prawdopodobieństwo
$p(\text{tabela } t \text{ jest zainfekowana})$. Wynikiem nie~jest~decyzja
binarna przy~ustalonym progu, lecz~ranking tabel. Tabele o~najwyższym
wyniku są~pierwszymi kandydatami do~ręcznej inspekcji.

Etykiety referencyjne (\emph{ground truth}) generowane są~automatycznie.
Tabela jest zainfekowana wtedy i~tylko wtedy, gdy była incydentna
do~co~najmniej jednego usuniętego zadania. Zbiór kandydatów ograniczono
do~węzłów typu \emph{Data Table}. Węzły typu pole (\emph{field}) łączy
z~tabelami wyłącznie relacja \emph{PARENT\_CHILD}, dlatego nie~podlegają one detekcji.

Zasadnicza trudność zadania polega na~tym, że~osierocona tabela-następnik
po~usunięciu zadania wygląda jak~prawdziwa tabela źródłowa
($\mathrm{in\_df} = 0$). Rozróżnienie tych~dwóch przypadków wymaga cech
opisujących pozycję węzła w~grafie DAG oraz~uczenia indukcyjnego
(cross-graph), a~nie~samego stopnia węzła.

% ============================================================
\section{Zbiory danych}
\label{sec:dane}
% ============================================================

W~pracy wykorzystano dwa źródła danych o~odmiennym pochodzeniu i~przeznaczeniu.
Wszystkie eksperymenty główne, obejmujące trening i~ewaluację metod
w~protokole indukcyjnym, przeprowadzono na~rzeczywistych grafach zbioru
DLG-DG-23. Uzupełnieniem jest~zewnętrzny graf DUT, zbudowany na~podstawie
benchmarku Dutkiewicza et~al.~\cite{Dutkiewicz2026}, który służy wyłącznie
do~sprawdzenia przenośności modeli wytrenowanych na~DLG-DG-23. Oba źródła
opisano w~kolejnych podsekcjach.

\subsection{Zbiór DLG-DG-23}
\label{sec:dane-dlg}

Zbiór DLG-DG-23 pochodzi z~pracy Chen et~al.~\cite{Chen2024}
i~zawiera 18~grafów lineage pobranych ze~środowiska produkcyjnego
Huawei Cloud. Dane zostały zanonimizowane, a~nazwy tabel i~atrybutów
zastąpiono identyfikatorami w~formacie UUID. Schematów bazy danych
nie~udostępniono. Zbiór opublikowano w~2024~roku jako pierwszy
otwarty zestaw rzeczywistych grafów lineage.

Każdy graf zawiera węzły trzech typów: \emph{Data Job} (zapytanie SQL
lub~zadanie ETL), \emph{Data Table}
(tabela lub~widok przechowujący dane) oraz~\emph{field} (atrybut lub~kolumna
należąca do~tabeli). Łączą je~krawędzie dwóch typów: \emph{DATA\_FLOW}
(skierowana krawędź przepływu danych między tabelą a~zadaniem)
oraz~\emph{PARENT\_CHILD} (relacja przynależności pola do~tabeli). Pełny
graf jest~zatem \emph{heterogeniczny}. Przedmiotem detekcji są~wyłącznie
węzły typu \emph{Data Table}, a~sygnał czerpany jest z~podgrafu przepływu
danych (krawędzie \emph{DATA\_FLOW} między węzłami \emph{Data Table}
i~\emph{Data Job}). Podgraf ten~ma~strukturę dwudzielną i~skierowaną,
interpretowaną jako graf DAG. Grafy różnią się skalą, a~liczba węzłów waha się
od~278 do~17~085.

Zbiór DLG-DG-23 wybrano, ponieważ jest pierwszym publicznie udostępnionym
zestawem rzeczywistych grafów lineage~\cite{Chen2024}, a~w~trakcie prowadzenia
badań nie~zidentyfikowano innego zbioru tego rodzaju. Anonimizacja nazw
węzłów ma~bezpośredni wpływ na~dobór metod, ponieważ metryki podobieństwa
oparte na~nazwach obiektów (np.~miara Jaro–Winklera stosowana przez Boińskiego
et~al.~\cite{Boinski2025}) są~nieaplikowalne. W~niniejszej pracy sygnał
ograniczono zatem do~atrybutów topologicznych grafu. Ograniczenie
to~jest zarazem częścią pytania badawczego o~to, czy~sama struktura grafu
lineage wystarcza do~wskazania tabel o~zerwanym pochodzeniu danych.
Metoda skuteczna przy~tak ograniczonym wglądzie w~graf ma~przy~tym większą
szansę na~zastosowanie praktyczne niż~podejścia wymagające pełnych
metadanych.

Tabela~\ref{tab:dlg-stats} przedstawia statystyki poszczególnych
grafów zbioru DLG-DG-23.

\begin{table}[ht]
\centering
\caption{Charakterystyka grafów zbioru DLG-DG-23.}
\label{tab:dlg-stats}
\begin{threeparttable}
\begin{tabular}{llrl}
\toprule
\textbf{ID} & \textbf{Węzły} & \textbf{Krawędzie DATA\_FLOW} & \textbf{Skala} \\
\midrule
DLG1  &    298 &    24 & mały   \\
DLG2  &    464 &    37 & mały   \\
DLG3  &    603 &    57 & mały   \\
DLG4  &    415 &    32 & mały   \\
DLG5  & 1\,299 &   175 & średni \\
DLG6  & 13\,840 & 1\,443 & duży  \\
DLG7  &    574 &    43 & mały   \\
DLG8  &    546 &    34 & mały   \\
DLG9  &    756 &    62 & mały   \\
DLG10 & 5\,378 &   497 & duży   \\
DLG11 & 2\,024 &   121 & średni \\
DLG12 & 5\,645 &   469 & duży   \\
DLG13 & 2\,157 &   229 & średni \\
DLG14 &    453 &    24 & mały   \\
DLG15 &    558 &    41 & mały   \\
DLG16 & 5\,574 &   745 & duży   \\
DLG17 &    651 &    38 & mały   \\
DLG18 & 17\,085 & 1\,694 & duży \\
\bottomrule
\end{tabular}
\begin{tablenotes}
\small
\item Skala: mały $<$ 1\,000 węzłów, średni 1\,000–5\,000, duży $>$ 5\,000.
\end{tablenotes}
\end{threeparttable}
\end{table}

\subsection{Zewnętrzny graf DUT}
\label{sec:dane-dut}

Drugim źródłem danych jest~benchmark opracowany przez Dutkiewicza
et~al.~\cite{Dutkiewicz2026} w~zespole badawczym promotora
(sekcja~\ref{sec:literatura}). Obejmuje on~syntetyczne scenariusze broken
lineage, wygenerowane dla~przykładowej bazy \emph{Northwind}~\cite{Northwind}.
Lineage zapisany jest~w~nim na~poziomie pojedynczych wierszy, w~postaci grafu
wiedzy RDF oraz~plików CSV, w~których każdy rekord wiąże wartość źródłową
z~wartością docelową, która z~niej powstała. Jest to~zatem reprezentacja
drobnoziarnista (sekcja~\ref{sec:lineage_def}), niezgodna z~gruboziarnistym
modelem tabela–zadanie, na~którym operują metody opisane w~niniejszej pracy.

Aby~ocenić na~tych danych modele wytrenowane na~DLG-DG-23 bez~zmian w~ich
implementacji, opracowano adapter (moduł \texttt{dutkiewicz\_adapter.py}),
który agreguje zależności wierszowe do~poziomu tabel. Dla~każdej tabeli
docelowej tworzony jest~węzeł typu \emph{Data Job}, do~którego prowadzą
krawędzie \emph{DATA\_FLOW} od~wszystkich tabel źródłowych, z~których
pochodzą jej~wiersze. Z~węzła tego prowadzi następnie krawędź \emph{DATA\_FLOW}
do~samej tabeli docelowej. Kolumny kluczy występujące w~zapisie lineage
dodawane są~jako~węzły typu \emph{field}, połączone z~tabelą relacją
\emph{PARENT\_CHILD}. Wynikowy graf ma~te~same typy węzłów i~krawędzi
co~grafy DLG-DG-23, dzięki czemu przechodzi przez potok przetwarzania
identycznie jak~one.

Graf DUT liczy 21~tabel, 14~zadań i~33~krawędzie przepływu danych. Jest więc
o~rzędy wielkości mniejszy od~grafów DLG-DG-23 i~ma~odmienną, łańcuchową
topologię, w~której tabele rzadko współdzielą zadania. Adapter nie~odtwarza
przy~tym pełnych schematów tabel, ponieważ benchmark ich~nie~zawiera,
dlatego liczba pól jest~mniejsza niż~w~DLG-DG-23. Poprawność odwzorowania
weryfikuje sześć testów jednostkowych (tabela~\ref{tab:testy}). Graf DUT
nie~jest~wykorzystywany w~treningu żadnej z~metod, a~wyniki uzyskane na~nim
przedstawiono w~sekcji~\ref{sec:cross_dataset}.

% ============================================================
\section{Metodyka symulacji broken lineage przez usuwanie zadań}
\label{sec:symulacja}
% ============================================================

Ponieważ zbiór DLG-DG-23 zawiera w~pełni zaobserwowane grafy lineage,
zjawisko broken lineage symulowane jest przez~kontrolowane usuwanie
zadań (nie~pojedynczych krawędzi). Podejście to~odwzorowuje
przyczynę broken lineage, ponieważ zniknięcie tabeli tymczasowej lub~ukrycie
logiki przez UDF odpowiada utracie całego węzła-zadania wraz z~jego krawędziami.

\subsection{Procedura symulacji i~konstrukcja etykiet}
\label{sec:procedura-symulacji}

Kandydatami do~usunięcia są~\emph{zadania łączące}, czyli zadania mające
co~najmniej jedną wejściową i~jedną wyjściową krawędź \emph{DATA\_FLOW}
($\mathrm{in\_df} > 0 \land \mathrm{out\_df} > 0$). Usunięcie takiego zadania
zrywa lineage po~obu stronach, ponieważ tabele karmiące je~tracą konsumenta,
a~tabele przez~nie~karmione tracą producenta.

Instancję klasyfikacji węzłów buduje się~z~grafu $G$ w~kilku krokach.
Najpierw losowany jest~podzbiór zadań łączących przeznaczonych do~usunięcia,
którego udział reguluje parametr \texttt{removal\_ratio}. W~eksperymentach
przyjęto wartość 10\%, a~wpływ tego parametru zbadano w~analizie wrażliwości
(sekcja~\ref{sec:przeciek}). Zbiór tabel zainfekowanych wyznacza się~jako~sumę
tabel incydentnych do~usuniętych zadań. Graf obserwowany $G_{\mathrm{obs}}$
powstaje następnie przez~usunięcie wybranych zadań wraz z~ich krawędziami.
Kandydatami do~detekcji są~wszystkie tabele pozostałe w~$G_{\mathrm{obs}}$,
przy~czym etykieta~1 oznacza tabelę zainfekowaną, a~etykieta~0 tabelę czystą.

Dodatkowo zaimplementowano wariant kontrolny, w~którym każda instancja
usuwa dokładnie jedno zadanie, co~daje czystą interpretację
przyczynową (jedno zdarzenie broken lineage $=$ jedno brakujące zadanie).

\subsection{Kontrola przecieku – wykluczenie węzłów izolowanych}
\label{sec:exclude-isolated}

Usunięcie zadania może~pozbawić tabelę wszystkich krawędzi
przepływu danych ($\mathrm{in\_df} = 0 \land \mathrm{out\_df} = 0$).
Ponieważ w~oryginalnych grafach tabele izolowane nie~występują, trywialna
reguła „tabela izolowana $\Rightarrow$ zainfekowana'' pozwala rozpoznać
znaczną część pozytywów bez~żadnego uczenia i~sztucznie zawyża wynik.
Aby~zapewnić uczciwą ewaluację, wprowadzono flagę \texttt{exclude\_isolated},
która usuwa z~puli kandydatów tabele w~pełni izolowane. Wariant ten~stawia
trudniejsze i~bardziej realistyczne zadanie, polegające na~wykryciu tabel
o~częściowo zerwanym lineage, które nadal mają część połączeń. W~pracy raportowane
są~oba warianty, a~wariant uczciwy (z~wykluczeniem izolowanych) traktowany
jest jako~główny.

Tabela~\ref{tab:symulacja} przedstawia skalę symulacji dla~czterech wartości
parametru \texttt{removal\_ratio}. Łącznie w~18~grafach zbioru DLG-DG-23
znajduje się 1453~zadania łączące (od~3 do~404 na~graf), spośród których usuwana
jest~losowa część. Liczby podano dla~wariantu przed wykluczeniem tabel
izolowanych, ponieważ tylko wtedy można zmierzyć ich~udział wśród pozytywów.

\begin{table}[ht]
\centering
\caption[Charakterystyka symulacji broken lineage]{Charakterystyka symulacji broken lineage w~zależności
od~parametru \texttt{removal\_ratio}. Wartości uśrednione po~18~grafach
i~trzech ziarnach losowości (42--44). W~ostatniej kolumnie podano odsetek
tabel zainfekowanych, które utraciły wszystkie krawędzie przepływu
danych i~są~rozpoznawalne trywialną regułą izolacji.}
\label{tab:symulacja}
\begin{tabular}{ccccc}
\toprule
\texttt{removal\_ratio} & \textbf{Usuniętych} & \textbf{Tabel} & \textbf{Udział} & \textbf{Pozytywów} \\
 & \textbf{zadań} & \textbf{zainfekowanych} & \textbf{pozytywów} & \textbf{izolowanych} \\
\midrule
0{,}05 & 4{,}2 & 15{,}1 & 11{,}4\% & 19{,}7\% \\
0{,}10 & 8{,}0 & 25{,}1 & 19{,}0\% & 22{,}9\% \\
0{,}20 & 16{,}2 & 45{,}2 & 34{,}1\% & 29{,}8\% \\
0{,}30 & 24{,}3 & 62{,}6 & 47{,}3\% & 34{,}5\% \\
\bottomrule
\end{tabular}
\end{table}

Dwie obserwacje z~tabeli mają~znaczenie dla~konstrukcji eksperymentu.
Przyjęta w~biegu kanonicznym wartość $0{,}10$ daje średnio 25~tabel
zainfekowanych na~graf przy~19\% pozytywów. Zadanie pozostaje więc~wyraźnie
niezrównoważone, ale~klasa pozytywna jest~dostatecznie liczna, by~metryki
rankingowe były wiarygodne. Udział pozytywów trywialnych (tabel w~pełni
izolowanych) rośnie z~kolei z~19{,}7\% do~34{,}5\% wraz ze~wzrostem
odsetka usuwanych zadań. Im~silniejsze zaburzenie, tym~większa część zadania
daje się~rozwiązać bez~żadnego uczenia. Liczby te~określają skalę przecieku,
któremu zapobiega opisana wyżej flaga \texttt{exclude\_isolated}. Przy~biegu
kanonicznym usuwa ona~z~puli kandydatów blisko jedną czwartą pozytywów,
i~to~tych najłatwiejszych.

% ============================================================
\section{Architektura rozwiązania}
\label{sec:architektura}
% ============================================================

Potok przetwarzania składa się z~czterech komponentów: modułu wczytywania
grafów (wraz z~adapterem grafu DUT), modułu symulacji broken lineage (budowa
instancji klasyfikacji węzłów), metod detekcji oraz~modułu ewaluacji indukcyjnej.

Rysunek~\ref{rys:potok} przedstawia schemat potoku przetwarzania, od~wczytania
grafów, przez symulację broken lineage i~ekstrakcję cech, po~scoring detektora
i~wyznaczenie metryk rankingowych. Oba źródła danych trafiają do~wspólnego
potoku. Grafy DLG-DG-23 wczytywane są~bezpośrednio z~plików JSON, natomiast
graf DUT przechodzi najpierw przez adapter, który sprowadza go~do~tego samego
modelu danych. Od~etapu symulacji oba rodzaje grafów przetwarzane są~identycznie,
różni je~jedynie rola w~ewaluacji.

\begin{figure}[ht]
  \centering
  \includegraphics[width=\linewidth]{rys_potok.pdf}
  \caption[Schemat potoku przetwarzania]{Schemat potoku przetwarzania. Grafy
           DLG-DG-23 oceniane są~w~protokole indukcyjnym leave-one-graph-out
           (trening na~17~grafach, test na~grafie niewidzianym). Graf DUT,
           sprowadzony adapterem do~modelu tabela–zadanie, służy wyłącznie
           do~testu modeli wytrenowanych na~DLG-DG-23.}
  \label{rys:potok}
\end{figure}

\subsection{Moduł wczytywania danych}

Dane wejściowe każdego grafu przechowywane są~w~dwóch plikach JSON:
\texttt{DLGi-node.json} (opisy węzłów z~atrybutem \texttt{asset\_type})
oraz~\texttt{DLGi-edge.json} (krawędzie z~atrybutami \texttt{relation\_type}
i~\texttt{relation\_id}). Moduł \texttt{loader.py} buduje z~nich obiekt
\texttt{nx.DiGraph} (NetworkX) i~zwraca słownik wszystkich 18~grafów.
Graf DUT wczytywany jest~przez moduł \texttt{dutkiewicz\_adapter.py}, który
zwraca obiekt tego samego typu, z~atrybutami węzłów i~krawędzi zgodnymi
z~konwencją DLG-DG-23 (sekcja~\ref{sec:dane-dut}). Dalsze moduły potoku
nie~rozróżniają więc~źródła grafu.

\subsection{Moduł symulacji usuwania zadań}

Moduł \texttt{job\_removal.py} realizuje metodykę z~sekcji~\ref{sec:symulacja}:
identyfikuje zadania łączące, usuwa ich podzbiór, wyznacza zbiór zainfekowanych
tabel i~zwraca instancję (graf $G_{\mathrm{obs}}$, listę kandydatów,
etykiety, listę usuniętych zadań). Obsługuje warianty \texttt{side},
określający stronę incydencji, po~której tabele uznaje się~za~zainfekowane
(strona producenta, konsumenta lub~obie), oraz~\texttt{exclude\_isolated}. Funkcja \texttt{inductive\_split} dzieli
identyfikatory grafów na~rozłączne zbiory treningowy i~testowy (cross-graph).

\subsection{Moduł ewaluacji i~metryki rankingowe}

Moduł \texttt{node\_metrics.py} wyznacza metryki rankingowe: AUC-ROC,
AUC-PR, Precision@k, Recall@k, Hits@k oraz~Brier score. Wyniki agregowane
są~po~grafach testowych w~protokole indukcyjnym. Dla~stabilności oszacowań
stosuje się powtórzenia z~różnymi ziarnami i~trymowanie wartości skrajnych
(min/max).

Ewaluację prowadzi się~w~protokole indukcyjnym cross-graph.
Model trenuje się~na~podzbiorze grafów, a~testuje na~grafach niewidzianych
(leave-one-graph-out lub~podział zbioru 18~grafów). Zbiory treningowy
i~testowy są~rozłączne, co~sprawdza zdolność metody do~generalizacji
na~nowe grafy, a~nie~jedynie do~zapamiętania jednego. Konsekwencją tego
protokołu jest~wykluczenie metody Node2Vec~\cite{Grover2016}, która uczy zanurzeń
węzłów transduktywnie, dla~konkretnego grafu, przez~co~embeddingi
są~geometrycznie nieporównywalne między grafami~\cite{Hamilton2017}.
Obserwacja ta~motywuje zastosowanie indukcyjnych
grafowych sieci neuronowych (GraphSAGE~\cite{Hamilton2017},
R-GCN~\cite{Schlichtkrull2018}, HGT~\cite{Hu2020}), wskazanych w~rozdziale~5 jako~kierunek dalszych badań.

% ============================================================
\section{Zaimplementowane metody detekcji}
\label{sec:algorytmy}
% ============================================================

Zaimplementowano trzy grupy metod detekcji zainfekowanych tabel, różniące się
ilością wykorzystywanej informacji i~stopniem złożoności. Heurystyki strukturalne
(sekcja~\ref{sec:heurystyki}) przypisują tabeli wynik na~podstawie prostych
reguł i~nie~wymagają uczenia, dlatego wyznaczają dolny punkt odniesienia.
Klasyczne klasyfikatory uczenia maszynowego (sekcja~\ref{sec:ml}) uczą się
na~grafach treningowych zależności między wektorem cech topologicznych
a~etykietą, co~pozwala ocenić, ile~wnosi samo uczenie indukcyjne. Autorski
detektor LineageDetector (sekcja~\ref{sec:lineage-detector}) rozszerza ten~wektor
o~relacyjny score kompletności i~uzupełnia model o~kalibrację prawdopodobieństw.

Wszystkie metody mają wspólny interfejs. Dla~każdej tabeli-kandydata zwracają
wynik z~przedziału $[0, 1]$, na~podstawie którego budowany jest~ranking. Dzięki
temu każda z~nich jest~oceniana w~tym samym potoku, na~tych samych instancjach
i~tymi samymi metrykami, a~różnice wyników można przypisać samym metodom.

% ============================================================
\subsection{Heurystyki strukturalne}
\label{sec:heurystyki}
% ============================================================

Heurystyki strukturalne przypisują tabeli wynik podejrzliwości wyłącznie
na~podstawie jej pozycji w~grafie, bez~etapu uczenia. Są~naturalnym
punktem odniesienia (\emph{baseline}) dla~metod uczonych. Miary oparte
na~wspólnych sąsiadach (Common Neighbors, Jaccard,
Adamic-Adar)~\cite{LibenNowell2007} są~w~grafie
dwudzielnym nieskuteczne, ponieważ węzły tego~samego typu nigdy nie~mają
wspólnych sąsiadów bezpośrednich. Z~tego powodu wykluczono je~z~implementacji.

W~pracy zaimplementowano cztery heurystyki bazowe. Piątym punktem
odniesienia w~tej~grupie jest~sam \emph{score kompletności}
(sekcja~\ref{sec:lineage-detector}), oceniany dodatkowo w~oderwaniu
od~detektora, co~pozwala rozdzielić wkład sygnału relacyjnego od~wkładu
uczenia. W~zestawieniach rozdziału~4 występuje on~zatem jako~osobna metoda
obok czterech poniższych:

\begin{itemize}
  \item \textbf{Anomalia stopnia} (\texttt{degree\_anomaly}): wynik
        $1/(1+\deg(t))$, gdzie $\deg(t)$ jest stopniem całkowitym tabeli.
        Heurystyka opiera się~na~założeniu, że~tabela, która utraciła
        zadanie, ma~mniej połączeń niż~przed zerwaniem lineage.
  \item \textbf{Brzeg grafu} (\texttt{boundary}): suma wskaźników
        $\mathbf{1}[\mathrm{in\_df}=0] + \mathbf{1}[\mathrm{out\_df}=0]$.
        Tabele na~brzegu przepływu (korzenie i~liście) są~podejrzane,
        ponieważ zerwanie lineage często przesuwa tabelę na~brzeg grafu.
  \item \textbf{Niska łączność z~zadaniami}
        (\texttt{low\_job\_connectivity}): wynik $1/(1+n_{J}(t))$,
        gdzie $n_{J}(t)$ to~liczba unikatowych zadań sąsiadujących
        z~tabelą przez krawędzie \emph{DATA\_FLOW}.
  \item \textbf{Reguła korzenia} (\texttt{rule\_root}): binarny wskaźnik
        $\mathbf{1}[\mathrm{in\_df}=0]$, zgodnie z~którym każda tabela
        bez~producenta jest uznawana za~podejrzaną. To~celowo trywialny punkt odniesienia,
        obrazujący siłę najprostszego możliwego symptomu.
\end{itemize}

Heurystyki te~mierzą wyłącznie lokalny symptom (stopień, rola).
Ich~ograniczenie, przede wszystkim nieodróżnianie osieroconej tabeli
od~prawdziwego źródła, motywuje metody uczone oraz~autorski detektor
LineageDetector (sekcja~\ref{sec:lineage-detector}), który wzbogaca opis
węzła o~relacyjny score kompletności (peer-consistency).

% ============================================================
\subsection{Klasyczne metody uczenia maszynowego}
\label{sec:ml}
% ============================================================

Klasyczne metody ML traktują detekcję jako~binarną klasyfikację węzła, w~której
wejściem jest wektor cech opisujących pozycję tabeli w~grafie DAG,
a~wyjściem prawdopodobieństwo zainfekowania. Ze~względu na~anonimizację
(UUID) jako~cechy przyjęto wyłącznie atrybuty topologiczne, ponieważ metryki
oparte na~nazwach są~nieaplikowalne.

\subsubsection{Wektor cech węzła}

Tabela~\ref{tab:cechy} przedstawia pełny wektor 13~cech topologicznych
wyznaczanych dla~każdej tabeli-kandydata. Cechy opisują pozycję węzła
w~grafie przepływu danych, bo~to~ona, a~nie~sam stopień, pozwala odróżnić
prawdziwą tabelę źródłową od~tabeli osieroconej po~usunięciu zadania.

\begin{table}[ht]
\centering
\caption[Wektor cech topologicznych węzła]{Wektor cech topologicznych węzła (moduł \texttt{node\_features}).}
\label{tab:cechy}
\small
\begin{tabular}{cl>{\raggedright\arraybackslash}p{5.2cm}>{\raggedright\arraybackslash}p{4.6cm}}
\toprule
\textbf{Nr} & \textbf{Cecha} & \textbf{Opis} & \textbf{Rola w~detekcji} \\
\midrule
1 & \texttt{in\_df} & Liczba wejściowych krawędzi \emph{DATA\_FLOW}
    (od~zadań) & Brak producenta ($=0$) to~główny symptom zerwania \\
2 & \texttt{out\_df} & Liczba wyjściowych krawędzi \emph{DATA\_FLOW}
    (do~zadań) & Symetryczny symptom, czyli utrata konsumenta \\
3 & \texttt{degree} & Stopień całkowity węzła (wszystkie typy krawędzi)
    & Ogólna łączność, spadająca po~usunięciu zadania \\
4 & \texttt{role} & Rola w~przepływie (źródło, pośrednia, ujście,
    poza przepływem) & Pozycja funkcjonalna tabeli w~DAG \\
5 & \texttt{flow\_depth} & Głębokość w~DAG (odległość od~źródeł)
    & Osierocona tabela ma~sztucznie małą głębokość \\
6 & \texttt{flow\_reach} & Zasięg w~dół DAG (odległość do~ujść)
    & Symetrycznie sztucznie mały zasięg \\
7 & \texttt{depth\_plus\_reach} & Suma głębokości i~zasięgu
    & Długość ścieżki przez węzeł, wykrywa „skrócone'' łańcuchy \\
8 & \texttt{is\_root} & Wskaźnik $\mathbf{1}[\mathrm{in\_df}=0]$
    & Binarna wersja symptomu utraty producenta \\
9 & \texttt{is\_leaf} & Wskaźnik $\mathbf{1}[\mathrm{out\_df}=0]$
    & Binarna wersja symptomu utraty konsumenta \\
10 & \texttt{n\_field\_children} & Liczba pól tabeli (krawędzie
    \emph{PARENT\_CHILD}) & Rozmiar tabeli, sygnał niezależny
    od~przepływu \\
11 & \texttt{n\_job\_neighbors} & Liczba unikatowych zadań sąsiadujących
    & Łączność z~warstwą zadań \\
12 & \texttt{avg\_job\_in\_df} & Średni stopień wejściowy sąsiednich zadań
    & Charakterystyka otoczenia węzła \\
13 & \texttt{avg\_job\_out\_df} & Średni stopień wyjściowy sąsiednich
    zadań & Jak~wyżej, od~strony wyjść \\
\bottomrule
\end{tabular}
\end{table}

Detektor LineageDetector (sekcja~\ref{sec:lineage-detector})
rozszerza ten wektor o~trzy cechy kompletności
(\texttt{peer\_producer\_score}, \texttt{peer\_consumer\_score},
\texttt{completeness\_anomaly}), opisane w~sekcji~\ref{sec:lineage-detector}.

\subsubsection{Random Forest}

Random Forest (RF)~\cite{Breiman2001} to~metoda uczenia zespołowego oparta na~uśrednianiu
predykcji drzew decyzyjnych trenowanych na~losowych podpróbkach danych
i~cech (\emph{bagging}). Dla~obsługi niezrównoważonych klas zastosowano
\texttt{class\_weight=\textquotedbl{}balanced\textquotedbl{}}.

\subsubsection{RUSBoost}

RUSBoost (\emph{Random Under-Sampling Boosting})~\cite{Seiffert2010} łączy boosting AdaBoost
z~losowym podpróbkowaniem klasy większościowej przed~każdą iteracją.
Zaimplementowano go~dla~bezpośredniego nawiązania do~wyników Boińskiego
et~al.~\cite{Boinski2025}, przy~czym w~niniejszej pracy klasyfikuje on~węzły,
a~nie~pary obiektów.

\subsubsection{LightGBM}

LightGBM~\cite{Ke2017} to~implementacja gradientowego boostingu drzew decyzyjnych,
budująca kolejne drzewa sekwencyjnie, każde minimalizując błąd resztkowy
poprzednika. Niezbalansowanie klas kompensuje parametr
\texttt{class\_weight}~$=$~\texttt{balanced}, przypisujący obserwacjom wagi
odwrotnie proporcjonalne do~liczności klasy.

% ============================================================
\subsection{Autorski detektor LineageDetector}
\label{sec:lineage-detector}
% ============================================================

Autorski detektor \texttt{LineageDetector} (klasa
\texttt{LineageHealthDetector}) łączy w~jednym skalibrowanym modelu trzy
elementy. Podstawę stanowią cechy pozycji w~DAG, wyznaczane przez moduł
\texttt{node\_features} i~opisujące węzeł lokalnie. Uzupełnia je~score
kompletności \emph{peer-consistency}, opisany w~dalszej części sekcji,
który wnosi sygnał relacyjny i~stanowi główny element wkładu własnego pracy.
Całość zamyka kalibracja prawdopodobieństw metodą izotoniczną
(\texttt{CalibratedClassifierCV}), dzięki której wynik detektora można
interpretować jako~oszacowanie $p(\text{zainfekowana})$, a~nie~jedynie
jako~wartość porządkującą ranking.

\subsubsection{Score kompletności peer-consistency}

Zaproponowany w~pracy score kompletności opiera się na~założeniu,
że~tabela o~zerwanym lineage jest~\emph{strukturalnie niespójna}
ze~swoimi sąsiadami (peerami) w~grafie DAG. Dla~tabeli-korzenia
($\mathrm{in\_df} = 0$) rozważa się~zbiór peerów-wejść, czyli inne tabele
wpadające do~tych~samych zadań downstream. Jeśli~peerzy mają
producenta, a~rozpatrywana tabela nie, to~jej brak producenta jest~podejrzany,
ponieważ prawdziwe źródło danych sąsiaduje zwykle z~innymi źródłami. Symetrycznie
dla~tabeli-liścia ($\mathrm{out\_df} = 0$) i~jej peerów po~stronie konsumpcji.
Formalnie dla~tabeli $t$ definiuje się:
%
\[
  \mathrm{prod}(t) = \frac{1}{|P(t)|}\sum_{p \in P(t)} \mathbf{1}[\mathrm{in\_df}(p) > 0],
  \qquad
  \mathrm{cons}(t) = \frac{1}{|C(t)|}\sum_{p \in C(t)} \mathbf{1}[\mathrm{out\_df}(p) > 0],
\]
%
gdzie $P(t)$ i~$C(t)$ oznaczają zbiory peerów odpowiednio po~stronie
producenta i~konsumenta, a~$\mathbf{1}[\cdot]$ funkcję wskaźnikową.
Bramkowany sygnał anomalii aktywuje się~wyłącznie na~brzegu grafu:
%
\[
  \mathrm{anomaly}(t) =
  \max\bigl(\mathbf{1}[\mathrm{in\_df}(t)=0]\cdot\mathrm{prod}(t),\;
            \mathbf{1}[\mathrm{out\_df}(t)=0]\cdot\mathrm{cons}(t)\bigr).
\]
%
Sygnał ten~ma~charakter relacyjny. Różni się~jakościowo od~lokalnego
stopnia węzła i~dostarcza informacji niedostępnej klasyfikatorowi
korzystającemu wyłącznie z~cech lokalnych. Definicja jest~symetryczna
względem obu~stron przepływu, jednak na~zbiorze DLG-DG-23 aktywna
okazuje się~wyłącznie strona producenta, ponieważ każde zadanie w~tym~zbiorze
zapisuje dokładnie jedną tabelę, przez~co~zbiór $C(t)$ jest~zawsze pusty.
Konsekwencje tej~własności omówiono w~sekcji~\ref{sec:waznosc}. Trzy wynikające z~niego cechy
(\texttt{peer\_producer\_score}, \texttt{peer\_consumer\_score},
\texttt{completeness\_anomaly}) rozszerzają wektor cech z~tabeli~\ref{tab:cechy}
i~są~dostępne wyłącznie detektorowi LineageDetector.

Bazowym estymatorem jest~domyślnie Random Forest, a~wariantem
alternatywnym LightGBM. Interfejs \texttt{fit}/\texttt{score} jest~zgodny z~pozostałymi
klasyfikatorami, co~umożliwia bezpośrednie porównanie w~tym~samym potoku.
Konstrukcja detektora umożliwia badania ablacyjne. Wyłączenie kalibracji
(\texttt{calibrate=False}) izoluje jej~wpływ na~Brier score, a~wyłączenie
cech kompletności (\texttt{use\_completeness=False}) wpływ sygnału
relacyjnego na~AUC.

% ============================================================
\section{Implementacja}
\label{sec:implementacja}
% ============================================================

Sekcja opisuje realizację techniczną potoku z~sekcji~\ref{sec:architektura}:
zastosowane narzędzia, organizację kodu oraz~te~fragmenty implementacji,
które są~istotne dla~zrozumienia metody, a~nie~wynikają wprost z~jej~opisu
formalnego.

\subsection{Stos technologiczny}

Całość zaimplementowano w~języku Python. Tabela~\ref{tab:stos} zestawia
wykorzystane biblioteki wraz z~wersjami, w~których wykonano bieg kanoniczny.
Podanie wersji jest~istotne dla~odtwarzalności, ponieważ implementacje algorytmów
uczenia maszynowego bywają modyfikowane między wydaniami, co~może~wpływać
na~wyniki liczbowe.

\begin{table}[ht]
\centering
\caption{Stos technologiczny i~wersje bibliotek użyte w~biegu kanonicznym.}
\label{tab:stos}
\begin{tabular}{lll}
\toprule
\textbf{Składnik} & \textbf{Wersja} & \textbf{Rola w~pracy} \\
\midrule
Python~\cite{Python}                 & 3.10.2  & Środowisko uruchomieniowe \\
NetworkX~\cite{NetworkX}             & 3.4.2   & Reprezentacja i~przetwarzanie grafów \\
NumPy~\cite{NumPy}                   & 1.26.4  & Operacje macierzowe, wektory cech \\
pandas~\cite{pandas}                 & 1.5.1   & Agregacja wyników, eksport CSV \\
scikit-learn~\cite{scikit-learn}     & 1.6.1   & Random Forest, kalibracja izotoniczna, metryki \\
LightGBM~\cite{Ke2017}               & 4.6.0   & Gradientowy boosting drzew \\
imbalanced-learn~\cite{imblearn}     & 0.14.1  & RUSBoost \\
SciPy~\cite{SciPy}                   & 1.15.3  & Test kolejności par Wilcoxona \\
Matplotlib~\cite{Matplotlib}         & 3.8.0   & Rysunki \\
pytest~\cite{pytest}                 & 9.0.2   & Testy jednostkowe \\
\bottomrule
\end{tabular}
\end{table}

\subsection{Organizacja repozytorium}

Kod podzielono na~pakiety odpowiadające etapom potoku, a~skrypty
uruchomieniowe oddzielono od~modułów bibliotecznych, aby~każdy eksperyment
z~rozdziału~4 dawał się~odtworzyć jednym poleceniem. Listing~\ref{lst:repo}
przedstawia strukturę katalogów.

\begin{lstlisting}[caption={Struktura repozytorium (fragmenty istotne
dla~detekcji węzłów).},label={lst:repo},numbers=none,language={},float=htbp]
src/
  data/
    loader.py               # wczytanie 18 grafow DLG-DG-23 z JSON
    dutkiewicz_adapter.py   # adapter benchmarku zewnetrznego (graf DUT)
  detection/
    job_removal.py          # symulacja broken lineage, etykiety, podzial
    node_features.py        # 13 cech topologicznych wezla
    completeness.py         # score peer-consistency (wklad wlasny)
    lineage_detector.py     # LineageHealthDetector
    node_classifier.py      # heurystyki + klasyczne klasyfikatory
    node_metrics.py         # metryki rankingowe
    run_detection_experiments.py   # bieg kanoniczny (rozdz. 4.2)
    run_ablation.py                # badanie ablacyjne (rozdz. 4.5)
    run_cross_dataset.py           # test cross-dataset (rozdz. 4.9)
    run_removal_sensitivity.py     # analiza wrazliwosci (rozdz. 4.6)
scripts/
  test_istotnosci.py        # testy Wilcoxona (rozdz. 4.4)
  tabela_symulacji.py       # tabela charakterystyki symulacji
tests/                      # 83 testy jednostkowe (pytest)
results/                    # wyniki biegow w formacie CSV
\end{lstlisting}

\subsection{Format danych i~wczytywanie}

Każdy graf zbioru DLG-DG-23 opisany jest~parą plików JSON: lista węzłów
z~atrybutem \texttt{asset\_type} oraz~lista krawędzi z~atrybutami
\texttt{relation\_id} i~\texttt{relation\_type} (listing~\ref{lst:json}).
Moduł \texttt{loader.py} przekształca je~w~obiekt \texttt{nx.DiGraph}
i~przenosi typ węzła oraz~typ relacji do~atrybutów grafu. Dzięki temu
pozostałe moduły filtrują podgraf przepływu danych bez~ponownego sięgania
do~plików.

\begin{lstlisting}[caption={Fragment plików wejściowych
\texttt{DLGi-node.json} i~\texttt{DLGi-edge.json} (identyfikatory skrócone).},
label={lst:json},numbers=none,language={},float=htbp]
// DLG1-node.json
{"nodes": [
  {"asset_id": "3f2a...c81", "asset_type": "Data Table"},
  {"asset_id": "9b07...4de", "asset_type": "Data Job"},
  {"asset_id": "a154...77f", "asset_type": "Data Field"}
]}

// DLG1-edge.json
{"edges": [
  {"relation_id": "r-0012", "relation_type": "DATA_FLOW",
   "source": "3f2a...c81", "target": "9b07...4de"},
  {"relation_id": "r-0013", "relation_type": "PARENT_CHILD",
   "source": "3f2a...c81", "target": "a154...77f"}
]}
\end{lstlisting}

% Listingi 3.1 i 3.2 maja zostac zlozone przed ta podsekcja, aby nie
% rozcinaly jej pierwszego akapitu.
\FloatBarrier
\subsection{Wyznaczanie score'u peer-consistency}

Najistotniejszym fragmentem implementacji jest~wyznaczanie relacyjnego
score'u kompletności, zdefiniowanego wzorami w~sekcji~\ref{sec:lineage-detector}.
Listing~\ref{lst:peer} pokazuje jego~rdzeń. Sam wzór nie~precyzuje, że~zbiory
peerów budowane są~przez zadania. Peerami tabeli
po~stronie producenta są~inne tabele wpadające do~tych~samych zadań
downstream, w~grafie dwudzielnym nie~istnieje bowiem bezpośrednia krawędź
tabela–tabela. Widoczne w~listingu bramkowanie wskaźnikami odpowiada wprost
drugiemu ze~wzorów z~sekcji~\ref{sec:lineage-detector}. Bez~niego score
mierzyłby jedynie własności otoczenia, a~nie~niespójność tabeli
z~jej~otoczeniem.

\begin{lstlisting}[caption={Rdzeń wyznaczania score'u peer-consistency
(moduł \texttt{completeness.py}).},label={lst:peer},language=Python,float=htbp]
# Strona producenta: peerzy = inne wejscia do downstream-jobow tej tabeli.
downstream_jobs = [v for _, v, d in G.out_edges(t, data=True)
                   if d.get("relation_type") == "DATA_FLOW"
                   and G.nodes[v].get("asset_type") == "Data Job"]
prod_peers = set()
for j in downstream_jobs:
    prod_peers.update(job_inputs.get(j, ()))
prod_peers.discard(t)
prod_score = (np.mean([1.0 if in_df[p] > 0 else 0.0 for p in prod_peers])
              if prod_peers else 0.0)

# Bramkowana anomalia: liczy sie tylko na brzegu (korzen/lisc).
anomaly = 0.0
if in_df[t] == 0:
    anomaly = max(anomaly, prod_score)
if out_df[t] == 0:
    anomaly = max(anomaly, cons_score)
\end{lstlisting}

\FloatBarrier
\subsection{Protokół indukcyjny w~kodzie}

Pętla ewaluacji realizuje schemat leave-one-graph-out, w~którym dla~każdego grafu
testowego model uczy się~na~wszystkich pozostałych. Listing~\ref{lst:logo}
przedstawia jej~szkielet. Kluczowa dla~poprawności protokołu jest~kolejność
operacji. Instancje wszystkich grafów budowane są~przed pętlą,
przy~wspólnym ziarnie, a~pula treningowa powstaje przez~wykluczenie grafu
testowego z~gotowego słownika. Model nie~ma~więc~dostępu do~żadnej
informacji z~grafu testowego, także pośrednio przez~dobór usuwanych zadań.

\begin{lstlisting}[caption={Szkielet pętli leave-one-graph-out
(moduł \texttt{run\_detection\_experiments.py}).},label={lst:logo},
language=Python,float=htbp]
datasets = build_datasets(graphs, removal_ratio, seed, side,
                          exclude_isolated=exclude_isolated)
for test_g in datasets:
    ds_test = datasets[test_g]
    train_pool = [datasets[g] for g in datasets if g != test_g]

    for name in FIT_MODELS:
        clf = _make_model(name, seed).fit(train_pool)
        m = node_detection_metrics(ds_test["labels"], clf.score(ds_test))
        rows.append({"graph": test_g, "algorithm": name, **m})
\end{lstlisting}

Wyniki z~kolejnych ziaren agregowane są~średnią z~trymowaniem wartości
skrajnych (odrzucenie najmniejszej i~największej), co~ogranicza wpływ
pojedynczego nietypowego losowania na~oszacowanie.

\subsection{Złożoność obliczeniowa}

Koszt przetwarzania jednego grafu jest~liniowy względem liczby krawędzi.
Wyznaczenie stopni \emph{DATA\_FLOW}, ekstrakcja cech oraz~budowa zbiorów
peerów wymagają stałej liczby przejść po~krawędziach, a~metadane przepływu
(głębokość i~zasięg w~DAG) wyznaczane są~przeszukiwaniem grafu, również
w~czasie $O(|V| + |E|)$. Dominującym składnikiem kosztu nie~jest~zatem
przetwarzanie grafu, lecz~trenowanie klasyfikatorów. W~protokole
leave-one-graph-out model uczony jest~od~nowa dla~każdego grafu testowego,
co~przy~18~grafach, trzech ziarnach i~czterech wariantach detektora daje
216~pełnych treningów. Bieg ablacyjny z~sekcji~\ref{sec:ablacja} zajął
z~tego~powodu 412~sekund, podczas gdy sama ekstrakcja cech dla~wszystkich
18~grafów trwa poniżej sekundy. Skala zbioru DLG-DG-23 (do~17~085~węzłów
na~graf) nie~stanowi więc~ograniczenia wydajnościowego.

% ============================================================
\section{Testy jednostkowe i~odtwarzalność}
\label{sec:testy}
% ============================================================

\subsection{Zakres testów}

Poprawność implementacji zweryfikowano zestawem 83~testów jednostkowych
uruchamianych narzędziem \texttt{pytest}. W~wersji implementacji
opisanej w~pracy wszystkie testy kończą się~powodzeniem.
Tabela~\ref{tab:testy} przedstawia ich~rozkład na~moduły.

Poza poprawnością techniczną testy sprawdzają dwie własności istotne
metodycznie. Pierwszą jest~brak przecieku między zbiorem treningowym
a~testowym, weryfikowany przez sprawdzenie, czy~identyfikatory grafów w~obu
zbiorach są~rozłączne. Drugą jest~determinizm symulacji, który oznacza,
że~przy~ustalonym ziarnie losowości powtórne wykonanie daje identyczny zbiór
usuniętych zadań i~identyczne etykiety. Bez~tej własności liczby raportowane
w~rozdziale~4 nie~byłyby~odtwarzalne.

\begin{table}[ht]
\centering
\caption{Rozkład testów jednostkowych na~moduły.}
\label{tab:testy}
\small
\begin{tabular}{lr>{\raggedright\arraybackslash}p{7.2cm}}
\toprule
\textbf{Plik testowy} & \textbf{Liczba} & \textbf{Weryfikowane własności} \\
\midrule
\texttt{test\_loader.py} & 29 &
  Poprawność typów węzłów i~krawędzi, zgodność liczby obiektów z~plikami
  źródłowymi, wczytanie wszystkich 18~grafów \\
\texttt{test\_splitter.py} & 17 &
  Rozłączność podziałów, zgodność liczności podzbiorów \\
\texttt{test\_job\_removal.py} & 12 &
  Identyfikacja zadań łączących, poprawność etykiet, rozłączność zbiorów
  treningowego i~testowego, działanie flagi \texttt{exclude\_isolated} \\
\texttt{test\_scenario\_splitter.py} & 9 &
  Scenariusze usuwania krawędzi (badanie wstępne) \\
\texttt{test\_dutkiewicz\_adapter.py} & 6 &
  Mapowanie benchmarku Dutkiewicza et~al.\ na~model
  tabela–zadanie–\emph{DATA\_FLOW} (graf DUT) \\
\texttt{test\_node\_features.py} & 4 &
  Wymiar wektora cech, zgodność z~\texttt{feature\_names()},
  wartości brzegowe dla~węzłów izolowanych \\
\texttt{test\_completeness.py} & 3 &
  Poprawność score'u peer-consistency, bramkowanie sygnału anomalii \\
\texttt{test\_lineage\_detector.py} & 3 &
  Interfejs \texttt{fit}/\texttt{score}, zakres zwracanych
  prawdopodobieństw, warianty ablacyjne \\
\midrule
\textbf{Razem} & \textbf{83} & \\
\bottomrule
\end{tabular}
\end{table}

\subsection{Odtwarzalność wyników}

Wszystkie liczby przedstawione w~rozdziale~4 pochodzą z~jednego
\emph{biegu kanonicznego} o~ustalonej konfiguracji: 18~grafów,
\texttt{removal\_ratio}~$=0{,}10$, \texttt{side}~$=$~\texttt{both},
trzy powtórzenia z~ziarnami 42--44, protokół indukcyjny leave-one-graph-out,
próg wiarygodności $\mathrm{MIN\_POS} = 5$, wariant uczciwy
(\texttt{exclude\_isolated}). Wyniki pośrednie zapisywane są~w~postaci plików
CSV w~katalogu \texttt{results/}, osobno dla~wariantu uczciwego
i~zawyżonego, wraz z~opisem konfiguracji. Eksperymenty poboczne
(ablacja, analiza wrażliwości, test cross-dataset) mają~własne skrypty
i~własne pliki wynikowe, wymienione w~listingu~\ref{lst:repo}, co~pozwala
odtworzyć każdą tabelę rozdziału~4 niezależnie od~pozostałych.

% ============================================================
\section{Dostępność implementacji i~pochodzenie algorytmów}
\label{sec:repozytorium}
% ============================================================

Całość implementacji opisanej w~niniejszej pracy, obejmującą moduły
wczytywania danych, symulacji broken lineage, ekstrakcji cech, metod
detekcji oraz~ewaluacji, wraz z~testami jednostkowymi i~skryptami
odtwarzającymi eksperymenty z~rozdziału~4, udostępniono w~publicznym
repozytorium Git
pod~adresem \texttt{https://github.com/Catcuss/magisterka-data-lineage}
(dostęp zweryfikowano 19~sierpnia 2026~r.). Repozytorium zawiera instrukcję
uruchomienia (dodatek~\ref{dod:uruchomienie}) oraz~pliki konfiguracyjne biegu
kanonicznego, co~zapewnia odtwarzalność raportowanych wyników.

Metody z~grupy heurystyk oraz~klasyczne klasyfikatory nie~są~oryginalnym
wkładem pracy, lecz~stanowią punkty odniesienia (\emph{baseline}) zaczerpnięte
z~literatury i~z~publicznych bibliotek. Wkładem własnym jest~ich~adaptacja
do~zadania detekcji zainfekowanych węzłów w~grafie lineage, autorski detektor
LineageDetector oraz~adapter grafu DUT. Pochodzenie i~zakres adaptacji
poszczególnych składników zestawiono poniżej:

\begin{itemize}
  \item \textbf{Heurystyki strukturalne}: inspirowane rodziną miar
        predykcji krawędzi w~grafach (Common Neighbors, Jaccard, Adamic–Adar,
        Preferential Attachment)~\cite{LibenNowell2007,Martinez2016}. Ponieważ graf
        przepływu danych jest~dwudzielny, miary oparte na~wspólnych sąsiadach
        są~w~nim nieskuteczne. Zaadaptowano je~więc~do~postaci scoringu
        pojedynczego węzła opartego na~stopniu i~roli w~DAG.
  \item \textbf{Random Forest}: implementacja z~biblioteki
        \texttt{scikit-learn}~\cite{scikit-learn}. Adaptacja: wejściem jest~
        autorski wektor cech topologicznych węzła (tab.~\ref{tab:cechy})
        zamiast cech pary obiektów, a~niezbalansowanie klas obsłużono przez
        \texttt{class\_weight}.
  \item \textbf{RUSBoost}: implementacja z~biblioteki
        \texttt{imbalanced-learn}~\cite{imblearn}, wybrana dla~bezpośredniego
        nawiązania do~pracy Boińskiego et~al.~\cite{Boinski2025}. Adaptacja:
        klasyfikacja węzła zamiast pary, na~cechach topologicznych.
  \item \textbf{LightGBM}: biblioteka gradientowego boostingu
        firmy Microsoft~\cite{Ke2017}. Adaptacja: klasyfikacja
        węzła na~cechach topologicznych, z~kompensacją rozkładu klas przez
        \texttt{class\_weight}. Hiperparametrów nie~strojono, a~wszystkie
        klasyfikatory porównywane są~w~konfiguracjach domyślnych
        (por.~sekcja~\ref{sec:trafnosc}).
  \item \textbf{LineageDetector}, wkład własny: łączy powyższe cechy
        z~autorskim relacyjnym score'em kompletności peer-consistency
        (sekcja~\ref{sec:lineage-detector}) i~kalibracją izotoniczną.
        Nie~ma~bezpośredniego odpowiednika w~literaturze.
  \item \textbf{Adapter grafu DUT}, wkład własny: przekształca benchmark
        Dutkiewicza et~al.~\cite{Dutkiewicz2026} z~wierszowej reprezentacji
        lineage do~modelu tabela–zadanie–\emph{DATA\_FLOW}
        (sekcja~\ref{sec:dane-dut}). Sam benchmark wraz ze~scenariuszami
        pochodzi z~pracy źródłowej.
\end{itemize}

Rozgraniczenie komponentów zapożyczonych, z~jawnie wskazanym źródłem,
od~elementów autorskich pozwala jednoznacznie ocenić rzeczywisty wkład
pracy i~odróżnić go~od~narzędzi ogólnego przeznaczenia. Ma~ono również
znaczenie dla~interpretacji wyników rozdziału~4. Klasyfikatory bazowe
i~LineageDetector korzystają z~tego samego wektora cech topologicznych,
tych samych implementacji bibliotecznych i~tego samego potoku ewaluacji.
Detektor różni się~od~klasyfikatora Random Forest, na~którym jest~oparty,
przede wszystkim elementami autorskimi, czyli score'em kompletności
i~kalibracją. Wpływ każdego z~tych elementów na~wynik zmierzono bezpośrednio
w~badaniu ablacyjnym (sekcja~\ref{sec:ablacja}), co~pozwala oddzielić
wkład własny od~skuteczności samych metod bibliotecznych.

% ============================================================

\chapter{Problem i~sposób jego rozwiązania}

Rozdział precyzuje formalną definicję problemu detekcji węzłów
o~zerwanym lineage (sekcja~\ref{sec:sformulowanie}). Sekcja~\ref{sec:dane}
charakteryzuje wykorzystane dane, czyli zbiór DLG-DG-23 oraz~zewnętrzny
graf DUT wraz z~adapterem, sekcja~\ref{sec:symulacja} opisuje
metodykę symulacji broken lineage przez~usuwanie zadań,
a~sekcja~\ref{sec:architektura} architekturę potoku przetwarzania.
Zaimplementowane metody detekcji omówiono w~sekcji~\ref{sec:algorytmy},
realizację techniczną w~sekcji~\ref{sec:implementacja}, testy jednostkowe
i~odtwarzalność wyników w~sekcji~\ref{sec:testy}, a~dostępność implementacji
i~pochodzenie algorytmów w~sekcji~\ref{sec:repozytorium}.

% ============================================================
\section{Sformułowanie problemu}
\label{sec:sformulowanie}
% ============================================================

Graf lineage modeluje się jako skierowany graf DAG przepływu danych
$G = (V, E)$, gdzie $V = V_{J} \cup V_{T} \cup V_{F}$ obejmuje węzły
typu zadanie (\emph{Data Job}), tabela (\emph{Data Table}) i~pole,
a~$E$ zawiera krawędzie przepływu danych (\emph{DATA\_FLOW}) między zadaniami
a~tabelami. Zjawisko broken lineage odpowiada usunięciu z~grafu pewnego
zadania $J$. Tabele z~nim incydentne, czyli jego poprzedniki (tabele
karmiące $J$) oraz~następniki (tabele karmione przez~$J$), tracą część
informacji o~swoim pochodzeniu i~stają się~\emph{zainfekowane}
(ang.~\emph{infected}).

Obserwowany graf ma~postać $G_{\mathrm{obs}} = G \setminus \{J_1,\dots,J_m\}$,
powstałą przez~usunięcie podzbioru zadań. Zadaniem detekcji jest wyznaczenie
funkcji, która każdej tabeli $t \in V_{T}$ przypisuje wynik
$s(t) \in [0, 1]$ interpretowany jako~oszacowane prawdopodobieństwo
$p(\text{tabela } t \text{ jest zainfekowana})$. Wynikiem nie~jest~decyzja
binarna przy~ustalonym progu, lecz~ranking tabel. Tabele o~najwyższym
wyniku są~pierwszymi kandydatami do~ręcznej inspekcji.

Etykiety referencyjne (\emph{ground truth}) generowane są~automatycznie.
Tabela jest zainfekowana wtedy i~tylko wtedy, gdy była incydentna
do~co~najmniej jednego usuniętego zadania. Zbiór kandydatów ograniczono
do~węzłów typu \emph{Data Table}. Węzły typu pole (\emph{field}) łączy
z~tabelami wyłącznie relacja \emph{PARENT\_CHILD}, dlatego nie~podlegają one detekcji.

Zasadnicza trudność zadania polega na~tym, że~osierocona tabela-następnik
po~usunięciu zadania wygląda jak~prawdziwa tabela źródłowa
($\mathrm{in\_df} = 0$). Rozróżnienie tych~dwóch przypadków wymaga cech
opisujących pozycję węzła w~grafie DAG oraz~uczenia indukcyjnego
(cross-graph), a~nie~samego stopnia węzła.

% ============================================================
\section{Zbiory danych}
\label{sec:dane}
% ============================================================

W~pracy wykorzystano dwa źródła danych o~odmiennym pochodzeniu i~przeznaczeniu.
Wszystkie eksperymenty główne, obejmujące trening i~ewaluację metod
w~protokole indukcyjnym, przeprowadzono na~rzeczywistych grafach zbioru
DLG-DG-23. Uzupełnieniem jest~zewnętrzny graf DUT, zbudowany na~podstawie
benchmarku Dutkiewicza et~al.~\cite{Dutkiewicz2026}, który służy wyłącznie
do~sprawdzenia przenośności modeli wytrenowanych na~DLG-DG-23. Oba źródła
opisano w~kolejnych podsekcjach.

\subsection{Zbiór DLG-DG-23}
\label{sec:dane-dlg}

Zbiór DLG-DG-23 pochodzi z~pracy Chen et~al.~\cite{Chen2024}
i~zawiera 18~grafów lineage pobranych ze~środowiska produkcyjnego
Huawei Cloud. Dane zostały zanonimizowane, a~nazwy tabel i~atrybutów
zastąpiono identyfikatorami w~formacie UUID. Schematów bazy danych
nie~udostępniono. Zbiór opublikowano w~2024~roku jako pierwszy
otwarty zestaw rzeczywistych grafów lineage.

Każdy graf zawiera węzły trzech typów: \emph{Data Job} (zapytanie SQL
lub~zadanie ETL), \emph{Data Table}
(tabela lub~widok przechowujący dane) oraz~\emph{field} (atrybut lub~kolumna
należąca do~tabeli). Łączą je~krawędzie dwóch typów: \emph{DATA\_FLOW}
(skierowana krawędź przepływu danych między tabelą a~zadaniem)
oraz~\emph{PARENT\_CHILD} (relacja przynależności pola do~tabeli). Pełny
graf jest~zatem \emph{heterogeniczny}. Przedmiotem detekcji są~wyłącznie
węzły typu \emph{Data Table}, a~sygnał czerpany jest z~podgrafu przepływu
danych (krawędzie \emph{DATA\_FLOW} między węzłami \emph{Data Table}
i~\emph{Data Job}). Podgraf ten~ma~strukturę dwudzielną i~skierowaną,
interpretowaną jako graf DAG. Grafy różnią się skalą, a~liczba węzłów waha się
od~278 do~17~085.

Zbiór DLG-DG-23 wybrano, ponieważ jest pierwszym publicznie udostępnionym
zestawem rzeczywistych grafów lineage~\cite{Chen2024}, a~w~trakcie prowadzenia
badań nie~zidentyfikowano innego zbioru tego rodzaju. Anonimizacja nazw
węzłów ma~bezpośredni wpływ na~dobór metod, ponieważ metryki podobieństwa
oparte na~nazwach obiektów (np.~miara Jaro–Winklera stosowana przez Boińskiego
et~al.~\cite{Boinski2025}) są~nieaplikowalne. W~niniejszej pracy sygnał
ograniczono zatem do~atrybutów topologicznych grafu. Ograniczenie
to~jest zarazem częścią pytania badawczego o~to, czy~sama struktura grafu
lineage wystarcza do~wskazania tabel o~zerwanym pochodzeniu danych.
Metoda skuteczna przy~tak ograniczonym wglądzie w~graf ma~przy~tym większą
szansę na~zastosowanie praktyczne niż~podejścia wymagające pełnych
metadanych.

Tabela~\ref{tab:dlg-stats} przedstawia statystyki poszczególnych
grafów zbioru DLG-DG-23.

\begin{table}[ht]
\centering
\caption{Charakterystyka grafów zbioru DLG-DG-23.}
\label{tab:dlg-stats}
\begin{threeparttable}
\begin{tabular}{llrl}
\toprule
\textbf{ID} & \textbf{Węzły} & \textbf{Krawędzie DATA\_FLOW} & \textbf{Skala} \\
\midrule
DLG1  &    298 &    24 & mały   \\
DLG2  &    464 &    37 & mały   \\
DLG3  &    603 &    57 & mały   \\
DLG4  &    415 &    32 & mały   \\
DLG5  & 1\,299 &   175 & średni \\
DLG6  & 13\,840 & 1\,443 & duży  \\
DLG7  &    574 &    43 & mały   \\
DLG8  &    546 &    34 & mały   \\
DLG9  &    756 &    62 & mały   \\
DLG10 & 5\,378 &   497 & duży   \\
DLG11 & 2\,024 &   121 & średni \\
DLG12 & 5\,645 &   469 & duży   \\
DLG13 & 2\,157 &   229 & średni \\
DLG14 &    453 &    24 & mały   \\
DLG15 &    558 &    41 & mały   \\
DLG16 & 5\,574 &   745 & duży   \\
DLG17 &    651 &    38 & mały   \\
DLG18 & 17\,085 & 1\,694 & duży \\
\bottomrule
\end{tabular}
\begin{tablenotes}
\small
\item Skala: mały $<$ 1\,000 węzłów, średni 1\,000–5\,000, duży $>$ 5\,000.
\end{tablenotes}
\end{threeparttable}
\end{table}

\subsection{Zewnętrzny graf DUT}
\label{sec:dane-dut}

Drugim źródłem danych jest~benchmark opracowany przez Dutkiewicza
et~al.~\cite{Dutkiewicz2026} w~zespole badawczym promotora
(sekcja~\ref{sec:literatura}). Obejmuje on~syntetyczne scenariusze broken
lineage, wygenerowane dla~przykładowej bazy \emph{Northwind}~\cite{Northwind}.
Lineage zapisany jest~w~nim na~poziomie pojedynczych wierszy, w~postaci grafu
wiedzy RDF oraz~plików CSV, w~których każdy rekord wiąże wartość źródłową
z~wartością docelową, która z~niej powstała. Jest to~zatem reprezentacja
drobnoziarnista (sekcja~\ref{sec:lineage_def}), niezgodna z~gruboziarnistym
modelem tabela–zadanie, na~którym operują metody opisane w~niniejszej pracy.

Aby~ocenić na~tych danych modele wytrenowane na~DLG-DG-23 bez~zmian w~ich
implementacji, opracowano adapter (moduł \texttt{dutkiewicz\_adapter.py}),
który agreguje zależności wierszowe do~poziomu tabel. Dla~każdej tabeli
docelowej tworzony jest~węzeł typu \emph{Data Job}, do~którego prowadzą
krawędzie \emph{DATA\_FLOW} od~wszystkich tabel źródłowych, z~których
pochodzą jej~wiersze. Z~węzła tego prowadzi następnie krawędź \emph{DATA\_FLOW}
do~samej tabeli docelowej. Kolumny kluczy występujące w~zapisie lineage
dodawane są~jako~węzły typu \emph{field}, połączone z~tabelą relacją
\emph{PARENT\_CHILD}. Wynikowy graf ma~te~same typy węzłów i~krawędzi
co~grafy DLG-DG-23, dzięki czemu przechodzi przez potok przetwarzania
identycznie jak~one.

Graf DUT liczy 21~tabel, 14~zadań i~33~krawędzie przepływu danych. Jest więc
o~rzędy wielkości mniejszy od~grafów DLG-DG-23 i~ma~odmienną, łańcuchową
topologię, w~której tabele rzadko współdzielą zadania. Adapter nie~odtwarza
przy~tym pełnych schematów tabel, ponieważ benchmark ich~nie~zawiera,
dlatego liczba pól jest~mniejsza niż~w~DLG-DG-23. Poprawność odwzorowania
weryfikuje sześć testów jednostkowych (tabela~\ref{tab:testy}). Graf DUT
nie~jest~wykorzystywany w~treningu żadnej z~metod, a~wyniki uzyskane na~nim
przedstawiono w~sekcji~\ref{sec:cross_dataset}.

% ============================================================
\section{Metodyka symulacji broken lineage przez usuwanie zadań}
\label{sec:symulacja}
% ============================================================

Ponieważ zbiór DLG-DG-23 zawiera w~pełni zaobserwowane grafy lineage,
zjawisko broken lineage symulowane jest przez~kontrolowane usuwanie
zadań (nie~pojedynczych krawędzi). Podejście to~odwzorowuje
przyczynę broken lineage, ponieważ zniknięcie tabeli tymczasowej lub~ukrycie
logiki przez UDF odpowiada utracie całego węzła-zadania wraz z~jego krawędziami.

\subsection{Procedura symulacji i~konstrukcja etykiet}
\label{sec:procedura-symulacji}

Kandydatami do~usunięcia są~\emph{zadania łączące}, czyli zadania mające
co~najmniej jedną wejściową i~jedną wyjściową krawędź \emph{DATA\_FLOW}
($\mathrm{in\_df} > 0 \land \mathrm{out\_df} > 0$). Usunięcie takiego zadania
zrywa lineage po~obu stronach, ponieważ tabele karmiące je~tracą konsumenta,
a~tabele przez~nie~karmione tracą producenta.

Instancję klasyfikacji węzłów buduje się~z~grafu $G$ w~kilku krokach.
Najpierw losowany jest~podzbiór zadań łączących przeznaczonych do~usunięcia,
którego udział reguluje parametr \texttt{removal\_ratio}. W~eksperymentach
przyjęto wartość 10\%, a~wpływ tego parametru zbadano w~analizie wrażliwości
(sekcja~\ref{sec:przeciek}). Zbiór tabel zainfekowanych wyznacza się~jako~sumę
tabel incydentnych do~usuniętych zadań. Graf obserwowany $G_{\mathrm{obs}}$
powstaje następnie przez~usunięcie wybranych zadań wraz z~ich krawędziami.
Kandydatami do~detekcji są~wszystkie tabele pozostałe w~$G_{\mathrm{obs}}$,
przy~czym etykieta~1 oznacza tabelę zainfekowaną, a~etykieta~0 tabelę czystą.

Dodatkowo zaimplementowano wariant kontrolny, w~którym każda instancja
usuwa dokładnie jedno zadanie, co~daje czystą interpretację
przyczynową (jedno zdarzenie broken lineage $=$ jedno brakujące zadanie).

\subsection{Kontrola przecieku – wykluczenie węzłów izolowanych}
\label{sec:exclude-isolated}

Usunięcie zadania może~pozbawić tabelę wszystkich krawędzi
przepływu danych ($\mathrm{in\_df} = 0 \land \mathrm{out\_df} = 0$).
Ponieważ w~oryginalnych grafach tabele izolowane nie~występują, trywialna
reguła „tabela izolowana $\Rightarrow$ zainfekowana'' pozwala rozpoznać
znaczną część pozytywów bez~żadnego uczenia i~sztucznie zawyża wynik.
Aby~zapewnić uczciwą ewaluację, wprowadzono flagę \texttt{exclude\_isolated},
która usuwa z~puli kandydatów tabele w~pełni izolowane. Wariant ten~stawia
trudniejsze i~bardziej realistyczne zadanie, polegające na~wykryciu tabel
o~częściowo zerwanym lineage, które nadal mają część połączeń. W~pracy raportowane
są~oba warianty, a~wariant uczciwy (z~wykluczeniem izolowanych) traktowany
jest jako~główny.

Tabela~\ref{tab:symulacja} przedstawia skalę symulacji dla~czterech wartości
parametru \texttt{removal\_ratio}. Łącznie w~18~grafach zbioru DLG-DG-23
znajduje się 1453~zadania łączące (od~3 do~404 na~graf), spośród których usuwana
jest~losowa część. Liczby podano dla~wariantu przed wykluczeniem tabel
izolowanych, ponieważ tylko wtedy można zmierzyć ich~udział wśród pozytywów.

\begin{table}[ht]
\centering
\caption[Charakterystyka symulacji broken lineage]{Charakterystyka symulacji broken lineage w~zależności
od~parametru \texttt{removal\_ratio}. Wartości uśrednione po~18~grafach
i~trzech ziarnach losowości (42--44). W~ostatniej kolumnie podano odsetek
tabel zainfekowanych, które utraciły wszystkie krawędzie przepływu
danych i~są~rozpoznawalne trywialną regułą izolacji.}
\label{tab:symulacja}
\begin{tabular}{ccccc}
\toprule
\texttt{removal\_ratio} & \textbf{Usuniętych} & \textbf{Tabel} & \textbf{Udział} & \textbf{Pozytywów} \\
 & \textbf{zadań} & \textbf{zainfekowanych} & \textbf{pozytywów} & \textbf{izolowanych} \\
\midrule
0{,}05 & 4{,}2 & 15{,}1 & 11{,}4\% & 19{,}7\% \\
0{,}10 & 8{,}0 & 25{,}1 & 19{,}0\% & 22{,}9\% \\
0{,}20 & 16{,}2 & 45{,}2 & 34{,}1\% & 29{,}8\% \\
0{,}30 & 24{,}3 & 62{,}6 & 47{,}3\% & 34{,}5\% \\
\bottomrule
\end{tabular}
\end{table}

Dwie obserwacje z~tabeli mają~znaczenie dla~konstrukcji eksperymentu.
Przyjęta w~biegu kanonicznym wartość $0{,}10$ daje średnio 25~tabel
zainfekowanych na~graf przy~19\% pozytywów. Zadanie pozostaje więc~wyraźnie
niezrównoważone, ale~klasa pozytywna jest~dostatecznie liczna, by~metryki
rankingowe były wiarygodne. Udział pozytywów trywialnych (tabel w~pełni
izolowanych) rośnie z~kolei z~19{,}7\% do~34{,}5\% wraz ze~wzrostem
odsetka usuwanych zadań. Im~silniejsze zaburzenie, tym~większa część zadania
daje się~rozwiązać bez~żadnego uczenia. Liczby te~określają skalę przecieku,
któremu zapobiega opisana wyżej flaga \texttt{exclude\_isolated}. Przy~biegu
kanonicznym usuwa ona~z~puli kandydatów blisko jedną czwartą pozytywów,
i~to~tych najłatwiejszych.

% ============================================================
\section{Architektura rozwiązania}
\label{sec:architektura}
% ============================================================

Potok przetwarzania składa się z~czterech komponentów: modułu wczytywania
grafów (wraz z~adapterem grafu DUT), modułu symulacji broken lineage (budowa
instancji klasyfikacji węzłów), metod detekcji oraz~modułu ewaluacji indukcyjnej.

Rysunek~\ref{rys:potok} przedstawia schemat potoku przetwarzania, od~wczytania
grafów, przez symulację broken lineage i~ekstrakcję cech, po~scoring detektora
i~wyznaczenie metryk rankingowych. Oba źródła danych trafiają do~wspólnego
potoku. Grafy DLG-DG-23 wczytywane są~bezpośrednio z~plików JSON, natomiast
graf DUT przechodzi najpierw przez adapter, który sprowadza go~do~tego samego
modelu danych. Od~etapu symulacji oba rodzaje grafów przetwarzane są~identycznie,
różni je~jedynie rola w~ewaluacji.

\begin{figure}[ht]
  \centering
  \includegraphics[width=\linewidth]{rys_potok.pdf}
  \caption[Schemat potoku przetwarzania]{Schemat potoku przetwarzania. Grafy
           DLG-DG-23 oceniane są~w~protokole indukcyjnym leave-one-graph-out
           (trening na~17~grafach, test na~grafie niewidzianym). Graf DUT,
           sprowadzony adapterem do~modelu tabela–zadanie, służy wyłącznie
           do~testu modeli wytrenowanych na~DLG-DG-23.}
  \label{rys:potok}
\end{figure}

\subsection{Moduł wczytywania danych}

Dane wejściowe każdego grafu przechowywane są~w~dwóch plikach JSON:
\texttt{DLGi-node.json} (opisy węzłów z~atrybutem \texttt{asset\_type})
oraz~\texttt{DLGi-edge.json} (krawędzie z~atrybutami \texttt{relation\_type}
i~\texttt{relation\_id}). Moduł \texttt{loader.py} buduje z~nich obiekt
\texttt{nx.DiGraph} (NetworkX) i~zwraca słownik wszystkich 18~grafów.
Graf DUT wczytywany jest~przez moduł \texttt{dutkiewicz\_adapter.py}, który
zwraca obiekt tego samego typu, z~atrybutami węzłów i~krawędzi zgodnymi
z~konwencją DLG-DG-23 (sekcja~\ref{sec:dane-dut}). Dalsze moduły potoku
nie~rozróżniają więc~źródła grafu.

\subsection{Moduł symulacji usuwania zadań}

Moduł \texttt{job\_removal.py} realizuje metodykę z~sekcji~\ref{sec:symulacja}:
identyfikuje zadania łączące, usuwa ich podzbiór, wyznacza zbiór zainfekowanych
tabel i~zwraca instancję (graf $G_{\mathrm{obs}}$, listę kandydatów,
etykiety, listę usuniętych zadań). Obsługuje warianty \texttt{side},
określający stronę incydencji, po~której tabele uznaje się~za~zainfekowane
(strona producenta, konsumenta lub~obie), oraz~\texttt{exclude\_isolated}. Funkcja \texttt{inductive\_split} dzieli
identyfikatory grafów na~rozłączne zbiory treningowy i~testowy (cross-graph).

\subsection{Moduł ewaluacji i~metryki rankingowe}

Moduł \texttt{node\_metrics.py} wyznacza metryki rankingowe: AUC-ROC,
AUC-PR, Precision@k, Recall@k, Hits@k oraz~Brier score. Wyniki agregowane
są~po~grafach testowych w~protokole indukcyjnym. Dla~stabilności oszacowań
stosuje się powtórzenia z~różnymi ziarnami i~trymowanie wartości skrajnych
(min/max).

Ewaluację prowadzi się~w~protokole indukcyjnym cross-graph.
Model trenuje się~na~podzbiorze grafów, a~testuje na~grafach niewidzianych
(leave-one-graph-out lub~podział zbioru 18~grafów). Zbiory treningowy
i~testowy są~rozłączne, co~sprawdza zdolność metody do~generalizacji
na~nowe grafy, a~nie~jedynie do~zapamiętania jednego. Konsekwencją tego
protokołu jest~wykluczenie metody Node2Vec~\cite{Grover2016}, która uczy zanurzeń
węzłów transduktywnie, dla~konkretnego grafu, przez~co~embeddingi
są~geometrycznie nieporównywalne między grafami~\cite{Hamilton2017}.
Obserwacja ta~motywuje zastosowanie indukcyjnych
grafowych sieci neuronowych (GraphSAGE~\cite{Hamilton2017},
R-GCN~\cite{Schlichtkrull2018}, HGT~\cite{Hu2020}), wskazanych w~rozdziale~5 jako~kierunek dalszych badań.

% ============================================================
\section{Zaimplementowane metody detekcji}
\label{sec:algorytmy}
% ============================================================

Zaimplementowano trzy grupy metod detekcji zainfekowanych tabel, różniące się
ilością wykorzystywanej informacji i~stopniem złożoności. Heurystyki strukturalne
(sekcja~\ref{sec:heurystyki}) przypisują tabeli wynik na~podstawie prostych
reguł i~nie~wymagają uczenia, dlatego wyznaczają dolny punkt odniesienia.
Klasyczne klasyfikatory uczenia maszynowego (sekcja~\ref{sec:ml}) uczą się
na~grafach treningowych zależności między wektorem cech topologicznych
a~etykietą, co~pozwala ocenić, ile~wnosi samo uczenie indukcyjne. Autorski
detektor LineageDetector (sekcja~\ref{sec:lineage-detector}) rozszerza ten~wektor
o~relacyjny score kompletności i~uzupełnia model o~kalibrację prawdopodobieństw.

Wszystkie metody mają wspólny interfejs. Dla~każdej tabeli-kandydata zwracają
wynik z~przedziału $[0, 1]$, na~podstawie którego budowany jest~ranking. Dzięki
temu każda z~nich jest~oceniana w~tym samym potoku, na~tych samych instancjach
i~tymi samymi metrykami, a~różnice wyników można przypisać samym metodom.

% ============================================================
\subsection{Heurystyki strukturalne}
\label{sec:heurystyki}
% ============================================================

Heurystyki strukturalne przypisują tabeli wynik podejrzliwości wyłącznie
na~podstawie jej pozycji w~grafie, bez~etapu uczenia. Są~naturalnym
punktem odniesienia (\emph{baseline}) dla~metod uczonych. Miary oparte
na~wspólnych sąsiadach (Common Neighbors, Jaccard,
Adamic-Adar)~\cite{LibenNowell2007} są~w~grafie
dwudzielnym nieskuteczne, ponieważ węzły tego~samego typu nigdy nie~mają
wspólnych sąsiadów bezpośrednich. Z~tego powodu wykluczono je~z~implementacji.

W~pracy zaimplementowano cztery heurystyki bazowe. Piątym punktem
odniesienia w~tej~grupie jest~sam \emph{score kompletności}
(sekcja~\ref{sec:lineage-detector}), oceniany dodatkowo w~oderwaniu
od~detektora, co~pozwala rozdzielić wkład sygnału relacyjnego od~wkładu
uczenia. W~zestawieniach rozdziału~4 występuje on~zatem jako~osobna metoda
obok czterech poniższych:

\begin{itemize}
  \item \textbf{Anomalia stopnia} (\texttt{degree\_anomaly}): wynik
        $1/(1+\deg(t))$, gdzie $\deg(t)$ jest stopniem całkowitym tabeli.
        Heurystyka opiera się~na~założeniu, że~tabela, która utraciła
        zadanie, ma~mniej połączeń niż~przed zerwaniem lineage.
  \item \textbf{Brzeg grafu} (\texttt{boundary}): suma wskaźników
        $\mathbf{1}[\mathrm{in\_df}=0] + \mathbf{1}[\mathrm{out\_df}=0]$.
        Tabele na~brzegu przepływu (korzenie i~liście) są~podejrzane,
        ponieważ zerwanie lineage często przesuwa tabelę na~brzeg grafu.
  \item \textbf{Niska łączność z~zadaniami}
        (\texttt{low\_job\_connectivity}): wynik $1/(1+n_{J}(t))$,
        gdzie $n_{J}(t)$ to~liczba unikatowych zadań sąsiadujących
        z~tabelą przez krawędzie \emph{DATA\_FLOW}.
  \item \textbf{Reguła korzenia} (\texttt{rule\_root}): binarny wskaźnik
        $\mathbf{1}[\mathrm{in\_df}=0]$, zgodnie z~którym każda tabela
        bez~producenta jest uznawana za~podejrzaną. To~celowo trywialny punkt odniesienia,
        obrazujący siłę najprostszego możliwego symptomu.
\end{itemize}

Heurystyki te~mierzą wyłącznie lokalny symptom (stopień, rola).
Ich~ograniczenie, przede wszystkim nieodróżnianie osieroconej tabeli
od~prawdziwego źródła, motywuje metody uczone oraz~autorski detektor
LineageDetector (sekcja~\ref{sec:lineage-detector}), który wzbogaca opis
węzła o~relacyjny score kompletności (peer-consistency).

% ============================================================
\subsection{Klasyczne metody uczenia maszynowego}
\label{sec:ml}
% ============================================================

Klasyczne metody ML traktują detekcję jako~binarną klasyfikację węzła, w~której
wejściem jest wektor cech opisujących pozycję tabeli w~grafie DAG,
a~wyjściem prawdopodobieństwo zainfekowania. Ze~względu na~anonimizację
(UUID) jako~cechy przyjęto wyłącznie atrybuty topologiczne, ponieważ metryki
oparte na~nazwach są~nieaplikowalne.

\subsubsection{Wektor cech węzła}

Tabela~\ref{tab:cechy} przedstawia pełny wektor 13~cech topologicznych
wyznaczanych dla~każdej tabeli-kandydata. Cechy opisują pozycję węzła
w~grafie przepływu danych, bo~to~ona, a~nie~sam stopień, pozwala odróżnić
prawdziwą tabelę źródłową od~tabeli osieroconej po~usunięciu zadania.

\begin{table}[ht]
\centering
\caption[Wektor cech topologicznych węzła]{Wektor cech topologicznych węzła (moduł \texttt{node\_features}).}
\label{tab:cechy}
\small
\begin{tabular}{cl>{\raggedright\arraybackslash}p{5.2cm}>{\raggedright\arraybackslash}p{4.6cm}}
\toprule
\textbf{Nr} & \textbf{Cecha} & \textbf{Opis} & \textbf{Rola w~detekcji} \\
\midrule
1 & \texttt{in\_df} & Liczba wejściowych krawędzi \emph{DATA\_FLOW}
    (od~zadań) & Brak producenta ($=0$) to~główny symptom zerwania \\
2 & \texttt{out\_df} & Liczba wyjściowych krawędzi \emph{DATA\_FLOW}
    (do~zadań) & Symetryczny symptom, czyli utrata konsumenta \\
3 & \texttt{degree} & Stopień całkowity węzła (wszystkie typy krawędzi)
    & Ogólna łączność, spadająca po~usunięciu zadania \\
4 & \texttt{role} & Rola w~przepływie (źródło, pośrednia, ujście,
    poza przepływem) & Pozycja funkcjonalna tabeli w~DAG \\
5 & \texttt{flow\_depth} & Głębokość w~DAG (odległość od~źródeł)
    & Osierocona tabela ma~sztucznie małą głębokość \\
6 & \texttt{flow\_reach} & Zasięg w~dół DAG (odległość do~ujść)
    & Symetrycznie sztucznie mały zasięg \\
7 & \texttt{depth\_plus\_reach} & Suma głębokości i~zasięgu
    & Długość ścieżki przez węzeł, wykrywa „skrócone'' łańcuchy \\
8 & \texttt{is\_root} & Wskaźnik $\mathbf{1}[\mathrm{in\_df}=0]$
    & Binarna wersja symptomu utraty producenta \\
9 & \texttt{is\_leaf} & Wskaźnik $\mathbf{1}[\mathrm{out\_df}=0]$
    & Binarna wersja symptomu utraty konsumenta \\
10 & \texttt{n\_field\_children} & Liczba pól tabeli (krawędzie
    \emph{PARENT\_CHILD}) & Rozmiar tabeli, sygnał niezależny
    od~przepływu \\
11 & \texttt{n\_job\_neighbors} & Liczba unikatowych zadań sąsiadujących
    & Łączność z~warstwą zadań \\
12 & \texttt{avg\_job\_in\_df} & Średni stopień wejściowy sąsiednich zadań
    & Charakterystyka otoczenia węzła \\
13 & \texttt{avg\_job\_out\_df} & Średni stopień wyjściowy sąsiednich
    zadań & Jak~wyżej, od~strony wyjść \\
\bottomrule
\end{tabular}
\end{table}

Detektor LineageDetector (sekcja~\ref{sec:lineage-detector})
rozszerza ten wektor o~trzy cechy kompletności
(\texttt{peer\_producer\_score}, \texttt{peer\_consumer\_score},
\texttt{completeness\_anomaly}), opisane w~sekcji~\ref{sec:lineage-detector}.

\subsubsection{Random Forest}

Random Forest (RF)~\cite{Breiman2001} to~metoda uczenia zespołowego oparta na~uśrednianiu
predykcji drzew decyzyjnych trenowanych na~losowych podpróbkach danych
i~cech (\emph{bagging}). Dla~obsługi niezrównoważonych klas zastosowano
\texttt{class\_weight=\textquotedbl{}balanced\textquotedbl{}}.

\subsubsection{RUSBoost}

RUSBoost (\emph{Random Under-Sampling Boosting})~\cite{Seiffert2010} łączy boosting AdaBoost
z~losowym podpróbkowaniem klasy większościowej przed~każdą iteracją.
Zaimplementowano go~dla~bezpośredniego nawiązania do~wyników Boińskiego
et~al.~\cite{Boinski2025}, przy~czym w~niniejszej pracy klasyfikuje on~węzły,
a~nie~pary obiektów.

\subsubsection{LightGBM}

LightGBM~\cite{Ke2017} to~implementacja gradientowego boostingu drzew decyzyjnych,
budująca kolejne drzewa sekwencyjnie, każde minimalizując błąd resztkowy
poprzednika. Niezbalansowanie klas kompensuje parametr
\texttt{class\_weight}~$=$~\texttt{balanced}, przypisujący obserwacjom wagi
odwrotnie proporcjonalne do~liczności klasy.

% ============================================================
\subsection{Autorski detektor LineageDetector}
\label{sec:lineage-detector}
% ============================================================

Autorski detektor \texttt{LineageDetector} (klasa
\texttt{LineageHealthDetector}) łączy w~jednym skalibrowanym modelu trzy
elementy. Podstawę stanowią cechy pozycji w~DAG, wyznaczane przez moduł
\texttt{node\_features} i~opisujące węzeł lokalnie. Uzupełnia je~score
kompletności \emph{peer-consistency}, opisany w~dalszej części sekcji,
który wnosi sygnał relacyjny i~stanowi główny element wkładu własnego pracy.
Całość zamyka kalibracja prawdopodobieństw metodą izotoniczną
(\texttt{CalibratedClassifierCV}), dzięki której wynik detektora można
interpretować jako~oszacowanie $p(\text{zainfekowana})$, a~nie~jedynie
jako~wartość porządkującą ranking.

\subsubsection{Score kompletności peer-consistency}

Zaproponowany w~pracy score kompletności opiera się na~założeniu,
że~tabela o~zerwanym lineage jest~\emph{strukturalnie niespójna}
ze~swoimi sąsiadami (peerami) w~grafie DAG. Dla~tabeli-korzenia
($\mathrm{in\_df} = 0$) rozważa się~zbiór peerów-wejść, czyli inne tabele
wpadające do~tych~samych zadań downstream. Jeśli~peerzy mają
producenta, a~rozpatrywana tabela nie, to~jej brak producenta jest~podejrzany,
ponieważ prawdziwe źródło danych sąsiaduje zwykle z~innymi źródłami. Symetrycznie
dla~tabeli-liścia ($\mathrm{out\_df} = 0$) i~jej peerów po~stronie konsumpcji.
Formalnie dla~tabeli $t$ definiuje się:
%
\[
  \mathrm{prod}(t) = \frac{1}{|P(t)|}\sum_{p \in P(t)} \mathbf{1}[\mathrm{in\_df}(p) > 0],
  \qquad
  \mathrm{cons}(t) = \frac{1}{|C(t)|}\sum_{p \in C(t)} \mathbf{1}[\mathrm{out\_df}(p) > 0],
\]
%
gdzie $P(t)$ i~$C(t)$ oznaczają zbiory peerów odpowiednio po~stronie
producenta i~konsumenta, a~$\mathbf{1}[\cdot]$ funkcję wskaźnikową.
Bramkowany sygnał anomalii aktywuje się~wyłącznie na~brzegu grafu:
%
\[
  \mathrm{anomaly}(t) =
  \max\bigl(\mathbf{1}[\mathrm{in\_df}(t)=0]\cdot\mathrm{prod}(t),\;
            \mathbf{1}[\mathrm{out\_df}(t)=0]\cdot\mathrm{cons}(t)\bigr).
\]
%
Sygnał ten~ma~charakter relacyjny. Różni się~jakościowo od~lokalnego
stopnia węzła i~dostarcza informacji niedostępnej klasyfikatorowi
korzystającemu wyłącznie z~cech lokalnych. Definicja jest~symetryczna
względem obu~stron przepływu, jednak na~zbiorze DLG-DG-23 aktywna
okazuje się~wyłącznie strona producenta, ponieważ każde zadanie w~tym~zbiorze
zapisuje dokładnie jedną tabelę, przez~co~zbiór $C(t)$ jest~zawsze pusty.
Konsekwencje tej~własności omówiono w~sekcji~\ref{sec:waznosc}. Trzy wynikające z~niego cechy
(\texttt{peer\_producer\_score}, \texttt{peer\_consumer\_score},
\texttt{completeness\_anomaly}) rozszerzają wektor cech z~tabeli~\ref{tab:cechy}
i~są~dostępne wyłącznie detektorowi LineageDetector.

Bazowym estymatorem jest~domyślnie Random Forest, a~wariantem
alternatywnym LightGBM. Interfejs \texttt{fit}/\texttt{score} jest~zgodny z~pozostałymi
klasyfikatorami, co~umożliwia bezpośrednie porównanie w~tym~samym potoku.
Konstrukcja detektora umożliwia badania ablacyjne. Wyłączenie kalibracji
(\texttt{calibrate=False}) izoluje jej~wpływ na~Brier score, a~wyłączenie
cech kompletności (\texttt{use\_completeness=False}) wpływ sygnału
relacyjnego na~AUC.

% ============================================================
\section{Implementacja}
\label{sec:implementacja}
% ============================================================

Sekcja opisuje realizację techniczną potoku z~sekcji~\ref{sec:architektura}:
zastosowane narzędzia, organizację kodu oraz~te~fragmenty implementacji,
które są~istotne dla~zrozumienia metody, a~nie~wynikają wprost z~jej~opisu
formalnego.

\subsection{Stos technologiczny}

Całość zaimplementowano w~języku Python. Tabela~\ref{tab:stos} zestawia
wykorzystane biblioteki wraz z~wersjami, w~których wykonano bieg kanoniczny.
Podanie wersji jest~istotne dla~odtwarzalności, ponieważ implementacje algorytmów
uczenia maszynowego bywają modyfikowane między wydaniami, co~może~wpływać
na~wyniki liczbowe.

\begin{table}[ht]
\centering
\caption{Stos technologiczny i~wersje bibliotek użyte w~biegu kanonicznym.}
\label{tab:stos}
\begin{tabular}{lll}
\toprule
\textbf{Składnik} & \textbf{Wersja} & \textbf{Rola w~pracy} \\
\midrule
Python~\cite{Python}                 & 3.10.2  & Środowisko uruchomieniowe \\
NetworkX~\cite{NetworkX}             & 3.4.2   & Reprezentacja i~przetwarzanie grafów \\
NumPy~\cite{NumPy}                   & 1.26.4  & Operacje macierzowe, wektory cech \\
pandas~\cite{pandas}                 & 1.5.1   & Agregacja wyników, eksport CSV \\
scikit-learn~\cite{scikit-learn}     & 1.6.1   & Random Forest, kalibracja izotoniczna, metryki \\
LightGBM~\cite{Ke2017}               & 4.6.0   & Gradientowy boosting drzew \\
imbalanced-learn~\cite{imblearn}     & 0.14.1  & RUSBoost \\
SciPy~\cite{SciPy}                   & 1.15.3  & Test kolejności par Wilcoxona \\
Matplotlib~\cite{Matplotlib}         & 3.8.0   & Rysunki \\
pytest~\cite{pytest}                 & 9.0.2   & Testy jednostkowe \\
\bottomrule
\end{tabular}
\end{table}

\subsection{Organizacja repozytorium}

Kod podzielono na~pakiety odpowiadające etapom potoku, a~skrypty
uruchomieniowe oddzielono od~modułów bibliotecznych, aby~każdy eksperyment
z~rozdziału~4 dawał się~odtworzyć jednym poleceniem. Listing~\ref{lst:repo}
przedstawia strukturę katalogów.

\begin{lstlisting}[caption={Struktura repozytorium (fragmenty istotne
dla~detekcji węzłów).},label={lst:repo},numbers=none,language={},float=htbp]
src/
  data/
    loader.py               # wczytanie 18 grafow DLG-DG-23 z JSON
    dutkiewicz_adapter.py   # adapter benchmarku zewnetrznego (graf DUT)
  detection/
    job_removal.py          # symulacja broken lineage, etykiety, podzial
    node_features.py        # 13 cech topologicznych wezla
    completeness.py         # score peer-consistency (wklad wlasny)
    lineage_detector.py     # LineageHealthDetector
    node_classifier.py      # heurystyki + klasyczne klasyfikatory
    node_metrics.py         # metryki rankingowe
    run_detection_experiments.py   # bieg kanoniczny (rozdz. 4.2)
    run_ablation.py                # badanie ablacyjne (rozdz. 4.5)
    run_cross_dataset.py           # test cross-dataset (rozdz. 4.9)
    run_removal_sensitivity.py     # analiza wrazliwosci (rozdz. 4.6)
scripts/
  test_istotnosci.py        # testy Wilcoxona (rozdz. 4.4)
  tabela_symulacji.py       # tabela charakterystyki symulacji
tests/                      # 83 testy jednostkowe (pytest)
results/                    # wyniki biegow w formacie CSV
\end{lstlisting}

\subsection{Format danych i~wczytywanie}

Każdy graf zbioru DLG-DG-23 opisany jest~parą plików JSON: lista węzłów
z~atrybutem \texttt{asset\_type} oraz~lista krawędzi z~atrybutami
\texttt{relation\_id} i~\texttt{relation\_type} (listing~\ref{lst:json}).
Moduł \texttt{loader.py} przekształca je~w~obiekt \texttt{nx.DiGraph}
i~przenosi typ węzła oraz~typ relacji do~atrybutów grafu. Dzięki temu
pozostałe moduły filtrują podgraf przepływu danych bez~ponownego sięgania
do~plików.

\begin{lstlisting}[caption={Fragment plików wejściowych
\texttt{DLGi-node.json} i~\texttt{DLGi-edge.json} (identyfikatory skrócone).},
label={lst:json},numbers=none,language={},float=htbp]
// DLG1-node.json
{"nodes": [
  {"asset_id": "3f2a...c81", "asset_type": "Data Table"},
  {"asset_id": "9b07...4de", "asset_type": "Data Job"},
  {"asset_id": "a154...77f", "asset_type": "Data Field"}
]}

// DLG1-edge.json
{"edges": [
  {"relation_id": "r-0012", "relation_type": "DATA_FLOW",
   "source": "3f2a...c81", "target": "9b07...4de"},
  {"relation_id": "r-0013", "relation_type": "PARENT_CHILD",
   "source": "3f2a...c81", "target": "a154...77f"}
]}
\end{lstlisting}

% Listingi 3.1 i 3.2 maja zostac zlozone przed ta podsekcja, aby nie
% rozcinaly jej pierwszego akapitu.
\FloatBarrier
\subsection{Wyznaczanie score'u peer-consistency}

Najistotniejszym fragmentem implementacji jest~wyznaczanie relacyjnego
score'u kompletności, zdefiniowanego wzorami w~sekcji~\ref{sec:lineage-detector}.
Listing~\ref{lst:peer} pokazuje jego~rdzeń. Sam wzór nie~precyzuje, że~zbiory
peerów budowane są~przez zadania. Peerami tabeli
po~stronie producenta są~inne tabele wpadające do~tych~samych zadań
downstream, w~grafie dwudzielnym nie~istnieje bowiem bezpośrednia krawędź
tabela–tabela. Widoczne w~listingu bramkowanie wskaźnikami odpowiada wprost
drugiemu ze~wzorów z~sekcji~\ref{sec:lineage-detector}. Bez~niego score
mierzyłby jedynie własności otoczenia, a~nie~niespójność tabeli
z~jej~otoczeniem.

\begin{lstlisting}[caption={Rdzeń wyznaczania score'u peer-consistency
(moduł \texttt{completeness.py}).},label={lst:peer},language=Python,float=htbp]
# Strona producenta: peerzy = inne wejscia do downstream-jobow tej tabeli.
downstream_jobs = [v for _, v, d in G.out_edges(t, data=True)
                   if d.get("relation_type") == "DATA_FLOW"
                   and G.nodes[v].get("asset_type") == "Data Job"]
prod_peers = set()
for j in downstream_jobs:
    prod_peers.update(job_inputs.get(j, ()))
prod_peers.discard(t)
prod_score = (np.mean([1.0 if in_df[p] > 0 else 0.0 for p in prod_peers])
              if prod_peers else 0.0)

# Bramkowana anomalia: liczy sie tylko na brzegu (korzen/lisc).
anomaly = 0.0
if in_df[t] == 0:
    anomaly = max(anomaly, prod_score)
if out_df[t] == 0:
    anomaly = max(anomaly, cons_score)
\end{lstlisting}

\FloatBarrier
\subsection{Protokół indukcyjny w~kodzie}

Pętla ewaluacji realizuje schemat leave-one-graph-out, w~którym dla~każdego grafu
testowego model uczy się~na~wszystkich pozostałych. Listing~\ref{lst:logo}
przedstawia jej~szkielet. Kluczowa dla~poprawności protokołu jest~kolejność
operacji. Instancje wszystkich grafów budowane są~przed pętlą,
przy~wspólnym ziarnie, a~pula treningowa powstaje przez~wykluczenie grafu
testowego z~gotowego słownika. Model nie~ma~więc~dostępu do~żadnej
informacji z~grafu testowego, także pośrednio przez~dobór usuwanych zadań.

\begin{lstlisting}[caption={Szkielet pętli leave-one-graph-out
(moduł \texttt{run\_detection\_experiments.py}).},label={lst:logo},
language=Python,float=htbp]
datasets = build_datasets(graphs, removal_ratio, seed, side,
                          exclude_isolated=exclude_isolated)
for test_g in datasets:
    ds_test = datasets[test_g]
    train_pool = [datasets[g] for g in datasets if g != test_g]

    for name in FIT_MODELS:
        clf = _make_model(name, seed).fit(train_pool)
        m = node_detection_metrics(ds_test["labels"], clf.score(ds_test))
        rows.append({"graph": test_g, "algorithm": name, **m})
\end{lstlisting}

Wyniki z~kolejnych ziaren agregowane są~średnią z~trymowaniem wartości
skrajnych (odrzucenie najmniejszej i~największej), co~ogranicza wpływ
pojedynczego nietypowego losowania na~oszacowanie.

\subsection{Złożoność obliczeniowa}

Koszt przetwarzania jednego grafu jest~liniowy względem liczby krawędzi.
Wyznaczenie stopni \emph{DATA\_FLOW}, ekstrakcja cech oraz~budowa zbiorów
peerów wymagają stałej liczby przejść po~krawędziach, a~metadane przepływu
(głębokość i~zasięg w~DAG) wyznaczane są~przeszukiwaniem grafu, również
w~czasie $O(|V| + |E|)$. Dominującym składnikiem kosztu nie~jest~zatem
przetwarzanie grafu, lecz~trenowanie klasyfikatorów. W~protokole
leave-one-graph-out model uczony jest~od~nowa dla~każdego grafu testowego,
co~przy~18~grafach, trzech ziarnach i~czterech wariantach detektora daje
216~pełnych treningów. Bieg ablacyjny z~sekcji~\ref{sec:ablacja} zajął
z~tego~powodu 412~sekund, podczas gdy sama ekstrakcja cech dla~wszystkich
18~grafów trwa poniżej sekundy. Skala zbioru DLG-DG-23 (do~17~085~węzłów
na~graf) nie~stanowi więc~ograniczenia wydajnościowego.

% ============================================================
\section{Testy jednostkowe i~odtwarzalność}
\label{sec:testy}
% ============================================================

\subsection{Zakres testów}

Poprawność implementacji zweryfikowano zestawem 83~testów jednostkowych
uruchamianych narzędziem \texttt{pytest}. W~wersji implementacji
opisanej w~pracy wszystkie testy kończą się~powodzeniem.
Tabela~\ref{tab:testy} przedstawia ich~rozkład na~moduły.

Poza poprawnością techniczną testy sprawdzają dwie własności istotne
metodycznie. Pierwszą jest~brak przecieku między zbiorem treningowym
a~testowym, weryfikowany przez sprawdzenie, czy~identyfikatory grafów w~obu
zbiorach są~rozłączne. Drugą jest~determinizm symulacji, który oznacza,
że~przy~ustalonym ziarnie losowości powtórne wykonanie daje identyczny zbiór
usuniętych zadań i~identyczne etykiety. Bez~tej własności liczby raportowane
w~rozdziale~4 nie~byłyby~odtwarzalne.

\begin{table}[ht]
\centering
\caption{Rozkład testów jednostkowych na~moduły.}
\label{tab:testy}
\small
\begin{tabular}{lr>{\raggedright\arraybackslash}p{7.2cm}}
\toprule
\textbf{Plik testowy} & \textbf{Liczba} & \textbf{Weryfikowane własności} \\
\midrule
\texttt{test\_loader.py} & 29 &
  Poprawność typów węzłów i~krawędzi, zgodność liczby obiektów z~plikami
  źródłowymi, wczytanie wszystkich 18~grafów \\
\texttt{test\_splitter.py} & 17 &
  Rozłączność podziałów, zgodność liczności podzbiorów \\
\texttt{test\_job\_removal.py} & 12 &
  Identyfikacja zadań łączących, poprawność etykiet, rozłączność zbiorów
  treningowego i~testowego, działanie flagi \texttt{exclude\_isolated} \\
\texttt{test\_scenario\_splitter.py} & 9 &
  Scenariusze usuwania krawędzi (badanie wstępne) \\
\texttt{test\_dutkiewicz\_adapter.py} & 6 &
  Mapowanie benchmarku Dutkiewicza et~al.\ na~model
  tabela–zadanie–\emph{DATA\_FLOW} (graf DUT) \\
\texttt{test\_node\_features.py} & 4 &
  Wymiar wektora cech, zgodność z~\texttt{feature\_names()},
  wartości brzegowe dla~węzłów izolowanych \\
\texttt{test\_completeness.py} & 3 &
  Poprawność score'u peer-consistency, bramkowanie sygnału anomalii \\
\texttt{test\_lineage\_detector.py} & 3 &
  Interfejs \texttt{fit}/\texttt{score}, zakres zwracanych
  prawdopodobieństw, warianty ablacyjne \\
\midrule
\textbf{Razem} & \textbf{83} & \\
\bottomrule
\end{tabular}
\end{table}

\subsection{Odtwarzalność wyników}

Wszystkie liczby przedstawione w~rozdziale~4 pochodzą z~jednego
\emph{biegu kanonicznego} o~ustalonej konfiguracji: 18~grafów,
\texttt{removal\_ratio}~$=0{,}10$, \texttt{side}~$=$~\texttt{both},
trzy powtórzenia z~ziarnami 42--44, protokół indukcyjny leave-one-graph-out,
próg wiarygodności $\mathrm{MIN\_POS} = 5$, wariant uczciwy
(\texttt{exclude\_isolated}). Wyniki pośrednie zapisywane są~w~postaci plików
CSV w~katalogu \texttt{results/}, osobno dla~wariantu uczciwego
i~zawyżonego, wraz z~opisem konfiguracji. Eksperymenty poboczne
(ablacja, analiza wrażliwości, test cross-dataset) mają~własne skrypty
i~własne pliki wynikowe, wymienione w~listingu~\ref{lst:repo}, co~pozwala
odtworzyć każdą tabelę rozdziału~4 niezależnie od~pozostałych.

% ============================================================
\section{Dostępność implementacji i~pochodzenie algorytmów}
\label{sec:repozytorium}
% ============================================================

Całość implementacji opisanej w~niniejszej pracy, obejmującą moduły
wczytywania danych, symulacji broken lineage, ekstrakcji cech, metod
detekcji oraz~ewaluacji, wraz z~testami jednostkowymi i~skryptami
odtwarzającymi eksperymenty z~rozdziału~4, udostępniono w~publicznym
repozytorium Git
pod~adresem \texttt{https://github.com/Catcuss/magisterka-data-lineage}
(dostęp zweryfikowano 19~sierpnia 2026~r.). Repozytorium zawiera instrukcję
uruchomienia (dodatek~\ref{dod:uruchomienie}) oraz~pliki konfiguracyjne biegu
kanonicznego, co~zapewnia odtwarzalność raportowanych wyników.

Metody z~grupy heurystyk oraz~klasyczne klasyfikatory nie~są~oryginalnym
wkładem pracy, lecz~stanowią punkty odniesienia (\emph{baseline}) zaczerpnięte
z~literatury i~z~publicznych bibliotek. Wkładem własnym jest~ich~adaptacja
do~zadania detekcji zainfekowanych węzłów w~grafie lineage, autorski detektor
LineageDetector oraz~adapter grafu DUT. Pochodzenie i~zakres adaptacji
poszczególnych składników zestawiono poniżej:

\begin{itemize}
  \item \textbf{Heurystyki strukturalne}: inspirowane rodziną miar
        predykcji krawędzi w~grafach (Common Neighbors, Jaccard, Adamic–Adar,
        Preferential Attachment)~\cite{LibenNowell2007,Martinez2016}. Ponieważ graf
        przepływu danych jest~dwudzielny, miary oparte na~wspólnych sąsiadach
        są~w~nim nieskuteczne. Zaadaptowano je~więc~do~postaci scoringu
        pojedynczego węzła opartego na~stopniu i~roli w~DAG.
  \item \textbf{Random Forest}: implementacja z~biblioteki
        \texttt{scikit-learn}~\cite{scikit-learn}. Adaptacja: wejściem jest~
        autorski wektor cech topologicznych węzła (tab.~\ref{tab:cechy})
        zamiast cech pary obiektów, a~niezbalansowanie klas obsłużono przez
        \texttt{class\_weight}.
  \item \textbf{RUSBoost}: implementacja z~biblioteki
        \texttt{imbalanced-learn}~\cite{imblearn}, wybrana dla~bezpośredniego
        nawiązania do~pracy Boińskiego et~al.~\cite{Boinski2025}. Adaptacja:
        klasyfikacja węzła zamiast pary, na~cechach topologicznych.
  \item \textbf{LightGBM}: biblioteka gradientowego boostingu
        firmy Microsoft~\cite{Ke2017}. Adaptacja: klasyfikacja
        węzła na~cechach topologicznych, z~kompensacją rozkładu klas przez
        \texttt{class\_weight}. Hiperparametrów nie~strojono, a~wszystkie
        klasyfikatory porównywane są~w~konfiguracjach domyślnych
        (por.~sekcja~\ref{sec:trafnosc}).
  \item \textbf{LineageDetector}, wkład własny: łączy powyższe cechy
        z~autorskim relacyjnym score'em kompletności peer-consistency
        (sekcja~\ref{sec:lineage-detector}) i~kalibracją izotoniczną.
        Nie~ma~bezpośredniego odpowiednika w~literaturze.
  \item \textbf{Adapter grafu DUT}, wkład własny: przekształca benchmark
        Dutkiewicza et~al.~\cite{Dutkiewicz2026} z~wierszowej reprezentacji
        lineage do~modelu tabela–zadanie–\emph{DATA\_FLOW}
        (sekcja~\ref{sec:dane-dut}). Sam benchmark wraz ze~scenariuszami
        pochodzi z~pracy źródłowej.
\end{itemize}

Rozgraniczenie komponentów zapożyczonych, z~jawnie wskazanym źródłem,
od~elementów autorskich pozwala jednoznacznie ocenić rzeczywisty wkład
pracy i~odróżnić go~od~narzędzi ogólnego przeznaczenia. Ma~ono również
znaczenie dla~interpretacji wyników rozdziału~4. Klasyfikatory bazowe
i~LineageDetector korzystają z~tego samego wektora cech topologicznych,
tych samych implementacji bibliotecznych i~tego samego potoku ewaluacji.
Detektor różni się~od~klasyfikatora Random Forest, na~którym jest~oparty,
przede wszystkim elementami autorskimi, czyli score'em kompletności
i~kalibracją. Wpływ każdego z~tych elementów na~wynik zmierzono bezpośrednio
w~badaniu ablacyjnym (sekcja~\ref{sec:ablacja}), co~pozwala oddzielić
wkład własny od~skuteczności samych metod bibliotecznych.

% ============================================================

\chapter{Problem i~sposób jego rozwiązania}

Rozdział precyzuje formalną definicję problemu detekcji węzłów
o~zerwanym lineage (sekcja~\ref{sec:sformulowanie}). Sekcja~\ref{sec:dane}
charakteryzuje wykorzystane dane, czyli zbiór DLG-DG-23 oraz~zewnętrzny
graf DUT wraz z~adapterem, sekcja~\ref{sec:symulacja} opisuje
metodykę symulacji broken lineage przez~usuwanie zadań,
a~sekcja~\ref{sec:architektura} architekturę potoku przetwarzania.
Zaimplementowane metody detekcji omówiono w~sekcji~\ref{sec:algorytmy},
realizację techniczną w~sekcji~\ref{sec:implementacja}, testy jednostkowe
i~odtwarzalność wyników w~sekcji~\ref{sec:testy}, a~dostępność implementacji
i~pochodzenie algorytmów w~sekcji~\ref{sec:repozytorium}.

% ============================================================
\section{Sformułowanie problemu}
\label{sec:sformulowanie}
% ============================================================

Graf lineage modeluje się jako skierowany graf DAG przepływu danych
$G = (V, E)$, gdzie $V = V_{J} \cup V_{T} \cup V_{F}$ obejmuje węzły
typu zadanie (\emph{Data Job}), tabela (\emph{Data Table}) i~pole,
a~$E$ zawiera krawędzie przepływu danych (\emph{DATA\_FLOW}) między zadaniami
a~tabelami. Zjawisko broken lineage odpowiada usunięciu z~grafu pewnego
zadania $J$. Tabele z~nim incydentne, czyli jego poprzedniki (tabele
karmiące $J$) oraz~następniki (tabele karmione przez~$J$), tracą część
informacji o~swoim pochodzeniu i~stają się~\emph{zainfekowane}
(ang.~\emph{infected}).

Obserwowany graf ma~postać $G_{\mathrm{obs}} = G \setminus \{J_1,\dots,J_m\}$,
powstałą przez~usunięcie podzbioru zadań. Zadaniem detekcji jest wyznaczenie
funkcji, która każdej tabeli $t \in V_{T}$ przypisuje wynik
$s(t) \in [0, 1]$ interpretowany jako~oszacowane prawdopodobieństwo
$p(\text{tabela } t \text{ jest zainfekowana})$. Wynikiem nie~jest~decyzja
binarna przy~ustalonym progu, lecz~ranking tabel. Tabele o~najwyższym
wyniku są~pierwszymi kandydatami do~ręcznej inspekcji.

Etykiety referencyjne (\emph{ground truth}) generowane są~automatycznie.
Tabela jest zainfekowana wtedy i~tylko wtedy, gdy była incydentna
do~co~najmniej jednego usuniętego zadania. Zbiór kandydatów ograniczono
do~węzłów typu \emph{Data Table}. Węzły typu pole (\emph{field}) łączy
z~tabelami wyłącznie relacja \emph{PARENT\_CHILD}, dlatego nie~podlegają one detekcji.

Zasadnicza trudność zadania polega na~tym, że~osierocona tabela-następnik
po~usunięciu zadania wygląda jak~prawdziwa tabela źródłowa
($\mathrm{in\_df} = 0$). Rozróżnienie tych~dwóch przypadków wymaga cech
opisujących pozycję węzła w~grafie DAG oraz~uczenia indukcyjnego
(cross-graph), a~nie~samego stopnia węzła.

% ============================================================
\section{Zbiory danych}
\label{sec:dane}
% ============================================================

W~pracy wykorzystano dwa źródła danych o~odmiennym pochodzeniu i~przeznaczeniu.
Wszystkie eksperymenty główne, obejmujące trening i~ewaluację metod
w~protokole indukcyjnym, przeprowadzono na~rzeczywistych grafach zbioru
DLG-DG-23. Uzupełnieniem jest~zewnętrzny graf DUT, zbudowany na~podstawie
benchmarku Dutkiewicza et~al.~\cite{Dutkiewicz2026}, który służy wyłącznie
do~sprawdzenia przenośności modeli wytrenowanych na~DLG-DG-23. Oba źródła
opisano w~kolejnych podsekcjach.

\subsection{Zbiór DLG-DG-23}
\label{sec:dane-dlg}

Zbiór DLG-DG-23 pochodzi z~pracy Chen et~al.~\cite{Chen2024}
i~zawiera 18~grafów lineage pobranych ze~środowiska produkcyjnego
Huawei Cloud. Dane zostały zanonimizowane, a~nazwy tabel i~atrybutów
zastąpiono identyfikatorami w~formacie UUID. Schematów bazy danych
nie~udostępniono. Zbiór opublikowano w~2024~roku jako pierwszy
otwarty zestaw rzeczywistych grafów lineage.

Każdy graf zawiera węzły trzech typów: \emph{Data Job} (zapytanie SQL
lub~zadanie ETL), \emph{Data Table}
(tabela lub~widok przechowujący dane) oraz~\emph{field} (atrybut lub~kolumna
należąca do~tabeli). Łączą je~krawędzie dwóch typów: \emph{DATA\_FLOW}
(skierowana krawędź przepływu danych między tabelą a~zadaniem)
oraz~\emph{PARENT\_CHILD} (relacja przynależności pola do~tabeli). Pełny
graf jest~zatem \emph{heterogeniczny}. Przedmiotem detekcji są~wyłącznie
węzły typu \emph{Data Table}, a~sygnał czerpany jest z~podgrafu przepływu
danych (krawędzie \emph{DATA\_FLOW} między węzłami \emph{Data Table}
i~\emph{Data Job}). Podgraf ten~ma~strukturę dwudzielną i~skierowaną,
interpretowaną jako graf DAG. Grafy różnią się skalą, a~liczba węzłów waha się
od~278 do~17~085.

Zbiór DLG-DG-23 wybrano, ponieważ jest pierwszym publicznie udostępnionym
zestawem rzeczywistych grafów lineage~\cite{Chen2024}, a~w~trakcie prowadzenia
badań nie~zidentyfikowano innego zbioru tego rodzaju. Anonimizacja nazw
węzłów ma~bezpośredni wpływ na~dobór metod, ponieważ metryki podobieństwa
oparte na~nazwach obiektów (np.~miara Jaro–Winklera stosowana przez Boińskiego
et~al.~\cite{Boinski2025}) są~nieaplikowalne. W~niniejszej pracy sygnał
ograniczono zatem do~atrybutów topologicznych grafu. Ograniczenie
to~jest zarazem częścią pytania badawczego o~to, czy~sama struktura grafu
lineage wystarcza do~wskazania tabel o~zerwanym pochodzeniu danych.
Metoda skuteczna przy~tak ograniczonym wglądzie w~graf ma~przy~tym większą
szansę na~zastosowanie praktyczne niż~podejścia wymagające pełnych
metadanych.

Tabela~\ref{tab:dlg-stats} przedstawia statystyki poszczególnych
grafów zbioru DLG-DG-23.

\begin{table}[ht]
\centering
\caption{Charakterystyka grafów zbioru DLG-DG-23.}
\label{tab:dlg-stats}
\begin{threeparttable}
\begin{tabular}{llrl}
\toprule
\textbf{ID} & \textbf{Węzły} & \textbf{Krawędzie DATA\_FLOW} & \textbf{Skala} \\
\midrule
DLG1  &    298 &    24 & mały   \\
DLG2  &    464 &    37 & mały   \\
DLG3  &    603 &    57 & mały   \\
DLG4  &    415 &    32 & mały   \\
DLG5  & 1\,299 &   175 & średni \\
DLG6  & 13\,840 & 1\,443 & duży  \\
DLG7  &    574 &    43 & mały   \\
DLG8  &    546 &    34 & mały   \\
DLG9  &    756 &    62 & mały   \\
DLG10 & 5\,378 &   497 & duży   \\
DLG11 & 2\,024 &   121 & średni \\
DLG12 & 5\,645 &   469 & duży   \\
DLG13 & 2\,157 &   229 & średni \\
DLG14 &    453 &    24 & mały   \\
DLG15 &    558 &    41 & mały   \\
DLG16 & 5\,574 &   745 & duży   \\
DLG17 &    651 &    38 & mały   \\
DLG18 & 17\,085 & 1\,694 & duży \\
\bottomrule
\end{tabular}
\begin{tablenotes}
\small
\item Skala: mały $<$ 1\,000 węzłów, średni 1\,000–5\,000, duży $>$ 5\,000.
\end{tablenotes}
\end{threeparttable}
\end{table}

\subsection{Zewnętrzny graf DUT}
\label{sec:dane-dut}

Drugim źródłem danych jest~benchmark opracowany przez Dutkiewicza
et~al.~\cite{Dutkiewicz2026} w~zespole badawczym promotora
(sekcja~\ref{sec:literatura}). Obejmuje on~syntetyczne scenariusze broken
lineage, wygenerowane dla~przykładowej bazy \emph{Northwind}~\cite{Northwind}.
Lineage zapisany jest~w~nim na~poziomie pojedynczych wierszy, w~postaci grafu
wiedzy RDF oraz~plików CSV, w~których każdy rekord wiąże wartość źródłową
z~wartością docelową, która z~niej powstała. Jest to~zatem reprezentacja
drobnoziarnista (sekcja~\ref{sec:lineage_def}), niezgodna z~gruboziarnistym
modelem tabela–zadanie, na~którym operują metody opisane w~niniejszej pracy.

Aby~ocenić na~tych danych modele wytrenowane na~DLG-DG-23 bez~zmian w~ich
implementacji, opracowano adapter (moduł \texttt{dutkiewicz\_adapter.py}),
który agreguje zależności wierszowe do~poziomu tabel. Dla~każdej tabeli
docelowej tworzony jest~węzeł typu \emph{Data Job}, do~którego prowadzą
krawędzie \emph{DATA\_FLOW} od~wszystkich tabel źródłowych, z~których
pochodzą jej~wiersze. Z~węzła tego prowadzi następnie krawędź \emph{DATA\_FLOW}
do~samej tabeli docelowej. Kolumny kluczy występujące w~zapisie lineage
dodawane są~jako~węzły typu \emph{field}, połączone z~tabelą relacją
\emph{PARENT\_CHILD}. Wynikowy graf ma~te~same typy węzłów i~krawędzi
co~grafy DLG-DG-23, dzięki czemu przechodzi przez potok przetwarzania
identycznie jak~one.

Graf DUT liczy 21~tabel, 14~zadań i~33~krawędzie przepływu danych. Jest więc
o~rzędy wielkości mniejszy od~grafów DLG-DG-23 i~ma~odmienną, łańcuchową
topologię, w~której tabele rzadko współdzielą zadania. Adapter nie~odtwarza
przy~tym pełnych schematów tabel, ponieważ benchmark ich~nie~zawiera,
dlatego liczba pól jest~mniejsza niż~w~DLG-DG-23. Poprawność odwzorowania
weryfikuje sześć testów jednostkowych (tabela~\ref{tab:testy}). Graf DUT
nie~jest~wykorzystywany w~treningu żadnej z~metod, a~wyniki uzyskane na~nim
przedstawiono w~sekcji~\ref{sec:cross_dataset}.

% ============================================================
\section{Metodyka symulacji broken lineage przez usuwanie zadań}
\label{sec:symulacja}
% ============================================================

Ponieważ zbiór DLG-DG-23 zawiera w~pełni zaobserwowane grafy lineage,
zjawisko broken lineage symulowane jest przez~kontrolowane usuwanie
zadań (nie~pojedynczych krawędzi). Podejście to~odwzorowuje
przyczynę broken lineage, ponieważ zniknięcie tabeli tymczasowej lub~ukrycie
logiki przez UDF odpowiada utracie całego węzła-zadania wraz z~jego krawędziami.

\subsection{Procedura symulacji i~konstrukcja etykiet}
\label{sec:procedura-symulacji}

Kandydatami do~usunięcia są~\emph{zadania łączące}, czyli zadania mające
co~najmniej jedną wejściową i~jedną wyjściową krawędź \emph{DATA\_FLOW}
($\mathrm{in\_df} > 0 \land \mathrm{out\_df} > 0$). Usunięcie takiego zadania
zrywa lineage po~obu stronach, ponieważ tabele karmiące je~tracą konsumenta,
a~tabele przez~nie~karmione tracą producenta.

Instancję klasyfikacji węzłów buduje się~z~grafu $G$ w~kilku krokach.
Najpierw losowany jest~podzbiór zadań łączących przeznaczonych do~usunięcia,
którego udział reguluje parametr \texttt{removal\_ratio}. W~eksperymentach
przyjęto wartość 10\%, a~wpływ tego parametru zbadano w~analizie wrażliwości
(sekcja~\ref{sec:przeciek}). Zbiór tabel zainfekowanych wyznacza się~jako~sumę
tabel incydentnych do~usuniętych zadań. Graf obserwowany $G_{\mathrm{obs}}$
powstaje następnie przez~usunięcie wybranych zadań wraz z~ich krawędziami.
Kandydatami do~detekcji są~wszystkie tabele pozostałe w~$G_{\mathrm{obs}}$,
przy~czym etykieta~1 oznacza tabelę zainfekowaną, a~etykieta~0 tabelę czystą.

Dodatkowo zaimplementowano wariant kontrolny, w~którym każda instancja
usuwa dokładnie jedno zadanie, co~daje czystą interpretację
przyczynową (jedno zdarzenie broken lineage $=$ jedno brakujące zadanie).

\subsection{Kontrola przecieku – wykluczenie węzłów izolowanych}
\label{sec:exclude-isolated}

Usunięcie zadania może~pozbawić tabelę wszystkich krawędzi
przepływu danych ($\mathrm{in\_df} = 0 \land \mathrm{out\_df} = 0$).
Ponieważ w~oryginalnych grafach tabele izolowane nie~występują, trywialna
reguła „tabela izolowana $\Rightarrow$ zainfekowana'' pozwala rozpoznać
znaczną część pozytywów bez~żadnego uczenia i~sztucznie zawyża wynik.
Aby~zapewnić uczciwą ewaluację, wprowadzono flagę \texttt{exclude\_isolated},
która usuwa z~puli kandydatów tabele w~pełni izolowane. Wariant ten~stawia
trudniejsze i~bardziej realistyczne zadanie, polegające na~wykryciu tabel
o~częściowo zerwanym lineage, które nadal mają część połączeń. W~pracy raportowane
są~oba warianty, a~wariant uczciwy (z~wykluczeniem izolowanych) traktowany
jest jako~główny.

Tabela~\ref{tab:symulacja} przedstawia skalę symulacji dla~czterech wartości
parametru \texttt{removal\_ratio}. Łącznie w~18~grafach zbioru DLG-DG-23
znajduje się 1453~zadania łączące (od~3 do~404 na~graf), spośród których usuwana
jest~losowa część. Liczby podano dla~wariantu przed wykluczeniem tabel
izolowanych, ponieważ tylko wtedy można zmierzyć ich~udział wśród pozytywów.

\begin{table}[ht]
\centering
\caption[Charakterystyka symulacji broken lineage]{Charakterystyka symulacji broken lineage w~zależności
od~parametru \texttt{removal\_ratio}. Wartości uśrednione po~18~grafach
i~trzech ziarnach losowości (42--44). W~ostatniej kolumnie podano odsetek
tabel zainfekowanych, które utraciły wszystkie krawędzie przepływu
danych i~są~rozpoznawalne trywialną regułą izolacji.}
\label{tab:symulacja}
\begin{tabular}{ccccc}
\toprule
\texttt{removal\_ratio} & \textbf{Usuniętych} & \textbf{Tabel} & \textbf{Udział} & \textbf{Pozytywów} \\
 & \textbf{zadań} & \textbf{zainfekowanych} & \textbf{pozytywów} & \textbf{izolowanych} \\
\midrule
0{,}05 & 4{,}2 & 15{,}1 & 11{,}4\% & 19{,}7\% \\
0{,}10 & 8{,}0 & 25{,}1 & 19{,}0\% & 22{,}9\% \\
0{,}20 & 16{,}2 & 45{,}2 & 34{,}1\% & 29{,}8\% \\
0{,}30 & 24{,}3 & 62{,}6 & 47{,}3\% & 34{,}5\% \\
\bottomrule
\end{tabular}
\end{table}

Dwie obserwacje z~tabeli mają~znaczenie dla~konstrukcji eksperymentu.
Przyjęta w~biegu kanonicznym wartość $0{,}10$ daje średnio 25~tabel
zainfekowanych na~graf przy~19\% pozytywów. Zadanie pozostaje więc~wyraźnie
niezrównoważone, ale~klasa pozytywna jest~dostatecznie liczna, by~metryki
rankingowe były wiarygodne. Udział pozytywów trywialnych (tabel w~pełni
izolowanych) rośnie z~kolei z~19{,}7\% do~34{,}5\% wraz ze~wzrostem
odsetka usuwanych zadań. Im~silniejsze zaburzenie, tym~większa część zadania
daje się~rozwiązać bez~żadnego uczenia. Liczby te~określają skalę przecieku,
któremu zapobiega opisana wyżej flaga \texttt{exclude\_isolated}. Przy~biegu
kanonicznym usuwa ona~z~puli kandydatów blisko jedną czwartą pozytywów,
i~to~tych najłatwiejszych.

% ============================================================
\section{Architektura rozwiązania}
\label{sec:architektura}
% ============================================================

Potok przetwarzania składa się z~czterech komponentów: modułu wczytywania
grafów (wraz z~adapterem grafu DUT), modułu symulacji broken lineage (budowa
instancji klasyfikacji węzłów), metod detekcji oraz~modułu ewaluacji indukcyjnej.

Rysunek~\ref{rys:potok} przedstawia schemat potoku przetwarzania, od~wczytania
grafów, przez symulację broken lineage i~ekstrakcję cech, po~scoring detektora
i~wyznaczenie metryk rankingowych. Oba źródła danych trafiają do~wspólnego
potoku. Grafy DLG-DG-23 wczytywane są~bezpośrednio z~plików JSON, natomiast
graf DUT przechodzi najpierw przez adapter, który sprowadza go~do~tego samego
modelu danych. Od~etapu symulacji oba rodzaje grafów przetwarzane są~identycznie,
różni je~jedynie rola w~ewaluacji.

\begin{figure}[ht]
  \centering
  \includegraphics[width=\linewidth]{rys_potok.pdf}
  \caption[Schemat potoku przetwarzania]{Schemat potoku przetwarzania. Grafy
           DLG-DG-23 oceniane są~w~protokole indukcyjnym leave-one-graph-out
           (trening na~17~grafach, test na~grafie niewidzianym). Graf DUT,
           sprowadzony adapterem do~modelu tabela–zadanie, służy wyłącznie
           do~testu modeli wytrenowanych na~DLG-DG-23.}
  \label{rys:potok}
\end{figure}

\subsection{Moduł wczytywania danych}

Dane wejściowe każdego grafu przechowywane są~w~dwóch plikach JSON:
\texttt{DLGi-node.json} (opisy węzłów z~atrybutem \texttt{asset\_type})
oraz~\texttt{DLGi-edge.json} (krawędzie z~atrybutami \texttt{relation\_type}
i~\texttt{relation\_id}). Moduł \texttt{loader.py} buduje z~nich obiekt
\texttt{nx.DiGraph} (NetworkX) i~zwraca słownik wszystkich 18~grafów.
Graf DUT wczytywany jest~przez moduł \texttt{dutkiewicz\_adapter.py}, który
zwraca obiekt tego samego typu, z~atrybutami węzłów i~krawędzi zgodnymi
z~konwencją DLG-DG-23 (sekcja~\ref{sec:dane-dut}). Dalsze moduły potoku
nie~rozróżniają więc~źródła grafu.

\subsection{Moduł symulacji usuwania zadań}

Moduł \texttt{job\_removal.py} realizuje metodykę z~sekcji~\ref{sec:symulacja}:
identyfikuje zadania łączące, usuwa ich podzbiór, wyznacza zbiór zainfekowanych
tabel i~zwraca instancję (graf $G_{\mathrm{obs}}$, listę kandydatów,
etykiety, listę usuniętych zadań). Obsługuje warianty \texttt{side},
określający stronę incydencji, po~której tabele uznaje się~za~zainfekowane
(strona producenta, konsumenta lub~obie), oraz~\texttt{exclude\_isolated}. Funkcja \texttt{inductive\_split} dzieli
identyfikatory grafów na~rozłączne zbiory treningowy i~testowy (cross-graph).

\subsection{Moduł ewaluacji i~metryki rankingowe}

Moduł \texttt{node\_metrics.py} wyznacza metryki rankingowe: AUC-ROC,
AUC-PR, Precision@k, Recall@k, Hits@k oraz~Brier score. Wyniki agregowane
są~po~grafach testowych w~protokole indukcyjnym. Dla~stabilności oszacowań
stosuje się powtórzenia z~różnymi ziarnami i~trymowanie wartości skrajnych
(min/max).

Ewaluację prowadzi się~w~protokole indukcyjnym cross-graph.
Model trenuje się~na~podzbiorze grafów, a~testuje na~grafach niewidzianych
(leave-one-graph-out lub~podział zbioru 18~grafów). Zbiory treningowy
i~testowy są~rozłączne, co~sprawdza zdolność metody do~generalizacji
na~nowe grafy, a~nie~jedynie do~zapamiętania jednego. Konsekwencją tego
protokołu jest~wykluczenie metody Node2Vec~\cite{Grover2016}, która uczy zanurzeń
węzłów transduktywnie, dla~konkretnego grafu, przez~co~embeddingi
są~geometrycznie nieporównywalne między grafami~\cite{Hamilton2017}.
Obserwacja ta~motywuje zastosowanie indukcyjnych
grafowych sieci neuronowych (GraphSAGE~\cite{Hamilton2017},
R-GCN~\cite{Schlichtkrull2018}, HGT~\cite{Hu2020}), wskazanych w~rozdziale~5 jako~kierunek dalszych badań.

% ============================================================
\section{Zaimplementowane metody detekcji}
\label{sec:algorytmy}
% ============================================================

Zaimplementowano trzy grupy metod detekcji zainfekowanych tabel, różniące się
ilością wykorzystywanej informacji i~stopniem złożoności. Heurystyki strukturalne
(sekcja~\ref{sec:heurystyki}) przypisują tabeli wynik na~podstawie prostych
reguł i~nie~wymagają uczenia, dlatego wyznaczają dolny punkt odniesienia.
Klasyczne klasyfikatory uczenia maszynowego (sekcja~\ref{sec:ml}) uczą się
na~grafach treningowych zależności między wektorem cech topologicznych
a~etykietą, co~pozwala ocenić, ile~wnosi samo uczenie indukcyjne. Autorski
detektor LineageDetector (sekcja~\ref{sec:lineage-detector}) rozszerza ten~wektor
o~relacyjny score kompletności i~uzupełnia model o~kalibrację prawdopodobieństw.

Wszystkie metody mają wspólny interfejs. Dla~każdej tabeli-kandydata zwracają
wynik z~przedziału $[0, 1]$, na~podstawie którego budowany jest~ranking. Dzięki
temu każda z~nich jest~oceniana w~tym samym potoku, na~tych samych instancjach
i~tymi samymi metrykami, a~różnice wyników można przypisać samym metodom.

% ============================================================
\subsection{Heurystyki strukturalne}
\label{sec:heurystyki}
% ============================================================

Heurystyki strukturalne przypisują tabeli wynik podejrzliwości wyłącznie
na~podstawie jej pozycji w~grafie, bez~etapu uczenia. Są~naturalnym
punktem odniesienia (\emph{baseline}) dla~metod uczonych. Miary oparte
na~wspólnych sąsiadach (Common Neighbors, Jaccard,
Adamic-Adar)~\cite{LibenNowell2007} są~w~grafie
dwudzielnym nieskuteczne, ponieważ węzły tego~samego typu nigdy nie~mają
wspólnych sąsiadów bezpośrednich. Z~tego powodu wykluczono je~z~implementacji.

W~pracy zaimplementowano cztery heurystyki bazowe. Piątym punktem
odniesienia w~tej~grupie jest~sam \emph{score kompletności}
(sekcja~\ref{sec:lineage-detector}), oceniany dodatkowo w~oderwaniu
od~detektora, co~pozwala rozdzielić wkład sygnału relacyjnego od~wkładu
uczenia. W~zestawieniach rozdziału~4 występuje on~zatem jako~osobna metoda
obok czterech poniższych:

\begin{itemize}
  \item \textbf{Anomalia stopnia} (\texttt{degree\_anomaly}): wynik
        $1/(1+\deg(t))$, gdzie $\deg(t)$ jest stopniem całkowitym tabeli.
        Heurystyka opiera się~na~założeniu, że~tabela, która utraciła
        zadanie, ma~mniej połączeń niż~przed zerwaniem lineage.
  \item \textbf{Brzeg grafu} (\texttt{boundary}): suma wskaźników
        $\mathbf{1}[\mathrm{in\_df}=0] + \mathbf{1}[\mathrm{out\_df}=0]$.
        Tabele na~brzegu przepływu (korzenie i~liście) są~podejrzane,
        ponieważ zerwanie lineage często przesuwa tabelę na~brzeg grafu.
  \item \textbf{Niska łączność z~zadaniami}
        (\texttt{low\_job\_connectivity}): wynik $1/(1+n_{J}(t))$,
        gdzie $n_{J}(t)$ to~liczba unikatowych zadań sąsiadujących
        z~tabelą przez krawędzie \emph{DATA\_FLOW}.
  \item \textbf{Reguła korzenia} (\texttt{rule\_root}): binarny wskaźnik
        $\mathbf{1}[\mathrm{in\_df}=0]$, zgodnie z~którym każda tabela
        bez~producenta jest uznawana za~podejrzaną. To~celowo trywialny punkt odniesienia,
        obrazujący siłę najprostszego możliwego symptomu.
\end{itemize}

Heurystyki te~mierzą wyłącznie lokalny symptom (stopień, rola).
Ich~ograniczenie, przede wszystkim nieodróżnianie osieroconej tabeli
od~prawdziwego źródła, motywuje metody uczone oraz~autorski detektor
LineageDetector (sekcja~\ref{sec:lineage-detector}), który wzbogaca opis
węzła o~relacyjny score kompletności (peer-consistency).

% ============================================================
\subsection{Klasyczne metody uczenia maszynowego}
\label{sec:ml}
% ============================================================

Klasyczne metody ML traktują detekcję jako~binarną klasyfikację węzła, w~której
wejściem jest wektor cech opisujących pozycję tabeli w~grafie DAG,
a~wyjściem prawdopodobieństwo zainfekowania. Ze~względu na~anonimizację
(UUID) jako~cechy przyjęto wyłącznie atrybuty topologiczne, ponieważ metryki
oparte na~nazwach są~nieaplikowalne.

\subsubsection{Wektor cech węzła}

Tabela~\ref{tab:cechy} przedstawia pełny wektor 13~cech topologicznych
wyznaczanych dla~każdej tabeli-kandydata. Cechy opisują pozycję węzła
w~grafie przepływu danych, bo~to~ona, a~nie~sam stopień, pozwala odróżnić
prawdziwą tabelę źródłową od~tabeli osieroconej po~usunięciu zadania.

\begin{table}[ht]
\centering
\caption[Wektor cech topologicznych węzła]{Wektor cech topologicznych węzła (moduł \texttt{node\_features}).}
\label{tab:cechy}
\small
\begin{tabular}{cl>{\raggedright\arraybackslash}p{5.2cm}>{\raggedright\arraybackslash}p{4.6cm}}
\toprule
\textbf{Nr} & \textbf{Cecha} & \textbf{Opis} & \textbf{Rola w~detekcji} \\
\midrule
1 & \texttt{in\_df} & Liczba wejściowych krawędzi \emph{DATA\_FLOW}
    (od~zadań) & Brak producenta ($=0$) to~główny symptom zerwania \\
2 & \texttt{out\_df} & Liczba wyjściowych krawędzi \emph{DATA\_FLOW}
    (do~zadań) & Symetryczny symptom, czyli utrata konsumenta \\
3 & \texttt{degree} & Stopień całkowity węzła (wszystkie typy krawędzi)
    & Ogólna łączność, spadająca po~usunięciu zadania \\
4 & \texttt{role} & Rola w~przepływie (źródło, pośrednia, ujście,
    poza przepływem) & Pozycja funkcjonalna tabeli w~DAG \\
5 & \texttt{flow\_depth} & Głębokość w~DAG (odległość od~źródeł)
    & Osierocona tabela ma~sztucznie małą głębokość \\
6 & \texttt{flow\_reach} & Zasięg w~dół DAG (odległość do~ujść)
    & Symetrycznie sztucznie mały zasięg \\
7 & \texttt{depth\_plus\_reach} & Suma głębokości i~zasięgu
    & Długość ścieżki przez węzeł, wykrywa „skrócone'' łańcuchy \\
8 & \texttt{is\_root} & Wskaźnik $\mathbf{1}[\mathrm{in\_df}=0]$
    & Binarna wersja symptomu utraty producenta \\
9 & \texttt{is\_leaf} & Wskaźnik $\mathbf{1}[\mathrm{out\_df}=0]$
    & Binarna wersja symptomu utraty konsumenta \\
10 & \texttt{n\_field\_children} & Liczba pól tabeli (krawędzie
    \emph{PARENT\_CHILD}) & Rozmiar tabeli, sygnał niezależny
    od~przepływu \\
11 & \texttt{n\_job\_neighbors} & Liczba unikatowych zadań sąsiadujących
    & Łączność z~warstwą zadań \\
12 & \texttt{avg\_job\_in\_df} & Średni stopień wejściowy sąsiednich zadań
    & Charakterystyka otoczenia węzła \\
13 & \texttt{avg\_job\_out\_df} & Średni stopień wyjściowy sąsiednich
    zadań & Jak~wyżej, od~strony wyjść \\
\bottomrule
\end{tabular}
\end{table}

Detektor LineageDetector (sekcja~\ref{sec:lineage-detector})
rozszerza ten wektor o~trzy cechy kompletności
(\texttt{peer\_producer\_score}, \texttt{peer\_consumer\_score},
\texttt{completeness\_anomaly}), opisane w~sekcji~\ref{sec:lineage-detector}.

\subsubsection{Random Forest}

Random Forest (RF)~\cite{Breiman2001} to~metoda uczenia zespołowego oparta na~uśrednianiu
predykcji drzew decyzyjnych trenowanych na~losowych podpróbkach danych
i~cech (\emph{bagging}). Dla~obsługi niezrównoważonych klas zastosowano
\texttt{class\_weight=\textquotedbl{}balanced\textquotedbl{}}.

\subsubsection{RUSBoost}

RUSBoost (\emph{Random Under-Sampling Boosting})~\cite{Seiffert2010} łączy boosting AdaBoost
z~losowym podpróbkowaniem klasy większościowej przed~każdą iteracją.
Zaimplementowano go~dla~bezpośredniego nawiązania do~wyników Boińskiego
et~al.~\cite{Boinski2025}, przy~czym w~niniejszej pracy klasyfikuje on~węzły,
a~nie~pary obiektów.

\subsubsection{LightGBM}

LightGBM~\cite{Ke2017} to~implementacja gradientowego boostingu drzew decyzyjnych,
budująca kolejne drzewa sekwencyjnie, każde minimalizując błąd resztkowy
poprzednika. Niezbalansowanie klas kompensuje parametr
\texttt{class\_weight}~$=$~\texttt{balanced}, przypisujący obserwacjom wagi
odwrotnie proporcjonalne do~liczności klasy.

% ============================================================
\subsection{Autorski detektor LineageDetector}
\label{sec:lineage-detector}
% ============================================================

Autorski detektor \texttt{LineageDetector} (klasa
\texttt{LineageHealthDetector}) łączy w~jednym skalibrowanym modelu trzy
elementy. Podstawę stanowią cechy pozycji w~DAG, wyznaczane przez moduł
\texttt{node\_features} i~opisujące węzeł lokalnie. Uzupełnia je~score
kompletności \emph{peer-consistency}, opisany w~dalszej części sekcji,
który wnosi sygnał relacyjny i~stanowi główny element wkładu własnego pracy.
Całość zamyka kalibracja prawdopodobieństw metodą izotoniczną
(\texttt{CalibratedClassifierCV}), dzięki której wynik detektora można
interpretować jako~oszacowanie $p(\text{zainfekowana})$, a~nie~jedynie
jako~wartość porządkującą ranking.

\subsubsection{Score kompletności peer-consistency}

Zaproponowany w~pracy score kompletności opiera się na~założeniu,
że~tabela o~zerwanym lineage jest~\emph{strukturalnie niespójna}
ze~swoimi sąsiadami (peerami) w~grafie DAG. Dla~tabeli-korzenia
($\mathrm{in\_df} = 0$) rozważa się~zbiór peerów-wejść, czyli inne tabele
wpadające do~tych~samych zadań downstream. Jeśli~peerzy mają
producenta, a~rozpatrywana tabela nie, to~jej brak producenta jest~podejrzany,
ponieważ prawdziwe źródło danych sąsiaduje zwykle z~innymi źródłami. Symetrycznie
dla~tabeli-liścia ($\mathrm{out\_df} = 0$) i~jej peerów po~stronie konsumpcji.
Formalnie dla~tabeli $t$ definiuje się:
%
\[
  \mathrm{prod}(t) = \frac{1}{|P(t)|}\sum_{p \in P(t)} \mathbf{1}[\mathrm{in\_df}(p) > 0],
  \qquad
  \mathrm{cons}(t) = \frac{1}{|C(t)|}\sum_{p \in C(t)} \mathbf{1}[\mathrm{out\_df}(p) > 0],
\]
%
gdzie $P(t)$ i~$C(t)$ oznaczają zbiory peerów odpowiednio po~stronie
producenta i~konsumenta, a~$\mathbf{1}[\cdot]$ funkcję wskaźnikową.
Bramkowany sygnał anomalii aktywuje się~wyłącznie na~brzegu grafu:
%
\[
  \mathrm{anomaly}(t) =
  \max\bigl(\mathbf{1}[\mathrm{in\_df}(t)=0]\cdot\mathrm{prod}(t),\;
            \mathbf{1}[\mathrm{out\_df}(t)=0]\cdot\mathrm{cons}(t)\bigr).
\]
%
Sygnał ten~ma~charakter relacyjny. Różni się~jakościowo od~lokalnego
stopnia węzła i~dostarcza informacji niedostępnej klasyfikatorowi
korzystającemu wyłącznie z~cech lokalnych. Definicja jest~symetryczna
względem obu~stron przepływu, jednak na~zbiorze DLG-DG-23 aktywna
okazuje się~wyłącznie strona producenta, ponieważ każde zadanie w~tym~zbiorze
zapisuje dokładnie jedną tabelę, przez~co~zbiór $C(t)$ jest~zawsze pusty.
Konsekwencje tej~własności omówiono w~sekcji~\ref{sec:waznosc}. Trzy wynikające z~niego cechy
(\texttt{peer\_producer\_score}, \texttt{peer\_consumer\_score},
\texttt{completeness\_anomaly}) rozszerzają wektor cech z~tabeli~\ref{tab:cechy}
i~są~dostępne wyłącznie detektorowi LineageDetector.

Bazowym estymatorem jest~domyślnie Random Forest, a~wariantem
alternatywnym LightGBM. Interfejs \texttt{fit}/\texttt{score} jest~zgodny z~pozostałymi
klasyfikatorami, co~umożliwia bezpośrednie porównanie w~tym~samym potoku.
Konstrukcja detektora umożliwia badania ablacyjne. Wyłączenie kalibracji
(\texttt{calibrate=False}) izoluje jej~wpływ na~Brier score, a~wyłączenie
cech kompletności (\texttt{use\_completeness=False}) wpływ sygnału
relacyjnego na~AUC.

% ============================================================
\section{Implementacja}
\label{sec:implementacja}
% ============================================================

Sekcja opisuje realizację techniczną potoku z~sekcji~\ref{sec:architektura}:
zastosowane narzędzia, organizację kodu oraz~te~fragmenty implementacji,
które są~istotne dla~zrozumienia metody, a~nie~wynikają wprost z~jej~opisu
formalnego.

\subsection{Stos technologiczny}

Całość zaimplementowano w~języku Python. Tabela~\ref{tab:stos} zestawia
wykorzystane biblioteki wraz z~wersjami, w~których wykonano bieg kanoniczny.
Podanie wersji jest~istotne dla~odtwarzalności, ponieważ implementacje algorytmów
uczenia maszynowego bywają modyfikowane między wydaniami, co~może~wpływać
na~wyniki liczbowe.

\begin{table}[ht]
\centering
\caption{Stos technologiczny i~wersje bibliotek użyte w~biegu kanonicznym.}
\label{tab:stos}
\begin{tabular}{lll}
\toprule
\textbf{Składnik} & \textbf{Wersja} & \textbf{Rola w~pracy} \\
\midrule
Python~\cite{Python}                 & 3.10.2  & Środowisko uruchomieniowe \\
NetworkX~\cite{NetworkX}             & 3.4.2   & Reprezentacja i~przetwarzanie grafów \\
NumPy~\cite{NumPy}                   & 1.26.4  & Operacje macierzowe, wektory cech \\
pandas~\cite{pandas}                 & 1.5.1   & Agregacja wyników, eksport CSV \\
scikit-learn~\cite{scikit-learn}     & 1.6.1   & Random Forest, kalibracja izotoniczna, metryki \\
LightGBM~\cite{Ke2017}               & 4.6.0   & Gradientowy boosting drzew \\
imbalanced-learn~\cite{imblearn}     & 0.14.1  & RUSBoost \\
SciPy~\cite{SciPy}                   & 1.15.3  & Test kolejności par Wilcoxona \\
Matplotlib~\cite{Matplotlib}         & 3.8.0   & Rysunki \\
pytest~\cite{pytest}                 & 9.0.2   & Testy jednostkowe \\
\bottomrule
\end{tabular}
\end{table}

\subsection{Organizacja repozytorium}

Kod podzielono na~pakiety odpowiadające etapom potoku, a~skrypty
uruchomieniowe oddzielono od~modułów bibliotecznych, aby~każdy eksperyment
z~rozdziału~4 dawał się~odtworzyć jednym poleceniem. Listing~\ref{lst:repo}
przedstawia strukturę katalogów.

\begin{lstlisting}[caption={Struktura repozytorium (fragmenty istotne
dla~detekcji węzłów).},label={lst:repo},numbers=none,language={},float=htbp]
src/
  data/
    loader.py               # wczytanie 18 grafow DLG-DG-23 z JSON
    dutkiewicz_adapter.py   # adapter benchmarku zewnetrznego (graf DUT)
  detection/
    job_removal.py          # symulacja broken lineage, etykiety, podzial
    node_features.py        # 13 cech topologicznych wezla
    completeness.py         # score peer-consistency (wklad wlasny)
    lineage_detector.py     # LineageHealthDetector
    node_classifier.py      # heurystyki + klasyczne klasyfikatory
    node_metrics.py         # metryki rankingowe
    run_detection_experiments.py   # bieg kanoniczny (rozdz. 4.2)
    run_ablation.py                # badanie ablacyjne (rozdz. 4.5)
    run_cross_dataset.py           # test cross-dataset (rozdz. 4.9)
    run_removal_sensitivity.py     # analiza wrazliwosci (rozdz. 4.6)
scripts/
  test_istotnosci.py        # testy Wilcoxona (rozdz. 4.4)
  tabela_symulacji.py       # tabela charakterystyki symulacji
tests/                      # 83 testy jednostkowe (pytest)
results/                    # wyniki biegow w formacie CSV
\end{lstlisting}

\subsection{Format danych i~wczytywanie}

Każdy graf zbioru DLG-DG-23 opisany jest~parą plików JSON: lista węzłów
z~atrybutem \texttt{asset\_type} oraz~lista krawędzi z~atrybutami
\texttt{relation\_id} i~\texttt{relation\_type} (listing~\ref{lst:json}).
Moduł \texttt{loader.py} przekształca je~w~obiekt \texttt{nx.DiGraph}
i~przenosi typ węzła oraz~typ relacji do~atrybutów grafu. Dzięki temu
pozostałe moduły filtrują podgraf przepływu danych bez~ponownego sięgania
do~plików.

\begin{lstlisting}[caption={Fragment plików wejściowych
\texttt{DLGi-node.json} i~\texttt{DLGi-edge.json} (identyfikatory skrócone).},
label={lst:json},numbers=none,language={},float=htbp]
// DLG1-node.json
{"nodes": [
  {"asset_id": "3f2a...c81", "asset_type": "Data Table"},
  {"asset_id": "9b07...4de", "asset_type": "Data Job"},
  {"asset_id": "a154...77f", "asset_type": "Data Field"}
]}

// DLG1-edge.json
{"edges": [
  {"relation_id": "r-0012", "relation_type": "DATA_FLOW",
   "source": "3f2a...c81", "target": "9b07...4de"},
  {"relation_id": "r-0013", "relation_type": "PARENT_CHILD",
   "source": "3f2a...c81", "target": "a154...77f"}
]}
\end{lstlisting}

% Listingi 3.1 i 3.2 maja zostac zlozone przed ta podsekcja, aby nie
% rozcinaly jej pierwszego akapitu.
\FloatBarrier
\subsection{Wyznaczanie score'u peer-consistency}

Najistotniejszym fragmentem implementacji jest~wyznaczanie relacyjnego
score'u kompletności, zdefiniowanego wzorami w~sekcji~\ref{sec:lineage-detector}.
Listing~\ref{lst:peer} pokazuje jego~rdzeń. Sam wzór nie~precyzuje, że~zbiory
peerów budowane są~przez zadania. Peerami tabeli
po~stronie producenta są~inne tabele wpadające do~tych~samych zadań
downstream, w~grafie dwudzielnym nie~istnieje bowiem bezpośrednia krawędź
tabela–tabela. Widoczne w~listingu bramkowanie wskaźnikami odpowiada wprost
drugiemu ze~wzorów z~sekcji~\ref{sec:lineage-detector}. Bez~niego score
mierzyłby jedynie własności otoczenia, a~nie~niespójność tabeli
z~jej~otoczeniem.

\begin{lstlisting}[caption={Rdzeń wyznaczania score'u peer-consistency
(moduł \texttt{completeness.py}).},label={lst:peer},language=Python,float=htbp]
# Strona producenta: peerzy = inne wejscia do downstream-jobow tej tabeli.
downstream_jobs = [v for _, v, d in G.out_edges(t, data=True)
                   if d.get("relation_type") == "DATA_FLOW"
                   and G.nodes[v].get("asset_type") == "Data Job"]
prod_peers = set()
for j in downstream_jobs:
    prod_peers.update(job_inputs.get(j, ()))
prod_peers.discard(t)
prod_score = (np.mean([1.0 if in_df[p] > 0 else 0.0 for p in prod_peers])
              if prod_peers else 0.0)

# Bramkowana anomalia: liczy sie tylko na brzegu (korzen/lisc).
anomaly = 0.0
if in_df[t] == 0:
    anomaly = max(anomaly, prod_score)
if out_df[t] == 0:
    anomaly = max(anomaly, cons_score)
\end{lstlisting}

\FloatBarrier
\subsection{Protokół indukcyjny w~kodzie}

Pętla ewaluacji realizuje schemat leave-one-graph-out, w~którym dla~każdego grafu
testowego model uczy się~na~wszystkich pozostałych. Listing~\ref{lst:logo}
przedstawia jej~szkielet. Kluczowa dla~poprawności protokołu jest~kolejność
operacji. Instancje wszystkich grafów budowane są~przed pętlą,
przy~wspólnym ziarnie, a~pula treningowa powstaje przez~wykluczenie grafu
testowego z~gotowego słownika. Model nie~ma~więc~dostępu do~żadnej
informacji z~grafu testowego, także pośrednio przez~dobór usuwanych zadań.

\begin{lstlisting}[caption={Szkielet pętli leave-one-graph-out
(moduł \texttt{run\_detection\_experiments.py}).},label={lst:logo},
language=Python,float=htbp]
datasets = build_datasets(graphs, removal_ratio, seed, side,
                          exclude_isolated=exclude_isolated)
for test_g in datasets:
    ds_test = datasets[test_g]
    train_pool = [datasets[g] for g in datasets if g != test_g]

    for name in FIT_MODELS:
        clf = _make_model(name, seed).fit(train_pool)
        m = node_detection_metrics(ds_test["labels"], clf.score(ds_test))
        rows.append({"graph": test_g, "algorithm": name, **m})
\end{lstlisting}

Wyniki z~kolejnych ziaren agregowane są~średnią z~trymowaniem wartości
skrajnych (odrzucenie najmniejszej i~największej), co~ogranicza wpływ
pojedynczego nietypowego losowania na~oszacowanie.

\subsection{Złożoność obliczeniowa}

Koszt przetwarzania jednego grafu jest~liniowy względem liczby krawędzi.
Wyznaczenie stopni \emph{DATA\_FLOW}, ekstrakcja cech oraz~budowa zbiorów
peerów wymagają stałej liczby przejść po~krawędziach, a~metadane przepływu
(głębokość i~zasięg w~DAG) wyznaczane są~przeszukiwaniem grafu, również
w~czasie $O(|V| + |E|)$. Dominującym składnikiem kosztu nie~jest~zatem
przetwarzanie grafu, lecz~trenowanie klasyfikatorów. W~protokole
leave-one-graph-out model uczony jest~od~nowa dla~każdego grafu testowego,
co~przy~18~grafach, trzech ziarnach i~czterech wariantach detektora daje
216~pełnych treningów. Bieg ablacyjny z~sekcji~\ref{sec:ablacja} zajął
z~tego~powodu 412~sekund, podczas gdy sama ekstrakcja cech dla~wszystkich
18~grafów trwa poniżej sekundy. Skala zbioru DLG-DG-23 (do~17~085~węzłów
na~graf) nie~stanowi więc~ograniczenia wydajnościowego.

% ============================================================
\section{Testy jednostkowe i~odtwarzalność}
\label{sec:testy}
% ============================================================

\subsection{Zakres testów}

Poprawność implementacji zweryfikowano zestawem 83~testów jednostkowych
uruchamianych narzędziem \texttt{pytest}. W~wersji implementacji
opisanej w~pracy wszystkie testy kończą się~powodzeniem.
Tabela~\ref{tab:testy} przedstawia ich~rozkład na~moduły.

Poza poprawnością techniczną testy sprawdzają dwie własności istotne
metodycznie. Pierwszą jest~brak przecieku między zbiorem treningowym
a~testowym, weryfikowany przez sprawdzenie, czy~identyfikatory grafów w~obu
zbiorach są~rozłączne. Drugą jest~determinizm symulacji, który oznacza,
że~przy~ustalonym ziarnie losowości powtórne wykonanie daje identyczny zbiór
usuniętych zadań i~identyczne etykiety. Bez~tej własności liczby raportowane
w~rozdziale~4 nie~byłyby~odtwarzalne.

\begin{table}[ht]
\centering
\caption{Rozkład testów jednostkowych na~moduły.}
\label{tab:testy}
\small
\begin{tabular}{lr>{\raggedright\arraybackslash}p{7.2cm}}
\toprule
\textbf{Plik testowy} & \textbf{Liczba} & \textbf{Weryfikowane własności} \\
\midrule
\texttt{test\_loader.py} & 29 &
  Poprawność typów węzłów i~krawędzi, zgodność liczby obiektów z~plikami
  źródłowymi, wczytanie wszystkich 18~grafów \\
\texttt{test\_splitter.py} & 17 &
  Rozłączność podziałów, zgodność liczności podzbiorów \\
\texttt{test\_job\_removal.py} & 12 &
  Identyfikacja zadań łączących, poprawność etykiet, rozłączność zbiorów
  treningowego i~testowego, działanie flagi \texttt{exclude\_isolated} \\
\texttt{test\_scenario\_splitter.py} & 9 &
  Scenariusze usuwania krawędzi (badanie wstępne) \\
\texttt{test\_dutkiewicz\_adapter.py} & 6 &
  Mapowanie benchmarku Dutkiewicza et~al.\ na~model
  tabela–zadanie–\emph{DATA\_FLOW} (graf DUT) \\
\texttt{test\_node\_features.py} & 4 &
  Wymiar wektora cech, zgodność z~\texttt{feature\_names()},
  wartości brzegowe dla~węzłów izolowanych \\
\texttt{test\_completeness.py} & 3 &
  Poprawność score'u peer-consistency, bramkowanie sygnału anomalii \\
\texttt{test\_lineage\_detector.py} & 3 &
  Interfejs \texttt{fit}/\texttt{score}, zakres zwracanych
  prawdopodobieństw, warianty ablacyjne \\
\midrule
\textbf{Razem} & \textbf{83} & \\
\bottomrule
\end{tabular}
\end{table}

\subsection{Odtwarzalność wyników}

Wszystkie liczby przedstawione w~rozdziale~4 pochodzą z~jednego
\emph{biegu kanonicznego} o~ustalonej konfiguracji: 18~grafów,
\texttt{removal\_ratio}~$=0{,}10$, \texttt{side}~$=$~\texttt{both},
trzy powtórzenia z~ziarnami 42--44, protokół indukcyjny leave-one-graph-out,
próg wiarygodności $\mathrm{MIN\_POS} = 5$, wariant uczciwy
(\texttt{exclude\_isolated}). Wyniki pośrednie zapisywane są~w~postaci plików
CSV w~katalogu \texttt{results/}, osobno dla~wariantu uczciwego
i~zawyżonego, wraz z~opisem konfiguracji. Eksperymenty poboczne
(ablacja, analiza wrażliwości, test cross-dataset) mają~własne skrypty
i~własne pliki wynikowe, wymienione w~listingu~\ref{lst:repo}, co~pozwala
odtworzyć każdą tabelę rozdziału~4 niezależnie od~pozostałych.

% ============================================================
\section{Dostępność implementacji i~pochodzenie algorytmów}
\label{sec:repozytorium}
% ============================================================

Całość implementacji opisanej w~niniejszej pracy, obejmującą moduły
wczytywania danych, symulacji broken lineage, ekstrakcji cech, metod
detekcji oraz~ewaluacji, wraz z~testami jednostkowymi i~skryptami
odtwarzającymi eksperymenty z~rozdziału~4, udostępniono w~publicznym
repozytorium Git
pod~adresem \texttt{https://github.com/Catcuss/magisterka-data-lineage}
(dostęp zweryfikowano 19~sierpnia 2026~r.). Repozytorium zawiera instrukcję
uruchomienia (dodatek~\ref{dod:uruchomienie}) oraz~pliki konfiguracyjne biegu
kanonicznego, co~zapewnia odtwarzalność raportowanych wyników.

Metody z~grupy heurystyk oraz~klasyczne klasyfikatory nie~są~oryginalnym
wkładem pracy, lecz~stanowią punkty odniesienia (\emph{baseline}) zaczerpnięte
z~literatury i~z~publicznych bibliotek. Wkładem własnym jest~ich~adaptacja
do~zadania detekcji zainfekowanych węzłów w~grafie lineage, autorski detektor
LineageDetector oraz~adapter grafu DUT. Pochodzenie i~zakres adaptacji
poszczególnych składników zestawiono poniżej:

\begin{itemize}
  \item \textbf{Heurystyki strukturalne}: inspirowane rodziną miar
        predykcji krawędzi w~grafach (Common Neighbors, Jaccard, Adamic–Adar,
        Preferential Attachment)~\cite{LibenNowell2007,Martinez2016}. Ponieważ graf
        przepływu danych jest~dwudzielny, miary oparte na~wspólnych sąsiadach
        są~w~nim nieskuteczne. Zaadaptowano je~więc~do~postaci scoringu
        pojedynczego węzła opartego na~stopniu i~roli w~DAG.
  \item \textbf{Random Forest}: implementacja z~biblioteki
        \texttt{scikit-learn}~\cite{scikit-learn}. Adaptacja: wejściem jest~
        autorski wektor cech topologicznych węzła (tab.~\ref{tab:cechy})
        zamiast cech pary obiektów, a~niezbalansowanie klas obsłużono przez
        \texttt{class\_weight}.
  \item \textbf{RUSBoost}: implementacja z~biblioteki
        \texttt{imbalanced-learn}~\cite{imblearn}, wybrana dla~bezpośredniego
        nawiązania do~pracy Boińskiego et~al.~\cite{Boinski2025}. Adaptacja:
        klasyfikacja węzła zamiast pary, na~cechach topologicznych.
  \item \textbf{LightGBM}: biblioteka gradientowego boostingu
        firmy Microsoft~\cite{Ke2017}. Adaptacja: klasyfikacja
        węzła na~cechach topologicznych, z~kompensacją rozkładu klas przez
        \texttt{class\_weight}. Hiperparametrów nie~strojono, a~wszystkie
        klasyfikatory porównywane są~w~konfiguracjach domyślnych
        (por.~sekcja~\ref{sec:trafnosc}).
  \item \textbf{LineageDetector}, wkład własny: łączy powyższe cechy
        z~autorskim relacyjnym score'em kompletności peer-consistency
        (sekcja~\ref{sec:lineage-detector}) i~kalibracją izotoniczną.
        Nie~ma~bezpośredniego odpowiednika w~literaturze.
  \item \textbf{Adapter grafu DUT}, wkład własny: przekształca benchmark
        Dutkiewicza et~al.~\cite{Dutkiewicz2026} z~wierszowej reprezentacji
        lineage do~modelu tabela–zadanie–\emph{DATA\_FLOW}
        (sekcja~\ref{sec:dane-dut}). Sam benchmark wraz ze~scenariuszami
        pochodzi z~pracy źródłowej.
\end{itemize}

Rozgraniczenie komponentów zapożyczonych, z~jawnie wskazanym źródłem,
od~elementów autorskich pozwala jednoznacznie ocenić rzeczywisty wkład
pracy i~odróżnić go~od~narzędzi ogólnego przeznaczenia. Ma~ono również
znaczenie dla~interpretacji wyników rozdziału~4. Klasyfikatory bazowe
i~LineageDetector korzystają z~tego samego wektora cech topologicznych,
tych samych implementacji bibliotecznych i~tego samego potoku ewaluacji.
Detektor różni się~od~klasyfikatora Random Forest, na~którym jest~oparty,
przede wszystkim elementami autorskimi, czyli score'em kompletności
i~kalibracją. Wpływ każdego z~tych elementów na~wynik zmierzono bezpośrednio
w~badaniu ablacyjnym (sekcja~\ref{sec:ablacja}), co~pozwala oddzielić
wkład własny od~skuteczności samych metod bibliotecznych.
